% !TEX root = ../../tree_notes.tex
% =====================================================================
%  Chapter 27 -- Offline Tree Queries: Mo's on Trees and Tarjan LCA  WRITTEN
%  (folder ch21_mos_virtual_trees; label prefix c21; PRINTS as
%  Chapter 27; folder number 21; build with:  python tools/build.py --only 21)
%
%  NUMBERING TRAP (the Ch.22/23/24/25/26 trap, mirrored): \chref{c21} prints
%  "Chapter 27" (THIS chapter).  Do NOT confuse the label with the printed
%  number: \chref{c30} prints "Chapter 26" (virtual trees, the PREVIOUS
%  chapter), \chref{c22} prints "Chapter 28" (isomorphism/hashing, the NEXT
%  chapter), \chref{c23} prints "Chapter 25" (link-cut), \chref{c11} prints
%  "Chapter 16" (LCA -- the ancestor oracle Mo-on-paths and Tarjan both lean
%  on), \chref{c24} prints "Chapter 5" (Euler tour / tin-tout -- the flatten
%  Mo rides on), \chref{c16} prints "Chapter 22" (HLD), \chref{c18} prints
%  "Chapter 24" (centroid), \chref{c10} prints "Chapter 13" (tree DP),
%  \chref{c09} prints "Chapter 12" (tree metrics), \chref{c28} prints
%  "Chapter 17" (DSU on tree / small-to-large), \chref{c20} prints "Chapter 15"
%  (rerooting), \chref{c14} prints "Chapter 18" (Fenwick trees -- owns the
%  solved CF 383C this chapter cites in prose only), \chref{c04} prints
%  "Chapter 4" (representation).
%  Every cite of THIS chapter is \chref{c21}.
%
%  Owns OFFLINE TREE QUERIES answered by reordering, not by a fold: Mo's
%  algorithm on a subtree (flatten with \tin/\tout so a subtree is a contiguous
%  array range, then sweep ranges with a sqrt-decomposition order that adds and
%  removes ONE node per step), Mo's algorithm on a PATH (a path is NOT
%  contiguous, so keep the doubled Euler tour and TOGGLE each node on/off; split
%  at \lca and handle the LCA node by an occurrence-parity special case), and
%  Tarjan's OFFLINE LCA (one DFS + DSU + "answer the queries whose second node
%  just turned black" -- O((n+q) alpha) with no lifting table).  It STARTS where
%  Ch.26 (virtual trees, \chref{c30}) ENDS: the virtual-tree boundary hands the
%  "many overlapping offline node-set queries with a non-invertible statistic"
%  case to Mo (Ch.27).  LCA, \tin/\tout and the distance formula are USED, never
%  re-derived (Ch.16/Ch.5); Tarjan is presented as an ALTERNATIVE offline LCA,
%  back-citing \chref{c11}, not a re-derivation of binary lifting.
%  Sections:
%    * 27.1 the Mo engine -- offline range queries on an ARRAY: block size
%      B ~ sqrt(n), sort queries by (left-block, right), move two pointers with
%      O(1) add/remove, total O((n+q) sqrt n).  Solved CF 617E "XOR and Favorite
%      Number" (2200, L4) -- count subarrays with (xor == k), the canonical Mo
%      warm-up; the freq-of-prefix-xor add/remove is the whole game.  Live
%      samples: [1 2 1 1 0 3], q (1,6),(3,5) -> 7,0 ; [1 1 1 1 1] -> 9,4,4.
%    * 27.2 Mo on SUBTREES -- flatten: the preorder array P[tin[v]] makes
%      subtree(v) the contiguous range [tin[v], tout[v]], so "aggregate over a
%      subtree" is a range query answered by the 27.1 engine.  Solved CF 463E
%      "Caisa and Tree" (2100, L4) -- deepest strict ancestor sharing a prime
%      with v, plus <=50 point updates: the DFS stack IS the root-path window
%      (per-prime last-seen table, O(1) push/undo, answer before push) and one
%      rebuild per static block; the sample is replayed to -1,1,2,-1,1.  The
%      invertible signed-update case CF 383C is cited in PROSE only -- its URL
%      and full solved treatment live in Ch.18 (\chref{c14}).  The
%      section then states the Mo-on-subtrees recipe (block-order the ranges when
%      the statistic has NO inverse, e.g. distinct-colour count -- CF 375D is
%      cited in PROSE only, its URL lives in Ch.17).
%    * 27.3 Mo on PATHS -- a root-to-node path is contiguous in the doubled
%      Euler tour but an arbitrary path u-v is not; keep E[1..2n] with each node
%      twice, TOGGLE membership as pointers cross an occurrence, split at
%      \lca(u,v): if the LCA is an endpoint use [st,st]; else use
%      [en[u],st[v]] and ADD the LCA by hand (its occurrences never both land in
%      the range).  WORKED example on the running tree: distinct colours on a
%      path, fully traced, including the parity trap.  NO solvedproblem (count
%      stays 2).
%    * 27.4 Tarjan offline LCA -- one DFS; a node is WHITE/GREY/BLACK, union a
%      child into its parent on return, and when v turns BLACK every query
%      (v,u) whose other endpoint u is already BLACK is answered by
%      ancestor[find(u)].  O((n+q) alpha(n)) with NO lifting table, but OFFLINE
%      only.  WORKED example: run it on the running tree to a batch of LCAs,
%      then contrast with Ch.16 binary lifting (online, O(log n)).  NO
%      solvedproblem.
%    * 27.5 recognition and the method boundary -- the Mo tell (offline batch of
%      overlapping subtree/path queries, an add/remove that is cheap but the
%      statistic has NO inverse so a BIT/segment tree cannot maintain it); and
%      the hand-offs: invertible range stats -> Euler+Fenwick (Ch.5/Ch.10);
%      ONLINE path queries -> HLD (Ch.22) or the LCT (Ch.25); per-set folds ->
%      virtual tree (Ch.26); small marked sets -> DSU-on-tree (Ch.17); whole-tree
%      distance distributions -> centroid (Ch.24).  WORKED examples only.
%  Chapter-end Summit: CF 1585E "Frequency Queries" (2400, L5) -- top of the
%  band.  Rooted tree, integer a_i on each node, q queries (v,l,k): walk the
%  path v -> root, tally each value's frequency, keep values with freq >= l,
%  order by frequency, report the k-th value (or -1).  It is the "non-invertible
%  offline path statistic" this chapter exists for: sweep root-to-node with a
%  freq[frequency] bucket structure (DSU-on-tree style add/remove / persistent
%  map along the root path), answering every query that lives on one root path.
%  Both official samples hand-replayed (1 -1 1 ; 1 1 2 1 -1).
%
%  Diagram-first: this chapter is TREE based.  The running example is the shared
%  rooted tree R used across sections (defined once, authoritatively, in
%  00_opener).  Rooted at 1, edges (parent -> child, 1-based):
%    1->2, 1->3, 2->4, 2->5, 3->6, 3->7, 4->8
%  (8 nodes).  Preorder P = [1,2,4,8,5,3,6,7] fixes tin(1)=1,tin(2)=2,tin(4)=3,
%  tin(8)=4,tin(5)=5,tin(3)=6,tin(6)=7,tin(7)=8; subtree(2)=P[2..5]={2,4,8,5}.
%  The doubled Euler tour E=[1,2,4,8,8,4,5,5,2,3,6,6,7,7,3,1] drives 27.3.
%  Colours val[1..8]=[2,1,3,1,2,3,1,2].  The Summit uses its OWN tiny cases (the
%  CF samples), never R, so no trace collides.
%
%  Anti-repetition contract (teach once, cross-reference by label):
%    * Ch.16 LCA (\chref{c11}) owns \lca, \dist, binary lifting; 27.3 USES \lca
%      to split a path, and 27.4 presents Tarjan as an OFFLINE ALTERNATIVE,
%      back-citing \thmref{c11-distf} but never re-deriving lifting.
%    * Ch.5 Euler tour (\chref{c24}) owns tin/tout and the flatten-to-array
%      trick; 27.2 reuses that flatten for the subtree range, 27.3 extends it to
%      the DOUBLED tour -- the single-visit flatten is cited, not re-taught.
%    * Ch.10 / Ch.5 Fenwick (\chref{c24}) owns invertible range add/point query;
%      27.2 cites Ch.~18's solved CF 383C (\chref{c14}) as that world's template
%      and 27.5 routes "invertible" queries back to it.
%    * Ch.17 DSU on tree (\chref{c28}) owns small-to-large and owns CF 375D;
%      27.2 mentions 375D in PROSE only (its URL lives there) and 27.5 draws the
%      Mo-vs-DSU-on-tree line.
%    * Ch.22 HLD (\chref{c16}) owns ONLINE path ranges; 27.5 hands every online
%      path query to it and says Mo is offline-only.
%    * Ch.26 virtual trees (\chref{c30}) is the PREVIOUS chapter and forward-
%      cited the "many overlapping offline node-set queries" case to THIS one.
%
%  Known URL collisions (confirmed by the urlsweep of chapters/ + the live CF
%  API; prose-cite only, NEVER re-link -- urlsweep must print an empty
%  DUPLICATES section):
%    * CF 375D "Tree and Queries" (1900) -> Ch.17 DSU-on-tree (\chref{c28}).
%      THE canonical subtree-distinct-colour problem; its URL lives there, PROSE
%      only here (named as the motivating Mo-on-subtree fold).
%    * CF 258E "Little Elephant and Tree" (2400) -> already owned elsewhere; the
%      path-colour Mo idea is described in 27.3 WITHOUT its URL.
%    * CF 86D "Powerful array" (2200) -> owned elsewhere; cited as the archetypal
%      sum(freq^2) Mo in 27.5 prose only, never re-linked.
%    * CF 613D "Kingdom and its Cities" (2800) -> Ch.~virtual-tree prose owner;
%      untouched here.
%    * CF 383C "Propagating tree" (2000) -> Ch.18 Fenwick trees (\chref{c14})
%      solves it fully (c14-prop); PROSE-cited here only, never re-linked.
%    * The three roster URLs below (617E, 463E, 1585E) and every mixed-set URL
%      were confirmed ABSENT from every other chapter by a fresh urlsweep and
%      ownership grep over chapters/ (the earlier owned_cf cache was stale -- do
%      NOT trust it; re-derive ownership from the chapter files themselves).
%
%  Roster (every title/rating/tag/sample live-verified against the Codeforces
%  problemset API / problem page BEFORE a trace was written; CF-forward):
%    solved: CF 617E XOR and Favorite Number (2200, L4) (27.1) -- array Mo; live
%            samples replayed -> 7,0 and 9,4,4.
%            CF 463E Caisa and Tree        (2100, L4) (27.2) -- root-path GCD
%            queries on the DFS stack + 51 rebuilds; live sample replayed ->
%            -1,1,2,-1,1.
%    worked: block-order sweep on a toy array (27.1); the preorder flatten of
%            subtree(2) and the parity split (27.2); the doubled-tour toggle and
%            the LCA-parity split for path(8,6) on R (27.3); the WHITE/GREY/BLACK
%            DSU run answering a batch of LCAs on R (27.4); the Mo-vs-five-tools
%            pick (27.5).
%    summit: CF 1585E Frequency Queries    (2400, L5) -- path v->root, keep
%            values with freq >= l, k-th by frequency; offline root-path sweep
%            with a freq-of-frequency bucket.  Both samples hand-replayed
%            (1 -1 1 ; 1 1 2 1 -1).
%    mixed (10, ascending by live-verified rating, every ID confirmed ABSENT from
%            every other chapter by re-deriving ownership from chapters/ --
%            LOCKED at the 99_chapter_end repair, WRITTEN): 739B(1900),
%            1328E(1900), 1702G1(1900), 1454E(2000), 665E(2100), 1790F(2100),
%            1000E(2100), 1806E(2200), 1399E2(2200), 1980G(2300).  A spread of
%            genuine offline-tree-query drills and deliberate REJECTION cases --
%            their hints name the TRUE tool (Fenwick/Euler, HLD, centroid,
%            DSU-on-tree, binary lifting, two pointers, reverse-time DSU) that
%            beats Mo.  (Earlier picks yielded away: 191C/1304E/519E/629E to
%            Ch.16, 1294F and 960E to the advanced-DP folder, 1902F to Ch.22,
%            2063E to the tree-DP-II folder and 1009F to Ch.17 -- one URL per
%            book; every final ID was re-diffed against a fresh scan of
%            chapters/ BEFORE linking.)
% =====================================================================
\bookchapter{Offline Tree Queries: Mo's on Trees and Tarjan LCA}{c21}
\chapterinfo{CF 1800--2400}{Written}

% =====================================================================
%  Chapter 27 opener -- "Chapter at a Glance", three reading paths + the
%  ladder staircase, THE running example (the shared 8-node rooted tree R,
%  its preorder flatten, its doubled Euler tour, and a colouring), the
%  who-owns-what comparison against Ch.16 (LCA), Ch.5 (Euler/tin), Ch.17
%  (DSU on tree), Ch.22 (HLD), Ch.24 (centroid) and Ch.26 (virtual trees),
%  the recognition guide, a constraint->technique table, and notation.
%  Exports fig:c21-tree, fig:c21-taxonomy.
%  No \term{lca} / \term{tin} / \term{Euler tour} here -- they are indexed by
%  Ch.16/Ch.5; \term the NEW objects "Mo's algorithm", "block order", "toggle",
%  "offline LCA" only.
% =====================================================================

\begin{chaptertheorybox}{Chapter at a Glance}
Chapter~5 (\chref{c24}) taught you to \emph{flatten} a tree: one DFS fixes
\(\tin,\tout\), and the preorder array turns every subtree into a contiguous
range.  Chapter~16 (\chref{c11}) taught you to \emph{climb}: binary lifting
answers \(\lca\) and \(\dist\) online in \(O(\log n)\) (\(\thmref{c11-distf}\)).
Both are single-query tools -- order the queries however you like, each is
answered on its own.

This chapter is about the opposite regime: a \emph{huge batch of offline
queries} over the \emph{same} tree, where the statistic you maintain has
\textbf{no inverse}.  "How many distinct colours are on this path?" you cannot
answer by subtracting two prefixes -- there is no \(\mathrm{pref}(r)-
\mathrm{pref}(l-1)\) for "distinct".  When add/remove of one element is cheap
but there is no undo-by-arithmetic, you stop trying to be online and instead
\emph{reorder} the queries so that consecutive ones differ by a single element:
that is Mo's algorithm, a \(\bigO{(n+q)\sqrt n}\) sweep with a
\term{block order}.  And when the queries only ask for \(\lca\), Tarjan's
offline DFS with a disjoint-set union answers all \(q\) of them in
\(\bigO{(n+q)\alpha(n)}\) -- no lifting table at all.

\textbf{You will learn to}
\begin{itemize}[nosep]
  \item run Mo's \term{block order}: sort \(q\) ranges by \((\lfloor L/B\rfloor, R)\),
        move two pointers with \(O(1)\) add/remove, and bound the work by
        \(\bigO{(n+q)\sqrt n}\) with \(B\approx\sqrt n\)
        (\(\secref{c21-moengine}\), solved CF 617E \emph{XOR and Favorite Number}); \hfill\lv{3}--\lv{4}
  \item answer \term{Mo on subtrees}: flatten subtree\((v)\) to \([\tin(v),\tout(v)]\)
        and reuse the engine; carry root-path statistics on the DFS stack
        (\(\secref{c21-subtree}\), solved CF 463E \emph{Caisa and Tree}); \hfill\lv{4}
  \item answer \term{Mo on paths}: keep the doubled Euler tour, \term{toggle}
        membership on each occurrence, and split at \(\lca\) with the LCA-parity
        special case (\(\secref{c21-path}\)); \hfill\lv{5}
  \item compute \term{offline LCA} with Tarjan's DFS+DSU in \(\bigO{(n+q)\alpha}\),
        and know when it beats Ch.~16 lifting (\(\secref{c21-tarjan}\)); \hfill\lv{4}--\lv{5}
  \item read a statement and pick Mo over Fenwick / HLD / DSU-on-tree / virtual
        tree / centroid exactly when the statistic has no inverse and the queries
        are offline (\(\secref{c21-boundary}\)). \hfill\lv{5}
\end{itemize}

\textbf{Prerequisites.} Chapter~5 (\chref{c24}) -- \(\tin/\tout\) and the
flatten-to-array trick: the range a subtree becomes.  Chapter~16 (\chref{c11}) --
\(\lca\), used to split a path (27.3) and the online target Tarjan replaces (27.4).
Chapter~13 (\chref{c10}) -- the intuition that a \(\bigO{n}\) fold per query is too
slow when there are \(q\) of them.  Chapter~17 (\chref{c28}) -- DSU on tree, the
\emph{other} offline answer to subtree statistics (27.5 contrasts them).
Chapter~26 (\chref{c30}) -- the \textbf{previous chapter}; its method boundary hands
the "many overlapping offline node-set queries, non-invertible statistic" case here.
\end{chaptertheorybox}

% ---------------- Reading paths + ladder position ----------------
\begin{chaptertheorybox}{How to Read This Chapter (three paths)}
Level badges \(\lv3\)--\lv{5}: \(\lv3\) is running the block order by hand on a
toy array and counting pointer moves; \(\lv4\) is the two coded solves --
array Mo (617E) and the root-path DFS-stack sweep (463E); \(\lv{5}\) is the doubled-tour
\emph{toggle} for paths and Tarjan's DFS+DSU, capped by the Summit (CF 1585E).

\textbf{Engine path (\(\lv3\)--\(\lv4\), ``make the sweep'').}
\secref{c21-moengine} in full -- block size, the two-pointer move, the
\(\bigO{(n+q)\sqrt n}\) bound, solved CF 617E -- then \secref{c21-subtree}, where
the same engine runs on the flattened subtree ranges (solved CF 463E).

\textbf{Tree-power path (\(\lv4\)--\(\lv{5}\), ``a statistic with no inverse, on
a path'').} \secref{c21-subtree} for the flatten, \secref{c21-path} for the
doubled Euler tour and the toggle + LCA split, then the Summit \(\probref{c21-summit}\),
where a frequency-of-frequency bucket answers offline root-path queries.

\textbf{Boundary path (\(\lv4\)--\(\lv{5}\), ``is this even a Mo job?'').}
\secref{c21-tarjan} (offline LCA vs Ch.~16 lifting) and \secref{c21-boundary}:
the tell (offline, overlapping, no inverse) and the look-alikes that are not --
invertible ranges for Fenwick, online paths for HLD, per-set folds for a virtual
tree, small marked sets for DSU-on-tree, all-pairs for centroid.
\end{chaptertheorybox}

\begin{figure}[H]
\centering
\begin{tikzpicture}[diagram, x=8.5mm, y=6.5mm]
  \foreach \st/\lab in {1/lv1,2/lv2,3/lv3,4/lv4,5/lv5,6/lv6} {
    \fill[\lab!22, draw=\lab, line width=0.8pt] (\st-1,0) rectangle (\st,\st*0.52);
    \node[font=\scriptsize\bfseries, text=\lab] at (\st-0.5,\st*0.52-0.42) {L\st};}
  \node[annot, anchor=south west, font=\scriptsize]
    at (0.05,0.62) {prereq (Ch.5/16): tin/tout flatten, lca -- the range and the split};
  \node[annot, anchor=south west, font=\scriptsize]
    at (2.05,1.66) {27.1 block order, two pointers, the (n+q)sqrt n bound};
  \node[annot, anchor=south west, font=\scriptsize]
    at (3.05,2.18) {27.1 solved CF 617E (array Mo) | 27.2 solved CF 463E (root path)};
  \node[annot, anchor=south west, font=\scriptsize, text=lv5]
    at (4.05,2.70) {27.3 doubled tour + toggle | 27.4 Tarjan offline LCA};
  \node[annot, anchor=south west, font=\scriptsize, text=lv6]
    at (5.05,3.22) {27.5 the pick; Summit CF 1585E (offline root-path frequency)};
\end{tikzpicture}
\caption{Chapter~27 on the six-step ladder: hand-running the block order and
counting moves is \(\lv3\); the coded array-Mo and DFS-stack-root-path solves (617E,
463E) are \(\lv4\); the doubled-tour toggle for paths and Tarjan's offline LCA
are \(\lv{5}\); the Summit (an offline root-path frequency query) is step~6.}
\end{figure}

% ---------------- The running example ----------------
\subsection*{The Running Example: One Tree, Three Flattenings}

Every section in this chapter runs on the same rooted tree \(R\) on \(n=8\)
nodes, fixed once here and authoritative (if a later trace disagrees, this
figure wins).  Rooted at \(1\); parent-to-child edges, 1-based:
\[
  1\!\to\!2,\quad 1\!\to\!3,\quad 2\!\to\!4,\quad 2\!\to\!5,\quad
  3\!\to\!6,\quad 3\!\to\!7,\quad 4\!\to\!8 .
\]
A left-to-right DFS gives the \term{preorder flatten} \(P\) (entry times
\(\tin\)): \(P=[\,1,2,4,8,5,3,6,7\,]\), so \(\tin(1){=}1\), \(\tin(2){=}2\),
\(\tin(4){=}3\), \(\tin(8){=}4\), \(\tin(5){=}5\), \(\tin(3){=}6\),
\(\tin(6){=}7\) and \(\tin(7){=}8\); subtree\((2)=\{2,4,8,5\}\) is therefore the
single contiguous range \(P[2..5]\) (\(\secref{c21-subtree}\)).  Depths are
\(\depth(1){=}0\), \(\depth(2){=}\depth(3){=}1\),
\(\depth(4){=}\depth(5){=}\depth(6){=}\depth(7){=}2\) and \(\depth(8){=}3\).

For path queries we instead keep the \emph{doubled} Euler tour \(E\) (a node is
written when we enter \emph{and} when we return to it), each node appearing twice:
\[
  E=[\,1,2,4,8,8,4,5,5,2,3,6,6,7,7,3,1\,] ,
\]
with \(st(v)\) the first index of \(v\) and \(en(v)\) its second, listed for the
nodes in preorder order \((1,2,4,8,5,3,6,7)\):
\(st{=}(1,2,3,4,7,10,11,13)\) and \(en{=}(16,9,6,5,8,15,12,14)\)
(\(\secref{c21-path}\)).  Finally a \emph{colouring}
\(\mathit{val}(v)=2,1,3,1,2,3,1,2\) for \(v=1..8\) makes "how many distinct colours on
this range/path?" a genuinely non-invertible statistic -- the reason Mo, and not
a Fenwick prefix difference, is the right tool.

\begin{figure}[H]
\centering
\begin{tikzpicture}[diagram, scale=0.95, transform shape,
   tnode/.style={circle, draw=primaryblue!70, fill=primaryblue!8, line width=0.7pt,
                 minimum size=6.6mm, font=\scriptsize\bfseries, inner sep=1pt},
   sub/.style={circle, draw=difforange, fill=difforange!16, line width=1.0pt,
             minimum size=6.6mm, font=\scriptsize\bfseries, inner sep=1pt},
   val/.style={font=\tiny, text=warningred}]
  \node[tnode] (n1) at (0,0) {1};
  \node[sub]   (n2) at (-2.4,-1.5) {2};
  \node[tnode] (n3) at (2.4,-1.5) {3};
  \node[sub]   (n4) at (-3.6,-3.0) {4};
  \node[sub]   (n5) at (-1.4,-3.0) {5};
  \node[tnode] (n6) at (1.4,-3.0) {6};
  \node[tnode] (n7) at (3.6,-3.0) {7};
  \node[sub]   (n8) at (-3.6,-4.5) {8};
  \draw (n1)--(n2) (n1)--(n3) (n2)--(n4) (n2)--(n5) (n3)--(n6) (n3)--(n7) (n4)--(n8);
  \node[val] at (0.45,0.35) {\tt 2};   \node[val] at (-1.95,-1.15) {\tt 1};
  \node[val] at (2.85,-1.15) {\tt 3};  \node[val] at (-4.05,-2.65) {\tt 1};
  \node[val] at (-0.95,-2.65) {\tt 2}; \node[val] at (1.85,-2.65) {\tt 3};
  \node[val] at (4.05,-2.65) {\tt 1};  \node[val] at (-4.05,-4.15) {\tt 2};
  \node[annot, font=\scriptsize, anchor=north west] at (-4.9,0.9)
    {\tt val(v) in red; orange = subtree(2) \{2,4,8,5\}=P[2..5]};
\end{tikzpicture}
\caption{The running example \(R\) (8 nodes, root 1).  Orange nodes are the
subtree of \(2\), which the preorder flatten makes the single contiguous range
\(P[\tin(2){=}2\,..\,\tout(2){=}5]\) (\(\secref{c21-subtree}\)).  Red labels are
\(\mathit{val}(v)\), the colour whose ``distinct count on a range/path'' statistic has no
inverse -- the whole reason this chapter exists.}
\label{fig:c21-tree}
\end{figure}

% ---------------- Which job, cost chart ----------------
\subsection*{Which Job? (the question-shape chart)}

\begin{center}
\small
\renewcommand{\arraystretch}{1.25}
\setlength{\tabcolsep}{3pt}
\begin{tabular}{@{}p{0.31\linewidth} p{0.17\linewidth} p{0.28\linewidth} >{\raggedright\arraybackslash}p{0.14\linewidth}@{}}
\toprule
\textbf{Question} & \textbf{Technique} & \textbf{Why not a Fenwick/segment prefix} & \textbf{Where} \\
\midrule
count subarrays with \(\mathit{xor}=k\), offline & array Mo, toggle prefix-xor freq & ``equal to \(k\) among freq of xors'' has no per-index inverse & \secref{c21-moengine} \\
distinct colours in a subtree, offline & flatten + Mo & \(\sum\!\mathrm{distinct}\) is not pref\((r)-\)pref\((l{-}1)\) & \secref{c21-subtree} \\
add to a node, sign-flip down its subtree & preorder + parity BIT & this one \emph{is} invertible: it goes to Fenwick & \secref{c21-subtree} \\
distinct colours on a path \(u\)-\(v\), offline & doubled tour + toggle & a path is two root-paths minus overlap; ranges do not reach it & \secref{c21-path} \\
answer \(q\) LCAs offline, no lifting table & Tarjan DFS+DSU & lifting is per-query \(O(\log n)\); Tarjan is \(O(\alpha)\) batch & \secref{c21-tarjan} \\
an \emph{online} path range with updates & \emph{go to HLD / LCT} & Mo is offline only; a moving/queried-now model breaks it & \chref{c16} \\
\bottomrule
\end{tabular}
\end{center}

\begin{insightbox}[One promise decides everything: reorder so neighbours differ by one]
A statistic with no inverse (distinct count, \(\sum \mathrm{freq}^2\), "how many
values appear \(\ge l\) times") cannot be answered from two prefixes, but
\emph{add/remove of one element} is cheap.  Mo's whole idea is to exploit that
asymmetry: instead of \(q\) independent answers, walk a window across the array
visiting every query range, ordering them so that the window moves as little as
possible.  Block ordering bounds the total movement to \(\bigO{(n+q)\sqrt n}\).
Every tree version is the same sweep on a smarter array: the preorder flatten
(subtrees), the doubled Euler tour with a \emph{toggle} (paths).  The one
chapter idea: \emph{when you can add cheaply but cannot subtract, change the
\emph{order} of the questions so you only ever add or drop one node.}
\end{insightbox}

% ---------------- Taxonomy chart ----------------
\begin{figure}[H]
\centering
\begin{tikzpicture}[diagram, scale=0.95, transform shape,
    mnode/.style={circle, draw=#1, line width=0.9pt, fill=#1!8, minimum size=9mm,
                  font=\scriptsize\bfseries, text=#1!60!black}]
  \node[mnode=primaryblue] (a) at (0,0)   {order};
  \node[annot, font=\scriptsize, anchor=west, text width=9.8cm, align=left] at (1.7,0) {block-sort the q offline ranges by (left block, right); a sweep answers all of them. B approx sqrt n, total O((n+q) sqrt n) (27.1)};
  \node[mnode=successgreen] (b) at (0,-1.1) {flatten};
  \node[annot, font=\scriptsize, anchor=west, text width=9.8cm, align=left] at (1.7,-1.1) {a subtree(v) is the contiguous preorder range [tin, tout]; run the same block sweep on it -- and when the queried sets are root paths, the DFS stack itself is the window (27.2)};
  \node[mnode=difforange]   (c) at (0,-2.2) {toggle};
  \node[annot, font=\scriptsize, anchor=west, text width=9.8cm, align=left] at (1.7,-2.2) {a path is not contiguous, so use the doubled Euler tour E[1..2n]; a node is in the path iff it appears exactly once in the current window; split at lca and add it by hand when neither endpoint is it (27.3)};
  \node[mnode=softpurple]   (d) at (0,-3.3) {offline};
  \node[annot, font=\scriptsize, anchor=west, text width=9.8cm, align=left] at (1.7,-3.3) {for LCA-only queries, Tarjan's single DFS with DSU answers all q in O((n+q) alpha) -- a WHITE/GREY/BLACK colouring, ancestor[find(u)] at the moment the second endpoint turns BLACK (27.4)};
  \node[mnode=warningred]   (e) at (0,-4.4) {reject};
  \node[annot, font=\scriptsize, anchor=west, text width=9.9cm, align=left] at (1.7,-4.4) {invertible range -> Fenwick (Ch.5/10); online path -> HLD (Ch.22) / LCT (Ch.25); per-set fold -> virtual tree (Ch.26); small marked set -> DSU-on-tree (Ch.17); all-pairs distance -> centroid (Ch.24) (27.5)};
\end{tikzpicture}
\caption{The chapter's five-state taxonomy.  \emph{order} is the array Mo sweep;
\emph{flatten} lifts it to subtrees; \emph{toggle} lifts it to paths via the
doubled tour; \emph{offline} is Tarjan's batched LCA; \emph{reject} hands the job
to the right neighbour.  Each solved problem sits on one row: CF 617E on
\emph{order}, CF 463E riding the recursion stack beside \emph{flatten}, the
Summit CF 1585E on \emph{toggle}+a frequency bucket.}
\label{fig:c21-taxonomy}
\end{figure}

% ---------------- Who owns what ----------------
\subsection*{Who Owns What (so nothing is learned twice)}

\begin{center}
\small
\renewcommand{\arraystretch}{1.25}
\setlength{\tabcolsep}{3pt}
\begin{tabular}{@{}p{0.28\linewidth} p{0.24\linewidth} p{0.44\linewidth}@{}}
\toprule
\textbf{Idea} & \textbf{Owner} & \textbf{This chapter's share} \\
\midrule
\(\tin/\tout\), subtree-as-range flatten & Ch.~5 (\chref{c24}) & the range Mo sweeps over; the single-visit flatten is reused, not re-taught \\
\(\lca\), \(\dist\), binary lifting & Ch.~16 (\chref{c11}) & 27.3 calls \(\lca\) to split a path; 27.4 offers Tarjan as an offline \emph{alternative}, never re-deriving lifting \\
invertible range add / point query & Ch.~5 / Ch.~10 Fenwick & 27.2 points back to it (Ch.~18 solves CF 383C with it, \chref{c14}); 27.5 routes every invertible query back to it \\
small-to-large / DSU on tree & Ch.~17 (\chref{c28}) & the rival offline answer for subtree stats (owns CF 375D); 27.5 draws the Mo-vs-DSU line \\
online path / subtree ranges & Ch.~22 HLD (\chref{c16}) & explicitly out; Mo is offline, a live path query with updates is HLD's job \\
per-marked-set folds on a fixed tree & Ch.~26 virtual tree (\chref{c30}) & the previous chapter hands \emph{overlapping} offline node-set queries to Mo \\
\bottomrule
\end{tabular}
\end{center}

% ---------------- Recognition guide ----------------
\begin{recognitionbox}[Pattern card -- ``a big offline batch of subtree/path queries, no inverse'']
\reckey{statements that fix one tree and hand you a \emph{large batch} of queries,
each asking a statistic over a subtree or a path, where the statistic has
\emph{no inverse}: number of distinct colours, \(\sum\mathrm{freq}^2\), ``how many
values occur \(\ge l\) times'', ``the \(k\)-th most frequent''.  The tell is
\(\sum\) over the ranges does not telescope, but adding or removing a single node
changes the answer in \(O(1)\).  All queries are known up front (offline) and no
query depends on the answer to an earlier one.}
\recuse{
  \item distinct-colour / frequency statistics over subtree ranges -- flatten + Mo (\(\secref{c21-subtree}\));
  \item the same over paths -- doubled Euler tour + toggle (\(\secref{c21-path}\));
  \item a pure batch of \(q\) LCA queries with no updates -- Tarjan offline (\(\secref{c21-tarjan}\)).
}
\recnotuse{
  \item the statistic \emph{does} have an inverse (a sum, an xor, a count of one fixed value) -- Euler + Fenwick / prefix array (\chref{c24});
  \item queries are \emph{online} or interleaved with updates -- HLD (\chref{c16}) or link-cut (\chref{c23});
  \item each query names a tiny marked set to fold over -- virtual tree (\chref{c30});
  \item you recompute a per-node subtree value for the \emph{whole} tree once -- tree DP / DSU-on-tree (\chref{c10}, \chref{c28});
  \item the answer is a distance distribution over all pairs -- centroid (\chref{c18}).
}
One-line test: \emph{offline batch + add/remove is cheap + the statistic does not
telescope -- then Mo; if it telescopes, use Fenwick, if it is online, use HLD.}
\end{recognitionbox}

% ---------------- Constraint -> technique ----------------
\subsection*{Reading the Constraints (which technique the limits want)}

\begin{center}
\small
\renewcommand{\arraystretch}{1.25}
\setlength{\tabcolsep}{3pt}
\begin{tabular}{@{}p{0.34\linewidth} p{0.40\linewidth} p{0.20\linewidth}@{}}
\toprule
\textbf{Constraint signature} & \textbf{What it rules in/out} & \textbf{Technique} \\
\midrule
\(n,q\le10^5\), queries known up front, ask "distinct / \(k\)-th / \(\ge l\)" over a range & telescoping prefix fails; a per-query fold is \(\bigO{nq}\) & \secref{c21-moengine} \\
subtree of \(v\), offline, non-invertible statistic & the subtree is a contiguous range once flattened & \secref{c21-subtree} \\
path \(u\)-\(v\), offline, non-invertible statistic & not contiguous -- doubled tour + toggle & \secref{c21-path} \\
\(q\) LCA queries, offline, no updates, want linear-ish total & no lifting table needed; one DFS+DSU & \secref{c21-tarjan} \\
updates interleaved, or queries must be answered online & Mo needs to see all queries first & \chref{c16} (HLD) \\
\bottomrule
\end{tabular}
\end{center}

% ---------------- Notation ----------------
\subsection*{Notation Used Only Here}

The fixed tree is \(R=(V,E)\), \(|V|=n\), rooted at \(1\).  \(P[\,]\) is the
preorder flatten (\(P[\tin(v)]=v\)); subtree\((v)\) is the range
\([\tin(v),\tout(v)]\).  \(E[\,]\) is the \emph{doubled} Euler tour of length
\(2n\), with \(st(v)\) and \(en(v)\) the two positions where \(v\) appears
(\(st\) on entry, \(en\) on exit).  A Mo window is the index interval
\([L,R]\) of the current array (either \(P\) or \(E\)); \texttt{add}/\texttt{rm}
are the \(O(1)\) steps that insert or delete the element at one end.  \(\mathit{val}(v)\)
is the colour/value at node \(v\); \(\mathit{cnt}[c]\) is how many times colour
\(c\) currently appears in the window; a \(\mathit{resp}\) structure (e.g.\
\(\mathit{bucket}[f]\)) holds the derived statistic so \texttt{add}/\texttt{rm}
stay \(O(1)\).  \(B\) is the Mo block size.  \(\alpha(\cdot)\) is the inverse
Ackermann function (Tarjan's DSU).  \(\depth,\tin,\tout,\lca,\dist\) are
Ch.~5/Ch.~16's, used unchanged.  Every section runs on the \(R\), \(P\), \(E\),
and \(\mathit{val}\) above; if a table here and a trace there disagree, this opener is
authoritative.

% =====================================================================
%  27.1  The Mo engine -- offline range queries on an ARRAY with a
%  sqrt-decomposition block order.  Exports: sec:c21-moengine
%  def:c21-mo  thm:c21-complexity  lst:c21-mo  prob:c21-xor  lst:c21-xor
%  Solved CF 617E "XOR and Favorite Number" (2200, L4), live-verified against
%  the CF problem page + problemset API (name, rating 2200, tag data structures;
%  both official samples).  DISCHARGES the Ch.26 (sec:c30-boundary) forward hand-
%  off "many overlapping offline node-set queries with a non-invertible statistic
%  -> Mo (Ch.27)".  This is the array engine; subtrees (27.2) and paths (27.3)
%  reuse it verbatim on a smarter array.
% =====================================================================
\section{Mo's Algorithm: The Block-Order Sweep}\label{sec:c21-moengine}

\begin{previouslybox}[Ch.~26 handed you the case this engine is built for]
Chapter~26 (\(\secref{c30-boundary}\)) closed by admitting a problem shape it
could not compress: \emph{one fixed tree, an offline batch of node-set queries, a
slowly-updated statistic with no inverse}.  Its pick table routed that shape here,
to Mo.  The reason a virtual tree (or a Fenwick prefix) fails on it is arithmetic:
a statistic like "how many distinct colours" does not telescope, so you cannot get
the answer for a range from two precomputed prefix values.  But \emph{adding or
removing one element} from a range \emph{is} cheap.  Mo's algorithm exploits that
asymmetry, and it does so on a plain array first -- this section, with subtrees and
paths layered on in 27.2--27.3.
\end{previouslybox}

\begin{topicintro}
You are given an array of length \(n\) and \(q\) offline range queries
\([L_i,R_i]\).  Maintain a sliding window \([curL,curR]\) with a statistic you can
update in \(O(1)\) when the window grows or shrinks by one at either end
(\texttt{add}/\texttt{rm}).  Naively jumping between arbitrary query ranges costs
\(\bigO{nq}\).  Mo's \term{block order} fixes it: sort queries by
\((\lfloor L/B\rfloor, R)\) for a block size \(B\), and the two pointers only ever
creep, never teleport.  With \(B\approx\sqrt n\) the total pointer movement -- hence
the total \texttt{add}/\texttt{rm} calls -- is \(\bigO{(n+q)\sqrt n}\)
(\(\thmref{c21-complexity}\)).  That single sweep answers all \(q\) queries.

\textbf{Prerequisites.} Nothing from earlier chapters is required for the engine
itself; you meet \(\tin/\tout\) (\chref{c24}) and \(\lca\) (\chref{c11}) in 27.2--27.3.
What you need here is comfort with \(O(1)\) amortised updates and a counting array.

\textbf{After this section you can}
\begin{itemize}[nosep]
  \item state the Mo block order and prove the \(\bigO{(n+q)\sqrt n}\) bound; \hfill\lv{4}
  \item write the two-pointer \texttt{add}/\texttt{rm} sweep cold; \hfill\lv{4}
  \item design an \(O(1)\) update for a non-telescoping statistic; \hfill\lv{4}
  \item count subarrays with a given xor offline -- solved CF 617E. \hfill\lv{4}
\end{itemize}
\end{topicintro}

% ---------------------------------------------------------------
\subsection{The block order and why it is fast}
% ---------------------------------------------------------------

\begin{definition}{Mo's algorithm (block order + two pointers)}{c21-mo}
Let every query be a range \([L_i,R_i]\) over an array of size \(n\), answered
offline.  Pick a block size \(B\) and \term{sort the queries} by the key
\((\lfloor L_i/B\rfloor,\ R_i)\) (within an odd block the \(R\) order is often
reversed -- the Hilbert/odd-even refinement -- but the plain order already gives
the bound).  Sweep a single window \([curL,curR]\) from query to query in that
order, moving the ends one step at a time with an \(O(1)\) \texttt{add}/\texttt{rm},
recording the window's statistic as each query's answer.
\end{definition}

\begin{theorem}{Mo is \(\bigO{(n+q)\sqrt n}\) with \(B=\sqrt n\)}{c21-complexity}
With block size \(B\), the right pointer moves \(O(n\cdot n/B)\) steps in total and
the left pointer moves \(O(q\cdot B)\) steps, so the total work is
\(\bigO(n^2/B + qB) = \bigO{(n+q)\sqrt n}\) when \(B=\Theta(\sqrt n)\) (more
precisely \(B=n/\sqrt q\) balances the two terms).
\end{theorem}
\begin{proofbox}
Split the \(q\) queries by the block of their \(L\). \emph{Right pointer.} Inside
one \(L\)-block the queries are sorted by \(R\), so while you answer that block the
\(R\) pointer only marches forward, covering at most \(n\) positions per block; with
\(n/B\) blocks the total \(R\)-movement is \(O(n\cdot n/B)=O(n^2/B)\). \emph{Left
pointer.} When a query starts, \(curL\) jumps to the new \(L\); within one block all
\(L\) lie in an interval of length \(B\), so each intra-block left move is \(O(B)\),
and across \(q\) queries the left pointer moves \(O(qB)\) steps. Summing
\(O(n^2/B+qB)\) and minimising (AM-GM, or set the terms equal: \(n^2/B=qB
\Rightarrow B=n/\sqrt q\); for \(q=\Theta(n)\) that is \(B=\sqrt n\)) gives
\(\bigO{(n+q)\sqrt n}\). Each step is one \(O(1)\) \texttt{add}/\texttt{rm}. \(\Box\)
\end{proofbox}

\begin{insightbox}[Why the sort key is "block of \(L\), then \(R\)" and not the reverse]
You pay for the pointer that \emph{jumps}; you almost pay for the one that
\emph{sweeps}.  Sorting by \(R\) inside a block makes the \(R\) pointer monotone
(sweeps, total \(O(n)\) per block, no resets).  Grouping by the \emph{block of
\(L\)} keeps the \(L\) pointer local (its jumps are \(\le B\)).  Swap the two keys
and neither pointer stays cheap -- you get \(O(nq)\) again.  The magic is entirely
in "one pointer monotone per group, the other confined to a small group."
\end{insightbox}

\begin{lstlisting}[caption={Generic Mo skeleton: block-sort the queries, sweep one window with O(1) add/rm, read off the statistic. ASCII only, identifiers without underscores}, label={lst:c21-mo}]
struct Query { int l, r, idx; };
int B;
bool cmp(const Query& a, const Query& b){
    if (a.l / B != b.l / B) return a.l / B < b.l / B;   // primary: block of L
    return a.r < b.r;                                   // secondary: R (odd-even: reverse)
}
// curL,curR = current window ends; cur = statistic of the window
for (auto &qu : queries){                    // queries already sorted, cur starts EMPTY
    while (curL > qu.l) add(--curL);         // grow left
    while (curR < qu.r) add(++curR);         // grow right
    while (curL < qu.l) rm(curL++);          // shrink left
    while (curR > qu.r) rm(curR--);          // shrink right
    ans[qu.idx] = cur;
}
\end{lstlisting}

\begin{tracebox}[The block order on \(n{=}9\), \(B{=}3\), three queries \(\lbrack 7,8\rbrack, \lbrack 1,3\rbrack, \lbrack 2,5\rbrack\)]
Block of \(L\): \([1,3]\) and \([2,5]\) are in block \(0\) (indices \(1..3\)), \([7,8]\)
in block \(2\).  Sorted: \([1,3],[2,5]\) (block 0, by \(R\)), then \([7,8]\).  Sweep from
empty \([1,0]\): to \([1,3]\) add positions \(1,2,3\) (3 moves).  To \([2,5]\): shrink
left once (\(rm(1)\)), grow right twice (\(add(4),add(5)\)) -- 3 moves, \(R\) never went
backwards.  To \([7,8]\): grow right to \(8\) (4 adds), then shrink left to \(7\) (3 rm)
-- 7 moves.  Total \(13 = O((n+q)\sqrt n)\): the \(R\) pointer ran \(3\!\to\!5\!\to\!8\)
monotonically and \(L\) only jumped within/across small blocks. \(\Box\)
\end{tracebox}

\begin{complexitybox}[The engine]
Sorting the \(q\) queries is \(\bigO{q\log q}\); the sweep is \(\bigO{(n+q)\sqrt n}\)
\texttt{add}/\texttt{rm} calls at \(O(1)\) each (\(\thmref{c21-complexity}\)). With
\(q=\Theta(n)\) the total is \(\bigO{n\sqrt n}\). \textbf{Space} \(\bigO{n+q}\) for
the array, the query list, and whatever the statistic needs. The whole method lives
or dies on \texttt{add}/\texttt{rm} being \(O(1)\) -- if updating costs \(\bigO{\log n}\)
you have Mo-with-a-BIT, usually slower than a direct structure.
\end{complexitybox}

% ---------------------------------------------------------------
\subsection{Problem Lab: CF 617E XOR and Favorite Number -- the canonical Mo warm-up}
% ---------------------------------------------------------------

\solvedproblem{c21-xor}{XOR and Favorite Number}{CF 617E}
  {https://codeforces.com/problemset/problem/617/E}{CF 2200 L4}
  {offline range queries with a block order; prefix-xor turns "subarray xor = k" into "count pairs (i,j) in range with px[i] xor px[j] = k"; a freq-of-prefix-xor window with O(1) add/rm}

\begin{problemstatement}[XOR and Favorite Number (statement, constraints and both samples verified against the live CF page)]
An array \(a_1,\dots,a_n\) (\(1\le n\le 10^5\), \(0\le a_i\le 10^6\)) and an integer
\(k\) (\(0\le k\le 10^6\)).  There are \(m\) questions (\(1\le m\le 10^5\)); each gives
\(l,r\) (\(1\le l\le r\le n\)) and asks: \emph{how many pairs \((i,j)\) with
\(l\le i\le j\le r\) satisfy \(a_i \oplus a_{i+1}\oplus\cdots\oplus a_j = k\)?}
Answer the questions in order.  (The count of pairs blows past \texttt{int}.)

\textbf{Sample 1.} \(n{=}6,m{=}2,k{=}3\); \(a=[1,2,1,1,0,3]\); queries \((1,6),(3,5)\)
\(\to\) \texttt{7} \texttt{0}. \textbf{Sample 2.} \(n{=}5,m{=}3,k{=}1\);
\(a=[1,1,1,1,1]\); queries \((1,5),(2,4),(1,3)\) \(\to\) \texttt{9} \texttt{4} \texttt{4}.
\end{problemstatement}

\begin{attemptgate}[Attempt it yourself first -- \lv{4}, give yourself 30 minutes]
A subarray xor is a difference of two prefix xors.  Once you see that, what is a
query range actually counting, and why can a Fenwick prefix not answer it directly?
\peek{let \(px[t]=a_1\oplus\cdots\oplus a_t\) (and \(px[0]=0\)); then
\(a_i\oplus\cdots\oplus a_j = px[j]\oplus px[i{-}1]\).  A query \([l,r]\) becomes
"how many pairs \(p<q\) with \(l{-}1\le p<q\le r\) have \(px[p]\oplus px[q]=k\)?" --
a range over the \emph{prefix} array \(px[0..n]\).}
\nudge{This is the "no inverse" tell: counting equal-xor pairs in a range does not
telescope.  But adding one index to the current window is \(O(1)\): if the window
already holds \(cnt[x]\) copies of each prefix value, inserting a new \(px=v\) forms
exactly \(cnt[v\oplus k]\) new good pairs.}
\fullhint{Run Mo over ranges \([L,R]=[l-1,\ r]\) of \(px[0..n]\).  Keep
\(\mathit{cnt}[v]\) = multiplicity of value \(v\) in the window and \(cur\) = number
of pairs in the window with xor \(k\).  \texttt{add(t)}: \(v=px[t];\ cur+\!\!\!=cnt[v\oplus k];\ cnt[v]++\).
\texttt{rm(t)}: \(v=px[t];\ cnt[v]--;\ cur-=cnt[v\oplus k]\).  \texttt{add}/\texttt{rm}
are symmetric because pair-counting is order-independent.  \textbf{Start the window
empty with \(cnt[px[0]]=0\)} and be careful to add index \(0\) (the empty prefix) when
\(L=0\): it is \(px[0]=0\) that pairs with a prefix equal to \(k\) to count subarrays
starting at \(l\).  Read answers in \(O(1)\) from \(cur\).  \(cur\) is \texttt{long
long}; the value domain is \(<2^{20}\), size \(\mathit{cnt}\) accordingly.}
\end{attemptgate}

\solstep{Understanding}
Two reductions do the work.  First, prefix-xor moves the problem from
\(O(\text{length})\) subarray xors to counting pairs in a range -- so a query is now
just a window \([l-1,r]\) over \(px\) and its answer is a pair count.  Second, that
pair count does not telescope (there is no per-index contribution you can subtract),
which is exactly the Mo regime: cheap to insert one prefix value, expensive to answer
a fresh range from scratch.  The window statistic is a frequency table \(\mathit{cnt}\)
plus a running pair total \(cur\).

\solstep{Key Observations}
\begin{observation}
Inserting a prefix value \(v\) into a window that already contains \(\mathit{cnt}[x]\)
of each \(x\) creates \(cnt[v\oplus k]\) new pairs summing (under xor) to \(k\); so
\texttt{add} is \(O(1)\) and its mirror \texttt{rm} is \(O(1)\) (decrement first, then
subtract -- the pair being removed no longer counts itself).
\end{observation}
\begin{observation}
Because a pair is an unordered \(\{p,q\}\) with \(p<q\), the two Mo pointers may move
in any order: the window always holds a set of prefix indices and \(cur\) counts
pairs among them, independent of insertion order.  This is what makes Mo's pointer
freedom legal here.
\end{observation}

\solstep{The Approach Ladder}
\begin{center}\renewcommand{\arraystretch}{1.25}\setlength{\tabcolsep}{3pt}
\begin{tabular}{@{}c p{0.26\linewidth} p{0.30\linewidth} l c@{}}
\toprule
\textbf{\#} & \textbf{Approach} & \textbf{Idea} & \textbf{Time} & \textbf{Verdict}\\
\midrule
1 & \lv{2}\ per query, double loop & count all subarrays in \([l,r]\) directly & \(\bigO{m n^2}\) & \tle \\
2 & \lv{3}\ per query, prefix-xor map & one pass, freq map over \([l-1,r]\) & \(\bigO{mn}\) & \tle \\
3 & \lv{4}\ \textbf{Mo over prefix xors} & block order, \(O(1)\) add/rm & \(\bigO{(n+m)\sqrt n}\) & \ac \\
\bottomrule
\end{tabular}\end{center}
Rung~2 is \(O(n)\) \emph{per query} and \(m\) is up to \(n\), so it is \(O(n^2)\);
rung~3 shares the frequency table across consecutive queries, paying only for the
pointers that move.

\solstep{Solution}
Build \(px[0..n]\).  Turn each question \((l,r)\) into a Mo range \([l-1, r]\) tagged
with its original index.  Sort by block order.  Sweep one window maintaining
\(\mathit{cnt}[\cdot]\) over prefix values and the running pair count \(cur\);
\texttt{add}/\texttt{rm} update both in \(O(1)\) as in the hint.  Store \(cur\) into
\(\mathit{ans}[\text{idx}]\) and print answers in input order.

\solstep{Implementation}
\begin{lstlisting}[caption={CF 617E via Mo over the prefix-xor array. add(t): cur += cnt[px[t] xor k] then cnt[px[t]]++. rm(t): cnt[px[t]]-- then cur -= cnt[px[t] xor k]. Block size about n divided by sqrt m. ASCII only}, label={lst:c21-xor}]
#include <bits/stdc++.h>
using namespace std;
typedef long long ll;
const int MAXN=100005;
int n,m,k, px[MAXN], B;
ll cur, ans[MAXN];
int cnt[1<<20];
struct Q{ int l,r,id; } qs[MAXN];
bool cmp(const Q&a,const Q&b){
  if(a.l/B!=b.l/B) return a.l/B<b.l/B;
  return a.r<b.r;
}
void add(int t){ int v=px[t]; cur+=cnt[v^k]; cnt[v]++; }
void rm (int t){ int v=px[t]; cnt[v]--; cur-=cnt[v^k]; }
int main(){
  scanf("%d%d%d",&n,&m,&k);
  px[0]=0;
  for(int i=1;i<=n;i++){ int a; scanf("%d",&a); px[i]=px[i-1]^a; }
  for(int i=0;i<m;i++){ int l,r; scanf("%d%d",&l,&r); qs[i]={l-1,r,i}; }
  B = n / max(1,(int)sqrt((double)n));      // block size
  sort(qs,qs+m,cmp);
  int curL=0, curR=-1; cur=0; memset(cnt,0,sizeof cnt);
  for(int qi=0;qi<m;qi++){
    Q &q=qs[qi];
    while(curL>q.l) add(--curL);
    while(curR<q.r) add(++curR);
    while(curL<q.l) rm(curL++);
    while(curR>q.r) rm(curR--);
    ans[q.id]=cur;
  }
  for(int i=0;i<m;i++) printf("%lld\n",ans[i]);
}
\end{lstlisting}

\solstep{Trace}
\begin{tracebox}[Sample 1: prefix xors, then the pair count for query \((1,6)\)]
\(px=[0,1,3,2,3,3,0]\) (indices \(0..6\)).  Query \((1,6)\to\) window \([0,6]\), all
seven prefix values.  Frequencies: \(\mathit{cnt}[0]{=}2,\mathit{cnt}[1]{=}1,
\mathit{cnt}[2]{=}1,\mathit{cnt}[3]{=}3\).  Pairs with \(x\oplus y=3\): the
\((0,3)\) family \(=\mathit{cnt}[0]\!\cdot\!\mathit{cnt}[3]=2\!\cdot\!3=6\); the
\((1,2)\) family \(=\mathit{cnt}[1]\!\cdot\!\mathit{cnt}[2]=1\).  Total
\(\boxed{7}\) -- the window's \(cur\).  Query \((3,5)\to[2,5]\) holds \(px\) values
\(\{3,2,3,3\}\): no value pairs to \(3\) (needs \(0\) with \(3\), needs \(1\) with
\(2\)), \(cur=\boxed{0}\).  Both match the judge. \(\Box\)
\end{tracebox}

\begin{warningbox}[Three ways Mo goes wrong]
\begin{enumerate}[nosep,leftmargin=1.6em]
  \item \textbf{Forgetting the empty prefix \(px[0]\).} A subarray starting at \(l\)
        pairs \(px[l{-}1]\) with a later prefix; if the window's \(L=l-1\) and you never
        add index \(0\), you under-count every subarray that starts at the left edge.
  \item \textbf{\texttt{rm} updating in the wrong order.} \texttt{rm} must
        \emph{decrement} \(\mathit{cnt}[v]\) \emph{before} subtracting
        \(\mathit{cnt}[v\oplus k]\), or you subtract the pair that is itself leaving.
  \item \textbf{Answer overflow / value-domain sizing.} \(cur\) is a pair count
        (\(\le n^2/2\)) -- \texttt{long long}.  Size \(\mathit{cnt}\) to the prefix-xor
        domain (\(<2^{20}\) here), not to \(n\).
\end{enumerate}
\end{warningbox}

\begin{hardeningbox}
\begin{itemize}[nosep,leftmargin=1.4em]
  \item \textbf{Block size is a tuning knob, not a theorem.} \(B=n/\sqrt q\) balances
        the two pointer terms; when in doubt \(B\approx\sqrt n\) works, and a fixed
        \(B\approx 700\) for \(n=10^5\) is a common contest shortcut.
  \item \textbf{Odd-even \(R\) order cuts constants.} Reversing \(R\) in odd \(L\)-blocks
        shortens the \(R\)-jump at block boundaries; correctness is untouched.
  \item \textbf{Mo is offline \emph{and} no-updates.} If updates are interleaved, the
        block order is invalid -- that is a different chapter (HLD, Ch.~22).
\end{itemize}
\end{hardeningbox}

\begin{recallbox}[27.1 in one breath]
\emph{Block-sort the offline ranges by \((\lfloor L/B\rfloor,R)\), sweep one window,
move the ends with an \(O(1)\) \texttt{add}/\texttt{rm}, and read the statistic -- all
\(q\) answers in \(\bigO{(n+q)\sqrt n}\) (\(\thmref{c21-complexity}\),
\(\coderef{c21-mo}\)).  CF 617E is the engine on prefix xors: a range is a pair-count
over \([l-1,r]\) of \(px\), inserting \(v\) forms \(cnt[v\oplus k]\) pairs, and the
sweep replays the samples to \(7,0\) and \(9,4,4\) (\(\probref{c21-xor}\)).  Subtrees
and paths reuse this \emph{exact} loop -- only the array changes.}
\selftest{
  \item Why sort by (block of \(L\), then \(R\)) and not the other way? \emph{(To keep one pointer monotone per group and the other confined to a width-\(B\) group; the bound \(O(n^2/B+qB)\) needs it. \(\thmref{c21-complexity}\).)}
  \item What makes CF 617E a Mo problem and not a Fenwick one? \emph{(The equal-xor pair count has no inverse/telescoping form, but a single insert/delete is \(O(1)\).)}
  \item Give the \texttt{add} step for the prefix-xor window. \emph{(\(cur += cnt[v\oplus k];\ cnt[v]++\). \(\coderef{c21-xor}\).)}
}
\end{recallbox}

% =====================================================================
%  27.2  Mo on subtrees -- the preorder flatten makes subtree(v) the
%  contiguous range [tin(v), tout(v)], so 27.1's engine runs verbatim;
%  plus the solved root-path batch CF 463E, where the window is the DFS
%  stack itself and the <=50 updates legalise one rebuild per block.
%  Exports: sec:c21-subtree  def:c21-flatten  thm:c21-subtree-range
%  lst:c21-flatten  prob:c21-caisa  lst:c21-caisa
%  Solved CF 463E "Caisa and Tree" (2100, L4), live-verified against the
%  CF problem page (name, root-1 strict-ancestor gcd query, type-2
%  updates capped at 50, n,q <= 100000, a_i <= 2*10^6; official sample
%  hand-replayed to -1 1 2 -1 1 BEFORE writing).
%  CF 383C (parity + difference BIT) is PROSE-ONLY here -- it is Ch.18's
%  own solved template (\chref{c14}); CF 375D's URL lives in Ch.17.
%  No URL is re-linked; urlsweep stays clean.
% =====================================================================
\section{Mo on Subtrees: the Preorder Flatten}\label{sec:c21-subtree}

\begin{previouslybox}[27.1 handed you the sweep; Ch.~5 handed you the array]
\secref{c21-moengine} built the engine: block-sort offline ranges, creep two
pointers, keep a statistic with \(O(1)\) \texttt{add}/\texttt{rm}, and every
range is answered in one \(\bigO{(n+q)\sqrt n}\) sweep.  What it never needed
was a tree -- so this section bolts the tree on, in the cheapest possible way.
Chapter~5 (\chref{c24}) already showed that one DFS fixes \(\tin(v),\tout(v)\),
and that the preorder array \(P\) with \(P[\tin(v)]=v\) turns \emph{every}
subtree into a contiguous segment.  A subtree query is therefore \emph{already}
an array range query: hand it to the 27.1 sweep unchanged.  And when the
statistic \emph{does} telescope (the sign-flip subtree update Ch.~18 solves as
CF 383C, \chref{c14}), you skip Mo entirely and spend a Fenwick difference
array instead.  The solved CF 463E below opens a \emph{third} door: the queried
set is not a subtree range but the current root path -- and a plain DFS carries
that window for free.  Knowing which world you live in is half the section.
\end{previouslybox}

\begin{topicintro}
Two lemmas do all the work.  \defref{c21-flatten} recalls the flatten: the DFS
entry/exit pair \([\tin(v),\tout(v)]\) is exactly the index span of
subtree\((v)\) inside preorder \(P\).  \thmref{c21-subtree-range} upgrades it:
running the \secref{c21-moengine} sweep on ranges \([\tin(v),\tout(v)]\) answers
\(q\) offline subtree statistics in \(\bigO{(n+q)\sqrt n}\) -- the engine does
not care that the array was baked from a tree.  The \emph{worked} running
example flattens \(R\) and reads subtree\((2)=P[2..5]\) off one line.  The
\solvedproblem{c21-caisa}{Caisa and Tree}{CF 463E}
  {https://codeforces.com/problemset/problem/463/E}{CF 2100 L4}
  {sieve of smallest prime factors; the DFS stack as the root-path window; a per-prime deepest-ancestor table with \(O(1)\) push/undo; one rebuild per block of the rare (\(\le 50\)) updates}
then opens the third world: "deepest strict ancestor sharing a prime factor
with \(v\)" has no inverse and no monotone lifting predicate -- but the queried
root paths are the recursion stack itself, so a plain DFS carries the window,
and the promise of at most \(50\) updates makes one rebuild per query block the
cheapest move on the board.

\textbf{Prerequisites.} \secref{c21-moengine} (the sweep, \texttt{add}/\texttt{rm});
\chref{c24} (\(\tin/\tout\), the flatten); \chref{c14} (Ch.~18, Fenwick trees --
where the invertible subtree-update template CF 383C is fully solved).

\textbf{After this section you can}
\begin{itemize}[nosep]
  \item flatten any rooted tree and name the range \([\tin(v),\tout(v)]\) of a subtree; \hfill\lv{3}
  \item run Mo on subtree ranges for a non-invertible statistic (distinct colours); \hfill\lv{4}
  \item carry a root-path statistic on the DFS stack with an \(O(1)\) push/undo table; \hfill\lv{4}
  \item solve CF 463E cold -- solved, traced to \(-1,1,2,-1,1\). \hfill\lv{4}
\end{itemize}
\end{topicintro}

% ---------------------------------------------------------------
\subsection{The flatten, and Mo rides it}
% ---------------------------------------------------------------

\begin{definition}{Preorder flatten and the subtree span}{c21-flatten}
Root the tree, run the left-to-right DFS, and timestamp \(\tin(v)\) = arrival
order (\(1\)..\(n\)) and \(\tout(v)\) = the last arrival inside \(v\)'s region.
Write \(P[\tin(v)]=v\).  Then
\[
  \subtree(v) \;=\; \{\,u : \tin(v)\le \tin(u)\le \tout(v)\,\},
\]
so \emph{aggregate over subtree\((v)\)} \(=\) \emph{aggregate over the array
range} \(P[\tin(v)\,..\,\tout(v)]\).
\end{definition}

\begin{theorem}{Subtree Mo is one flattened range sweep}{c21-subtree-range}
Given \(q\) offline queries of the form "statistic of subtree\((v_i)\)" where
the statistic supports \(O(1)\) element insert/delete but has no closed prefix
form, the block-order sweep of \secref{c21-moengine} over the ranges
\([\tin(v_i),\tout(v_i)]\) on array \(P\) answers all \(q\) in
\(\bigO{(n+q)\sqrt n}\).
\end{theorem}
\begin{proofbox}
The DFS visits every node of subtree\((v)\) between entering \(v\) and leaving
its region, and visits nothing else in between -- so subtree\((v)\) occupies
exactly the index span \([\tin(v),\tout(v)]\) of \(P\).  Under that bijection
node \(u\) is the single array position \(\tin(u)\): \texttt{add}/\texttt{rm}
touch one node of the current subtree-window per step.  The \emph{shape} of the
ranges is irrelevant to \thmref{c21-complexity} -- it bounded pointer movement
for any family of \(q\) ranges over an \(n\)-array.  Applying it verbatim gives
\(\bigO{(n+q)\sqrt n}\). \(\Box\)
\end{proofbox}

\begin{lstlisting}[caption={The flatten DFS plus the Mo-on-subtrees wiring: a query on subtree(v) is just the range tin[v] to tout[v] fed to the 27.1 block-order sweep. ASCII only, no underscores in identifiers}, label={lst:c21-flatten}]
int timer_, tin[MAXN], tout[MAXN], P[MAXN], dep[MAXN];
void dfs(int u, int p){
    tin[u] = ++timer_;  P[timer_] = u;
    for (int v : adj[u]) if (v != p){
        dep[v] = dep[u] + 1;
        dfs(v, u);
    }
    tout[u] = timer_;                    // subtree(u) = P[tin[u] .. tout[u]]
}
// Mo-on-subtrees: query i on subtree(v) becomes the range
//   l = tin[v],  r = tout[v]
// then sort by (l/B, r) and sweep with the add()/rm() of listing lst:c21-mo
// applied to P[pos] (the node at that preorder position).
\end{lstlisting}

\begin{tracebox}[Running example: subtree\((2)\) of \(R\) is one line of the flatten]
On \(R\) (\(\figref{c21-tree}\)) the DFS gives \(P=[1,2,4,8,5,3,6,7]\),
\(\tin(v)=(1,2,3,4,5,6,7,8)\) and \(\tout(v)=(8,5,6,4,5,7,8,4)\) for nodes
\(v=1..8\); in particular \(\tout(2)=5\).  subtree\((2)=\{2,4,8,5\}\) and indeed
\(P[2..5]=[2,4,8,5]\): a \emph{contiguous} range, though the four nodes are
scattered through the DFS.  A distinct-colour Mo would set
\(L,R\) targets inside \([2,5]\); the window's \texttt{add} of position \(i\)
inserts colour \(\mathit{val}(P[i])\): with \(\mathit{val}=(2,1,3,1,2,3,1,2)\)
from the opener, positions \(2,3,4,5\) hold nodes \(2,4,8,5\) with colours
\(1,3,2,2\) -- three distinct.  Nothing about the sweep changed; only the array
did. \(\Box\)
\end{tracebox}

\begin{insightbox}[The two worlds, decided by one question]
Ask: \emph{"can I get the answer on a range from two precomputed prefixes?"}
If yes (sums, xors, signed parity-corrected sums), a \(\bigO{\log n}\) BIT/
segment tree over the flatten crushes the problem and Mo is pure overhead.  If
no (distinct count, \(\sum \mathrm{freq}^2\), "how many values appear \(\ge l\)
times"), prefixes do not exist, but one-element \texttt{add}/\texttt{rm} do --
and that asymmetry is exactly the block-order sweep's salary.  CF~383C --
Ch.~18's solved case (\chref{c14}) -- is the study in spotting that a
\emph{scary-looking} update (sign flips with depth) lives in the first world.
CF~463E, solved below, lives in a \emph{third}: no inverse exists, yet no
reordering is needed either, because every queried root path is the current
recursion stack -- a nested window that a plain DFS sweeps for free.
\end{insightbox}

\begin{complexitybox}[Subtree Mo]
Flatten: \(\bigO{n}\) one DFS.  Per query: form its range in \(O(1)\).  Sweep:
\(\bigO{(n+q)\sqrt n}\) (\(\thmref{c21-subtree-range}\)).  Contrast: an
invertible statistic on the same flatten costs \(\bigO{(n+q)\log n}\) with a
difference BIT -- so \emph{never pay \(\sqrt n\) when a prefix exists}.  The
canonical non-invertible drill is CF~375D ("Tree and Queries": distinct colours
in subtrees, answered here by subtree ranges, in Ch.~17 by DSU-on-tree -- the
URL and full treatment live in \chref{c28}).
\end{complexitybox}

% ---------------------------------------------------------------
\subsection{Problem Lab: CF 463E Caisa and Tree -- the root-path stack and rare rebuilds}
% ---------------------------------------------------------------

\begin{problemstatement}[Caisa and Tree (statement, constraints and sample verified against the live CF page)]
A tree on \(n\le 100\,000\) nodes rooted at \(1\); node \(i\) has value
\(a_i\le 2\cdot 10^6\).  Process \(q\le 100\,000\) queries in order, two kinds:
\begin{itemize}[nosep]
  \item \texttt{1 v} -- write the root-to-\(v\) path as
        \(u_1,\dots,u_k\) (\(u_1=1\), \(u_k=v\)); output the vertex \(u_i\) with
        \(i<k\) and \(\gcd(a_{u_i}, a_v) > 1\) having the \emph{maximum} such
        \(i\); output \(-1\) if none exists;
  \item \texttt{2 v w} -- change \(a_v\) to \(w\).
\end{itemize}
The promise that shapes everything: \emph{at most \(50\) queries are of type 2}.

\textbf{Sample.} \(n{=}4,\ q{=}6\), \(a=[10,8,4,3]\), chain edges
\(1\!-\!2,\ 2\!-\!3,\ 3\!-\!4\); queries \(\texttt{1 1},\ \texttt{1 2},\
\texttt{1 3},\ \texttt{1 4},\ \texttt{2 1 9},\ \texttt{1 4}\) \(\to\)
\texttt{-1 1 2 -1 1}.
\end{problemstatement}

\begin{attemptgate}[Attempt it yourself first -- \lv{4}, give yourself 25 minutes]
\(\gcd(a_u,a_v)>1\) means "share \emph{some} prime".  You want the deepest such
ancestor for many \(v\), and the values barely change.  Can you avoid both
climbing per query and a lifting table?
\peek{Two numbers share a prime iff some prime \(p\mid\) both.  Factor \(a_v\)
into \emph{distinct} primes -- at most seven below \(2\cdot 10^6\).}
\nudge{For every prime \(p\) keep \(\mathit{last}[p]\): the deepest node on the
\emph{current root path} divisible by \(p\).  A DFS enters and leaves nodes, so
\(\mathit{last}\) is push-then-undo -- the \(O(1)\) \texttt{add}/\texttt{rm} of
27.1 again, this time riding the recursion stack instead of two pointers.}
\fullhint{Read \(v\)'s query \emph{before} pushing \(v\)'s primes -- the target
must be a strict ancestor (\(i<k\)).  Between two type-2 queries all values are
static, so cut the stream into \(\le 51\) blocks, run one DFS per block, apply
the single update, repeat.  Sieve smallest-prime-factors once up to
\(2\cdot 10^6\).  Cost: \(51\cdot n\) enters with \(\le 7\) table probes each
-- inside the 10-second limit with room to spare.  No Mo, no lifting: the
rebuild schedule \emph{is} the offline trade.}
\end{attemptgate}

\solstep{Understanding}
Two subproblems, each with a cheap answer.  (i) \emph{The statistic}:
"\(\gcd(a_u,a_v)>1\)" is existential over primes -- the deepest ancestor
divisible by \(p\), for some prime \(p\) of \(a_v\).  It is \emph{not} monotone
along the path (a grandparent may share a prime the parent does not), so binary
lifting has nothing to bite on -- that is ladder rung 2's WA.  A per-prime
\emph{last-seen} table answers it with at most seven probes.  (ii) \emph{The
schedule}: interleaved updates normally kill any offline sweep -- except there
are at most \(50\), so the query stream splits into \(\le 51\) static blocks and
a full re-DFS per block costs almost nothing compared to any sqrt-decomposition.

\solstep{Key Observations}
\begin{observation}
Every \(a\le 2\cdot 10^6\) has at most \(7\) distinct prime factors
(\(2{\cdot}3{\cdot}5{\cdot}7{\cdot}11{\cdot}13{\cdot}17 = 510510\), and
multiplying by the eighth prime exceeds \(2\cdot 10^6\)), so a factor list is a
constant-size key: \(O(1)\) probes and \(O(1)\) undos per node, always.
\end{observation}
\begin{observation}
During a DFS from the root the call stack \emph{is} the current root path:
writing \(\mathit{last}[p]=u\) on entry to \(u\) and restoring the saved value
on exit maintains "deepest ancestor divisible by \(p\)" for \emph{every} prime
at once.  Attach each type-1 query to its node and answer it at entry -- before
the node's own primes are pushed, so the match stays a strict ancestor.
\end{observation}

\solstep{The Approach Ladder}
\begin{center}\renewcommand{\arraystretch}{1.25}\setlength{\tabcolsep}{3pt}
\begin{tabular}{@{}c p{0.26\linewidth} p{0.30\linewidth} l c@{}}
\toprule
\textbf{\#} & \textbf{Approach} & \textbf{Idea} & \textbf{Time} & \textbf{Verdict}\\
\midrule
1 & \lv{2}\ climb the path per query & test \(\gcd(a_u,a_v)>1\) ancestor by ancestor & \(\bigO{qn}\) chain & \tle \\
2 & \lv{3}\ binary lifting on the path & lift to "deepest ancestor sharing a prime" & \(\bigO{q\log n}\) hope & \textbf{WA} \\
3 & \lv{4}\ \textbf{DFS stack + \(\mathit{last}[p]\) table} & one re-DFS per static block, \(O(1)\) push/undo & \(\bigO{51\,n\cdot 7+q}\) & \ac \\
\bottomrule
\end{tabular}\end{center}
Rung~2 is the trap the statement sets on purpose: binary lifting needs a
\emph{monotone} predicate, and "shares a prime with \(v\)" is not one -- a
grandparent can match where the parent fails, so the lift overshoots or
undershoots and no fix-up patch makes it right.  Rung~1 is honest but quadratic:
on the chain with \(n=q=10^5\) it is \(10^{10}\) gcds.  Rung~3 collects the two
gifts the statement hands you -- the queried window is always the current root
path (free to carry on the recursion stack) and updates number \(\le 50\) (so
the static-block rebuild is just \(51\) DFSes).

\solstep{Solution}
Sieve smallest-prime-factors \(\mathit{spf}[\cdot]\) up to \(2\cdot10^6\) once.
Cut the query stream into its \(\le 51\) maximal runs of type-1 queries
separated by type-2 updates.  For each run: attach every query to its vertex,
then DFS from \(1\) maintaining \(\mathit{last}[p]\) -- the deepest node on the
current stack divisible by prime \(p\).  On entry to \(u\): factor \(a_u\) into
\emph{distinct} primes; answer each query attached to \(u\) as the deepest
\(\mathit{last}[p]\) over those primes, \emph{read before} \(u\)'s own primes
are pushed (the match must be a strict ancestor); then save and overwrite
\(\mathit{last}[p]\leftarrow u\), recurse, and restore the saved values on
exit.  Print the block's answers in original order, apply the one update, and
start the next DFS.

\solstep{Implementation}
\begin{lstlisting}[caption={CF 463E: one DFS per static block. The recursion stack carries the root path; lastp[p] is the deepest stack node divisible by prime p, saved and undone around each node; queries are read before the push. ASCII only}, label={lst:c21-caisa}]
#include <bits/stdc++.h>
using namespace std;
const int MAXN=100005, MAXA=2000005;
int n,q,a[MAXN],spf[MAXA],dep[MAXN],ans[MAXN];
int typ[MAXN],qv[MAXN],qw[MAXN],lastp[MAXA];
vector<int> adj[MAXN],ask[MAXN];

void build_spf(){
  for(int i=2;i<MAXA;i++){
    if(!spf[i]){
      spf[i]=i;
      if((long long)i*i<MAXA)
        for(long long j=(long long)i*i;j<MAXA;j+=i)
          if(!spf[j]) spf[j]=i;
    }
  }
}
void factor(int x, vector<int>& out){      // distinct primes only
  out.clear();
  while(x>1){ int p=spf[x]; out.push_back(p); while(x%p==0) x/=p; }
}
void dfs(int u,int p){
  vector<int> f; factor(a[u],f);
  int best=0;
  for(int pr:f)                              // deepest over primes of a[u]
    if(lastp[pr] && (!best || dep[lastp[pr]]>dep[best])) best=lastp[pr];
  for(int id:ask[u]) ans[id] = best ? best : -1;   // BEFORE pushing: i<k
  vector<pair<int,int> > saved;
  for(int pr:f){ saved.push_back(make_pair(pr,lastp[pr])); lastp[pr]=u; }
  for(int v:adj[u]) if(v!=p){ dep[v]=dep[u]+1; dfs(v,u); }
  for(int k=0;k<(int)saved.size();k++)       // undo on exit
    lastp[saved[k].first]=saved[k].second;
}
int main(){
  build_spf();
  scanf("%d%d",&n,&q);
  for(int i=1;i<=n;i++) scanf("%d",&a[i]);
  for(int i=1;i<n;i++){ int u,v; scanf("%d%d",&u,&v);
                        adj[u].push_back(v); adj[v].push_back(u); }
  for(int i=0;i<q;i++){
    scanf("%d%d",&typ[i],&qv[i]);
    if(typ[i]==2) scanf("%d",&qw[i]);
  }
  int id=0;
  while(id<q){
    int start=id;
    for(int u=1;u<=n;u++) ask[u].clear();    // fresh attachment per block
    while(id<q && typ[id]==1){ ask[qv[id]].push_back(id); id++; }
    dep[1]=0; dfs(1,0);                      // whole static block, one DFS
    for(int j=start;j<id;j++) printf("%d\n",ans[j]);
    if(id<q){ a[qv[id]]=qw[id]; id++; }      // the single update between blocks
  }
  return 0;
}
\end{lstlisting}

\solstep{Trace}
\begin{tracebox}[Replaying the official sample to \(-1,1,2,-1,1\)]
Chain \(1\!-\!2\!-\!3\!-\!4\), \(a=[10,8,4,3]\).  Block~1 attaches queries
1--4 to nodes 1--4; \(\mathit{last}[\cdot]\) starts empty. \quad
\textbf{enter~1} (\(10=2{\cdot}5\)): probe \(\mathit{last}[2],\mathit{last}[5]\)
-- both empty \(\Rightarrow\) query~1 \(=\boxed{-1}\); push
\(\mathit{last}[2]=\mathit{last}[5]=1\). \quad
\textbf{enter~2} (\(8=2^3\to\{2\}\)): \(\mathit{last}[2]=1\Rightarrow\)
query~2 \(=\boxed{1}\); push \(\mathit{last}[2]=2\). \quad
\textbf{enter~3} (\(4\to\{2\}\)): \(\mathit{last}[2]=2\Rightarrow\) query~3
\(=\boxed{2}\); push \(3\). \quad
\textbf{enter~4} (\(3\to\{3\}\)): \(\mathit{last}[3]\) empty \(\Rightarrow\)
query~4 \(=\boxed{-1}\).  Every exit restores \(\mathit{last}\) to empty. \quad
Apply \texttt{2 1 9}: \(a_1=9=3^2\).  Block~2 attaches query~5 to node~4.
\textbf{enter~1}: push \(\mathit{last}[3]=1\); nodes 2,3 push
\(\mathit{last}[2]=2,3\); \textbf{enter~4} (\(\{3\}\)):
\(\mathit{last}[3]=1\Rightarrow\) query~5 \(=\boxed{1}\) -- the updated
\(a_1=9\) now shares the prime 3 with \(a_4=3\).  Output
\texttt{-1 1 2 -1 1} -- match. \(\Box\)
\end{tracebox}

\begin{warningbox}[Three ways CF 463E goes wrong]
\begin{enumerate}[nosep,leftmargin=1.6em]
  \item \textbf{Answering after the push.} If \(u\)'s own primes are already
        in the table when its queries are read, \(v\) matches \emph{itself} --
        the statement demands \(i<k\), so every non-\(-1\) answer degenerates
        to \(v\).
  \item \textbf{Factoring with multiplicity.} Pushing prime \(2\) three times
        for \(a_u=8\) makes the undo restore the \emph{middle} copy: after
        leaving \(u\), \(\mathit{last}[2]\) still points at \(u\), and every
        later probe in the DFS is contaminated by a node off the current path.
  \item \textbf{Stale block state.} Attachments (\(\mathit{ask}\)) not cleared
        between blocks fire twice; an update applied in the wrong order (after
        the DFS that should have seen it, or never applied) shifts whole
        answer runs by one version.
\end{enumerate}
\end{warningbox}

\begin{hardeningbox}
\begin{itemize}[nosep,leftmargin=1.4em]
  \item \textbf{The stack window is 27.1 in disguise.} Push-on-entry /
        undo-on-exit is \texttt{add}/\texttt{rm} over a \emph{nested} window.
        You cannot reorder DFS visits the way Mo reorders ranges, so the
        primitive pays exactly when the queried sets are root paths.  The
        Summit (\(\probref{c21-summit}\)) is the same skeleton carrying a full
        frequency table instead of a last-seen table.
  \item \textbf{Price the rebuild honestly.} \((U{+}1)\) full DFSes is the
        trade when updates number \(U\le 50\); it degrades gracefully to
        \(U=\bigO{\sqrt n}\), but \(U=\bigO{n}\) kills it -- that world is Mo
        with modifications or an online structure, not this chapter's door.
  \item \textbf{Table sizing.} \(\mathit{last}[\cdot]\) is indexed by
        \emph{prime value}, not by node: it needs \(2\cdot10^6\) slots, while
        everything else is \(\bigO{n}\).  Sizing it by \(n\) is an
        out-of-bounds write that samples will not catch.
  \item \textbf{Depth bookkeeping.} Comparing "deepest" across primes needs
        \(\depth\); fill it on the same DFS.  A \(10^5\)-deep chain is fine on
        Codeforces' stack, but the explicit-stack version is the
        belt-and-braces port (same caution as Tarjan, \(\secref{c21-tarjan}\)).
\end{itemize}
\end{hardeningbox}

% ---------------------------------------------------------------
\subsection{The Mo-on-subtrees recipe -- and the two doors that beat it}
% ---------------------------------------------------------------

\begin{recognitionbox}[Subtree queries: which tool, exactly]
\begin{itemize}[nosep,leftmargin=1.4em]
  \item \textbf{Distinct colours / freq-shape statistic over subtree\((v)\),
        all queries known upfront.}
        \reckey{run \secref{c21-moengine}'s sweep on ranges
        \([\tin(v),\tout(v)]\) of \(P\); CF 375D is the drill, treated in
        \chref{c28}.}
        \recuse{\item window \(=\) current subtree span; \texttt{add} \(=\)
        insert one node's colour.}
        \recnotuse{\item never re-DFS per query.}
  \item \textbf{Signed/invertible subtree update, point or range read.}
        \reckey{flatten + difference BIT, \(\bigO{\log n}\) per operation -- the
        solved template is Ch.~18's CF 383C (\(\chref{c14}\)).}
        \recuse{\item the parity factor \((-1)^{\depth}\) when signs alternate.}
        \recnotuse{\item Mo here is \(O(\sqrt n)\) tax on a problem with change.}
  \item \textbf{Online updates mixed with subtree aggregates, subtrees only.}
        \reckey{Euler + segment tree; the path version needs the doubled tour
        (\(\secref{c21-path}\)) or HLD (\chref{c16}).}
        \recuse{\item the segment tree inherits the same \([\tin,\tout]\)
        ranges.}
        \recnotuse{\item block order cannot answer query \(i\) after an update
        at query \(i' > i\).}
  \item \textbf{Root-path statistic, strict-ancestor flavour, rare updates.}
        \reckey{DFS-stack window + per-key push/undo table, one rebuild per
        static block (CF 463E, \(\probref{c21-caisa}\)).}
        \recuse{\item read the answer \emph{before} pushing the node's own keys.}
        \recnotuse{\item no binary lifting -- an existential predicate over
        primes is not monotone; no Mo -- a DFS already visits every root path
        in order.}
\end{itemize}
\end{recognitionbox}

\begin{notebox}[The one-line transfer principle]
Every array-Mo solution from 27.1 becomes a subtree-Mo solution by the
substitution \(l\to\tin(v),\ r\to\tout(v)\) -- the proof of
\thmref{c21-subtree-range} is that the sweep never inspects the ranges'
\emph{meaning}.  The next section performs the harder transfer, to
\emph{paths}, where no flatten makes the target contiguous and a
\term{toggle} must be invented.
\end{notebox}

\begin{recallbox}[27.2 in one breath]
\emph{Subtree\((v)\) is the range \([\tin(v),\tout(v)]\) of the preorder array
\(P\) (\(\defref{c21-flatten}\)) -- non-invertible subtree statistics ride the
27.1 block-order sweep unchanged (\(\thmref{c21-subtree-range}\),
\(\coderef{c21-flatten}\)); invertible ones skip Mo (Ch.~18's CF 383C,
\(\chref{c14}\)); and CF 463E opens the third door: the queried set is the
\emph{current root path}, so the recursion stack carries the window -- a
per-prime \(\mathit{last}[p]\) table with \(O(1)\) push/undo, answers read
\emph{before} the push, one re-DFS per block of the \(\le 50\) updates
(\(\probref{c21-caisa}\), \(\coderef{c21-caisa}\), sample replayed to
\(-1,1,2,-1,1\)).}
\selftest{
  \item Which world does "distinct colours in subtree\((v)\), offline" live in, and what answers it? \emph{(No inverse -- Mo's sweep over \([\tin(v),\tout(v)]\) ranges, or DSU-on-tree; \thmref{c21-subtree-range}. \secref{c21-boundary} contrasts them.)}
  \item Why does CF 463E need neither Mo nor binary lifting? \emph{(The DFS stack already visits every root path -- there is no reordering to gain; and "shares a prime with \(v\)" is non-monotone, so lifting has no valid predicate.)}
  \item Name the two ordering facts that make the push/undo table correct. \emph{(Read each node's queries before pushing its own primes -- strict ancestor \(i<k\); factor into distinct primes and undo to the saved values on exit -- one live entry per prime.)}
  \item What promise buys the \(51\) rebuilds, and what would kill the approach? \emph{(\(\le 50\) type-2 queries; at \(U=\bigO{n}\) updates the rebuild is \(\bigO{nq}\) -- that world needs Mo-with-modifications or an online structure.)}
}
\end{recallbox}

% =====================================================================
%  27.3  Mo on PATHS -- the doubled Euler tour E[1..2n] with the
%  membership TOGGLE, and the lca split (ancestor case vs general case,
%  the LCA-parity trap).  WORKED EXAMPLES ONLY -- no solvedproblem here
%  (chapter contract: exactly two, in 27.1 and 27.2).
%  Exports: sec:c21-path  def:c21-toggle  thm:c21-pathsplit
%  lst:c21-euler-toggle
%  Every trace below was replayed by hand against the authoritative
%  running example of 00_opener: E=[1,2,4,8,8,4,5,5,2,3,6,6,7,7,3,1],
%  st(v) node-ordered = (1,2,10,3,7,11,13,4), en(v) = (16,9,15,6,8,12,14,5),
%  val = (2,1,3,1,2,3,1,2), lca from Ch.16 usage (never re-derived).
%  CF 258E / 86D style ideas are PROSE-ONLY; their URLs live elsewhere.
% =====================================================================
\section{Mo on Paths: the Doubled Tour and the Toggle}\label{sec:c21-path}

\begin{previouslybox}[27.2 made subtrees ranges -- paths refuse that]
The flatten of \secref{c21-subtree} works because subtree\((v)\) is a
\emph{contiguous} block of preorder.  A path \(u\)-\(v\) is not: its nodes
\(u, \parent(u), \dots, \lca, \dots, \parent(v), v\) climb out of one subtree
and down another, and no re-ordering of a length-\(n\) array makes every one of
\(\bigO{n^2}\) candidate paths contiguous at once.  The escape is to stop
demanding "the range \emph{equals} the path" and settle for "the range
\emph{represents} the path" -- via the doubled Euler tour \(E\) of the opener,
where each node is written at entry and at exit, and membership is read off by
\textbf{parity of occurrence count}.
\end{previouslybox}

\begin{topicintro}
Length \(2n\) tour, two occurrences per node.  As a Mo pointer crosses a
position it \term{toggles} that node: appears once in the window means
"\emph{in} the represented set", twice means "\emph{out}".  \thmref{c21-pathsplit}
says exactly which window represents path \(u\)-\(v\): order so
\(st(u)<st(v)\); if \(u\) is an ancestor of \(v\) the window \([st(u),st(v)]\)
represents the whole path; otherwise \([en(u),st(v)]\) represents the path
\emph{minus} \(\lca(u,v)\) -- and the LCA must be fed to the statistic by hand,
counted once, then withdrawn.  That last clause is the famous parity trap, and
this section disarms it with three hand-replayed walks on the running tree \(R\):
path\((4,8)\) (ancestor case), path\((8,6)\) (general case), path\((2,3)\) (the
trap \emph{changing an answer}).

\textbf{Prerequisites.} \secref{c21-moengine} (window + \texttt{add}/\texttt{rm}),
\secref{c21-subtree} (the flatten mindset), \(\lca\) as a black box from
\chref{c11} (\(\thmref{c11-distf}\)'s chapter -- Tarjan gets its own offline
implementation in \secref{c21-tarjan}).

\textbf{After this section you can}
\begin{itemize}[nosep]
  \item build the doubled tour and toggle-window it in your sleep; \hfill\lv{4}
  \item state both windows of \thmref{c21-pathsplit} and prove the parity claim; \hfill\lv{5}
  \item handle the ancestor case and the add-the-LCA-by-hand case correctly; \hfill\lv{5}
  \item answer distinct-colour-on-path offline in \(\bigO{(n+q)\sqrt n}\). \hfill\lv{5}
\end{itemize}
\end{topicintro}

% ---------------------------------------------------------------
\subsection{The doubled tour and the toggle window}
% ---------------------------------------------------------------

\begin{definition}{The doubled Euler tour and the toggle}{c21-toggle}
Run the DFS writing \(v\) once when you \emph{enter} it and once when you
\emph{leave} it: \(E[1..2n]\), with every node appearing exactly twice, at
\(st(v)\) (first) and \(en(v)\) (second).  On the running tree \(R\):
\[
  E=[1,2,4,8,8,4,5,5,2,3,6,6,7,7,3,1],
\]
\(st(v)=(1,2,10,3,7,11,13,4)\) and \(en(v)=(16,9,15,6,8,12,14,5)\) for
\(v=1..8\).  A Mo window \([\ell,r]\) over \(E\) \textbf{toggles} node membership
as pointers cross occurrences: a node is \emph{represented} iff it occurs an
odd number of times (i.e.\ exactly once) in \(E[\ell..r]\).  \texttt{add}/\texttt{rm}
are the same test: \(\mathit{held}[x] \oplus= 1\); insert \(x\) into the
statistic when \(\mathit{held}\) turns \(1\), erase when it turns \(0\).
\end{definition}

\begin{theorem}{Two windows represent a path (the LCA split)}{c21-pathsplit}
Let \(st(u)<st(v)\).  (i) If \(\lca(u,v)=u\), the odd-occurrence set of
\(E[st(u)\,..\,st(v)]\) is exactly the node set of path \(u\)-\(v\).
(ii) Otherwise the odd-occurrence set of \(E[en(u)\,..\,st(v)]\) is the path
node set \emph{minus} \(\lca(u,v)\).
\end{theorem}
\begin{proofbox}
A node occurs oddly in a window iff exactly one of \(st,\ en\) lands inside it.
\emph{Case (i):} window \([st(u),st(v)]\), \(u\) an ancestor of \(v\).  Any
path node \(w\) (including \(u\) and \(v\)): \(st(w)\in[st(u),st(v)]\) since
\(w\) lies on the descent from \(u\) to \(v\), while \(en(w)\ge en(v)>st(v)\)
since \(w\) ancestors \(v\) -- exactly one occurrence inside.  A node \(y\)
outside subtree\((u)\): both occurrences left of \(st(u)\) or both right of
\(st(v)\).  A node \(y\in\subtree(u)\) off the path: \(y\) hangs in a side
branch of some path node \(w\), so its span \([st(y),en(y)]\) is disjoint from
\(st(v)\) -- either entirely inside \([st(u),st(v)]\) (two occurrences, even)
or entirely to its right (zero).  Odd set \(=\) path.  \emph{Case (ii):}
\(w=\lca(u,v)\neq u\), window \([en(u),\ st(v)]\).  Ancestor \(x\) of \(u\)
strictly below \(w\): \(st(x)<st(u)<en(u)\) (outside, left) and
\(en(u)<en(x)<st(v)\) (inside: \(x\)'s subtree closes after \(u\)'s does, and
since \(v\notin\subtree(x)\) while \(st(x)<st(v)\), branch \(x\) finishes
before \(v\) opens) -- once.  Symmetrically, ancestor \(y\) of \(v\) strictly
below \(w\): \(st(y)\in(en(u),st(v))\) with \(en(y)>st(v)\) -- once.  \(u\):
only \(en(u)\) in the window; \(v\): only \(st(v)\).  \(w\) itself:
\(st(w)<st(u)\) and \(en(w)>st(v)\) -- \emph{zero} occurrences, the missing
endpoint.  Every other node has both occurrences inside one subtree span
entirely left of \(en(u)\), entirely right of \(st(v)\), or nested and fully
contained in the window -- always an even count.  Odd set
\(=\) path \(\setminus\{w\}\). \(\Box\)
\end{proofbox}

\begin{lstlisting}[caption={Doubled-tour build, toggle add/rm, and the window setup for a path query. The lca node is fed to the statistic between answering steps in case (ii). ASCII only, no underscores}, label={lst:c21-euler-toggle}]
int eu[2*MAXN], ecnt, st[MAXN], en[MAXN], held[MAXN];
void dfs(int u, int p){
    st[u] = ++ecnt;  eu[ecnt] = u;
    for (int v : adj[u]) if (v != p) dfs(v, u);
    en[u] = ++ecnt;  eu[ecnt] = u;          // second write = exit
}
void add(int pos){                            // toggle node at tour position
    int x = eu[pos];
    held[x] ^= 1;
    if (held[x]) insert(x);                   // entered the represented set
    else       erase(x);                      // left it
}
void rm(int pos){ add(pos); }                 // toggle is its own inverse
// window for query (u,v):
int L, R; bool extraLca;
if (st[u] > st[v]) swap(u, v);                // keep st[u] < st[v]
if (is_ancestor(u, v)){ L = st[u]; R = st[v]; extraLca = false; }
else                  { L = en[u]; R = st[v]; extraLca = true;  }
// is_ancestor(a,b): tin/tout from the SAME dfs, tin[a]<=tin[b] && tout[b]<=tout[a]
// (or st[a]<=st[b] && en[b]<=en[a] on this tour).
// In case extraLca:  answer = statistic(window UNION {lca(u,v)})
//                    -- call insert(lca), read cur, erase(lca).
\end{lstlisting}

\begin{insightbox}[Why toggle is legal where insert is not]
On subtrees, each node sits at \emph{one} array position, so "in the window"
was unambiguous.  On the doubled tour a node can cross the pointer twice --
a plain \texttt{add} would count it twice and a distinct-count would lie.
Making \texttt{add} \(\equiv\) \texttt{rm} (an xor-toggle) means two crossings
annihilate: the window's multiset of \emph{distinct represented nodes} always
equals the odd-occurrence set, no matter how the pointers wandered.  This is
exactly why path Mo tolerates the same \(O(1)\) block-order sweep as array Mo
(\(\thmref{c21-complexity}\) applies verbatim with \(n \to 2n\)).
\end{insightbox}

% ---------------------------------------------------------------
\subsection{Worked example: distinct colours on a path of \(R\)}
% ---------------------------------------------------------------

The statistic: \(\mathit{distinct}=\) number of different \(\mathit{val}\)
colours among represented nodes, maintained the classic way -- a colour
frequency array \(\mathit{cf}\), a counter of colours with \(\mathit{cf}>0\);
\texttt{insert(x)}: \(\mathit{cf}[\mathit{val}(x)]{+}{+}\) (bumping \(0\to1\)
raises the count); \texttt{erase(x)} mirrors it.  Colours on \(R\):
\(\mathit{val}=(2,1,3,1,2,3,1,2)\) for nodes \(1..8\).

\begin{tracebox}[Walk 1 -- ancestor case, path\((4,8)\)]
\(st(4)=3 < st(8)=4\); \(\lca(4,8)=4=u\) \(\Rightarrow\) case (i), window
\([3,4]\).  \(E[3..4]=[4,8]\): both occur once \(\Rightarrow\) represented
\(\{4,8\}\) -- the path \(4\to8\) exactly, LCA \(=\) endpoint \(4\) already
inside.  Colours \(\mathit{val}(4)=1,\mathit{val}(8)=2\): \(\boxed{2}\)
distinct.  No extra handling. \(\Box\)
\end{tracebox}

\begin{tracebox}[Walk 2 -- general case, path\((8,6)\)]
\(st(8)=4 < st(6)=11\); \(\lca(8,6)=1\neq8\) \(\Rightarrow\) case (ii), window
\([en(8)=5,\ 11]\).  \(E[5..11]=[8,4,5,5,2,3,6]\).  Occurrence parity:
\(8{:}1\), \(4{:}1\), \(5{:}2\) (toggles out), \(2{:}1\), \(3{:}1\), \(6{:}1\)
\(\Rightarrow\) represented \(\{8,4,2,3,6\}\).  The true path is
\(8\to4\to2\to1\to3\to6\) -- \(\lca=1\) is missing, precisely as
\thmref{c21-pathsplit}(ii) promises.  Feed node \(1\) to \texttt{insert},
read, then \texttt{erase}: represented \(\{8,4,2,1,3,6\}\), colours
\(2,1,1,2,3,3\) \(\Rightarrow\) \(\boxed{3}\) distinct. \(\Box\)
\end{tracebox}

\begin{tracebox}[Walk 3 -- the trap bites: path\((2,3)\) with and without the LCA]
\(st(2)=2 < st(3)=10\); \(\lca(2,3)=1\neq2\) \(\Rightarrow\) window
\([en(2)=9,\ 10]\): \(E[9..10]=[2,3]\), both odd \(\Rightarrow\) represented
\(\{2,3\}\), colours \(1,3\): \(\mathit{distinct}=2\).  But the path is
\(2\to1\to3\) and node \(1\) brings colour \(2\) \(\Rightarrow\) true answer
\(\boxed{3}\).  Forgetting the extra \(\mathit{lca}\) \emph{changes the
answer} -- on walk~2 it happened not to (colour \(2\) was already carried by
node~8), on walk~3 it does.  A single missed trace of this shape is a WA on
CF path-Mo problems. \(\Box\)
\end{tracebox}

\begin{complexitybox}[Path Mo]
Tour build \(\bigO{n}\); each query's window computed in \(O(1)\) given
\(\lca\) (Ch.~16 lifting: \(\bigO{q\log n}\) total) or precomputed batch.
Sweep: \(\bigO{(2n+q)\sqrt{2n}} = \bigO{(n+q)\sqrt n}\) toggles, each with the
\(O(1)\) inner \texttt{insert}/\texttt{erase}.  The add/read/erase of the LCA
is \(O(1)\) per query. \textbf{Space} \(\bigO{n}\) tour + statistic
arrays.  The same skeleton answers "how many nodes with \(\mathit{val}\ge X\)",
"\(\sum \mathit{cf}^2\)"-style pair counts, mode-of-colours on a path --
anything the toggle can maintain at \(O(1)\).
\end{complexitybox}

\begin{warningbox}[The three path-Mo failure modes]
\begin{enumerate}[nosep,leftmargin=1.6em]
  \item \textbf{Forgetting case (ii)'s LCA} (walk~3).  Symptom: answers wrong
        exactly when the LCA's colour is new to the window.  Fix: wrap the
        read with \texttt{insert(lca)}/\texttt{erase(lca)} -- and \emph{only at
        read time}, never as a permanent window member.
  \item \textbf{Swapping endpoints wrongly.} The rules assume \(st(u)<st(v)\);
        compare \(st\), not node id.  Then case (i) is "\(\lca = u\)", the
        \emph{left} endpoint -- if your code tests the wrong endpoint the
        ancestor case silently degrades to a wrong range.
  \item \textbf{Double-adding within case (i).} When \(u\) is the ancestor,
        \(u\) is \emph{already} oddly represented; inserting the LCA again
        toggles it \emph{out}.  The extra insert belongs to case (ii) only.
\end{enumerate}
\end{warningbox}

\begin{hardeningbox}
\begin{itemize}[nosep,leftmargin=1.4em]
  \item \textbf{Query sorting uses the window endpoints} \((L_i/B, R_i)\)
        computed \emph{after} the split -- sorting by raw \((st,r)\) mixes the
        two cases and breaks nothing, but the standard form keeps every window
        a plain range over \(E[1..2n]\).
  \item \textbf{Block size on \(2n\):} \(B \approx 2n/\sqrt q\); with
        \(q=\Theta(n)\), \(B\approx\sqrt{2n}\).
  \item \textbf{The toggle pattern is the tree analogue of "xor is its own
        inverse"}: \texttt{rm = add} because the second occurrence undoes the
        first.  Problems like CF~86D's \(\sum\mathrm{freq}^2\) on arrays do not
        need it; paths cannot live without it.  (86D's URL lives in
        \(\secref{c21-boundary}\)'s prose -- no re-linking.)
\end{itemize}
\end{hardeningbox}

\begin{recallbox}[27.3 in one breath]
\emph{Double the tour, \term{toggle} on every crossing, and a path becomes a
range: \([st(u),st(v)]\) when \(u\) ancestors \(v\), else \([en(u),st(v)]\)
\emph{plus the LCA handed to the statistic at read time}
(\(\thmref{c21-pathsplit}\), \(\coderef{c21-euler-toggle}\)).  Walks on \(R\):
\((4,8)\to2\), \((8,6)\to3\) with the LCA, \((2,3)\to3\) vs.\ a lying \(2\)
without it.  Cost stays \(\bigO{(n+q)\sqrt n}\).}
\selftest{
  \item Why must the represented set be defined by \emph{odd} occurrence? \emph{(Two crossings of the same node must annihilate; parity is the only toggle-invariant that makes rm = add.)}
  \item Give the window for path\((5,6)\) of \(R\) and its odd set. \emph{(\([en(5)=8,\ st(6)=11]=\) E[8..11] \(=[5,2,3,6]\), odd set \(\{5,2,3,6\}\), plus lca \(1\) at read time -- the path \(5,2,1,3,6\).)}
  \item When is the extra LCA insert \emph{not} needed? \emph{(Case (i): lca equals the left endpoint and is already oddly represented.)}
}
\end{recallbox}

% =====================================================================
%  27.4  Tarjan's offline LCA -- one DFS + DSU, WHITE/GREY/BLACK, union a
%  child on return, answer a query when its second endpoint turns BLACK.
%  O((n+q) alpha(n)) with NO lifting table; OFFLINE only.
%  WORKED EXAMPLE ONLY -- no solvedproblem (chapter contract: exactly two,
%  already spent in 27.1/27.2).  Back-cites Ch.16 (\chref{c11}) as the
%  ONLINE owner: lifting is used, never re-derived; here Tarjan is the
%  batch alternative.
%  Exports: sec:c21-tarjan  def:c21-tarjan  thm:c21-tarjan  lst:c21-tarjan
%  Trace runs on the authoritative R of 00_opener with the batch
%  Q1=(8,6)->1, Q2=(4,5)->2, Q3=(2,8)->2, Q4=(6,7)->3, Q5=(5,7)->1;
%  every union/answer event was hand-replayed against the DFS order
%  P=[1,2,4,8,5,3,6,7] before writing.
% =====================================================================
\section{Tarjan's Offline LCA: One DFS and a Disjoint-Set Union}\label{sec:c21-tarjan}

\begin{previouslybox}[Ch.~16 answers one LCA at a time -- this answers the batch]
Chapter~16 (\chref{c11}) built \(\lca\) from binary lifting: \(O(n\log n)\)
preprocessing, then each query climbs in \(O(\log n)\), \emph{online}.  We have
been using that oracle as a black box (\(\secref{c21-path}\) splits windows by
\(\lca\)).  Now suppose the \(q\) LCA queries are all known in advance -- can the
batch be cheaper than \(q\) separate climbs?  Tarjan's 1979 answer says yes, and
it is the most beautiful five-line idea in this chapter: run \emph{one plain
DFS}, keep a disjoint-set union of the region you have \emph{finished}, and
whenever a node dies (turns BLACK) it settles every old query whose other
endpoint already died.  Total \(\bigO{(n+q)\alpha(n)}\) -- effectively linear,
no lifting table, no \(\log\).
\end{previouslybox}

\begin{topicintro}
Paint nodes \emph{WHITE} (unseen), \emph{GREY} (on the DFS stack -- entered,
children pending), \emph{BLACK} (finished).  When a child \(v\) returns to its
parent \(u\), \(\mathit{union}\) \(v\)'s set into \(u\)'s and stamp
\(\mathit{ancestor}[\mathit{find}(u)] = u\): each DSU set is exactly one
maximal BLACK region, and its stamp is the GREY parent that region is waiting
under.  Attach every query \((u_i,v_i)\) to \emph{both} endpoints' lists.  When
\(u\) turns BLACK, scan its list: a partner already BLACK means \emph{this} is
the moment the two endpoints' paths met -- the answer is
\(\mathit{ancestor}[\mathit{find}(\text{partner})]\).  \thmref{c21-tarjan} proves
it and bounds it.  The worked trace runs five queries on the running tree \(R\),
including both sibling cases and a strict ancestor case.

\textbf{Prerequisites.} DSU with path compression is used as a black box here
(\(\mathit{find}/\mathit{union}\), amortised \(\alpha\)); \chref{c11} owns what
an \(\lca\) \emph{is} and the online lifting alternative this section is
measured against; the WHITE/GREY/BLACK DFS vocabulary is the standard
cycle-detection language from the traversal chapters.

\textbf{After this section you can}
\begin{itemize}[nosep]
  \item run the three-colour DFS with union-on-return and state the invariant; \hfill\lv{4}
  \item answer a batch of \(q\) LCAs in \(\bigO{(n+q)\alpha}\) offline; \hfill\lv{4}
  \item trace \(\mathit{ancestor}[\mathit{find}(\cdot)]\) answers by hand (done: five on \(R\)); \hfill\lv{5}
  \item choose between Tarjan, lifting, and Euler+RMQ from the statement's modality. \hfill\lv{4}
\end{itemize}
\end{topicintro}

% ---------------------------------------------------------------
\subsection{The algorithm and its invariant}
% ---------------------------------------------------------------

\begin{definition}{Tarjan's offline LCA}{c21-tarjan}
Root the tree; attach each query \((u_i,v_i)\) to both \(u_i\) and \(v_i\).
DFS: on \emph{entering} \(u\) set \(\mathit{col}[u]=\)GREY, \(\mathit{fa}[u]=u\),
\(\mathit{anc}[u]=u\); after each child \(v\) returns,
\(\mathit{fa}[\mathit{find}(v)] := \mathit{find}(u)\) and
\(\mathit{anc}[\mathit{find}(u)] := u\); when all children are done, for each
attached query whose partner is BLACK record
\(\mathit{ans}_i = \mathit{anc}[\mathit{find}(\text{partner})]\), then set
\(\mathit{col}[u]:=\)BLACK.
\end{definition}

\begin{theorem}{Tarjan answers every offline LCA in \(\bigO{(n+q)\alpha(n)}\)}{c21-tarjan}
At every instant, each DSU set is precisely a maximal fully-BLACK connected
region, and \(\mathit{anc}[\mathit{find}(x)]\) is the deepest GREY node whose
subtree contains \(x\).  When a query's second endpoint \(u\) turns BLACK,
that stamp equals \(\lca(u,v)\).  Total time \(\bigO{(n+q)\alpha(n)}\).
\end{theorem}
\begin{proofbox}
\emph{Invariant.} Induction over DFS events.  A child region \(C\) merges into
\(u\)'s side exactly when \(C\) is maximal-BLACK (child returned) and \(u\) is
still GREY, so sets are maximal BLACK regions glued under their waiting GREY
parent; stamping \(\mathit{anc}[\mathit{find}(u)]=u\) after each merge keeps the
stamp that parent. \emph{Correctness.} Suppose \(u\) turns BLACK and partner
\(v\) is already BLACK (wlog \(v\) died first, or \(v\) is in an earlier branch).
At \(u\)'s death moment, \(\mathit{find}(v)\)'s region is the maximal subtree
rooted at some GREY node's completed child-chain, and its stamp \(w\) is the
deepest GREY ancestor of \(v\).  GREY nodes form exactly the current root-to-\(u\)
stack, so \(w\) is on that stack.  \(w\) ancestors \(v\) and \(w\) ancestors
\(u\) (it is the DFS-stack node currently processing the region that contains
\(u\)); any deeper stack node \(z\) child-of-\(w\) toward \(u\) cannot contain
\(v\) (if it did, \(v\)'s BLACK region would have merged up through \(z\) into
\(w\) only after \(z\)'s return -- and \(z\) is GREY now, with \(u\) inside it,
while \(v\) died inside an \emph{already returned} child of \(w\)).  Hence
\(w=\lca(u,v)\). \emph{Cost.} \(n-1\) unions, \(O(n+q)\) finds, DSU with path
compression (union here always points a child set at its parent set, so even
without rank the amortised bound is \(\bigO{\alpha}\) for the mixed batch) --
total \(\bigO{(n+q)\alpha(n)}\). \(\Box\)
\end{proofbox}

\begin{lstlisting}[caption={Tarjan offline LCA in full: two-endpoint query attachment, union-on-child-return, answer on second endpoint's death. ASCII only}, label={lst:c21-tarjan}]
#include <bits/stdc++.h>
using namespace std;
const int MAXN=200005;
vector<int> adj[MAXN];
vector<pair<int,int>> qlist[MAXN];      // (partner, query id), on BOTH ends
int fa[MAXN], anc[MAXN], col[MAXN];     // col: 0 WHITE, 1 GREY, 2 BLACK
int ans[MAXN];
int find(int x){ return fa[x]==x ? x : fa[x]=find(fa[x]); }
void dfs(int u, int p){
  col[u] = 1;  fa[u] = u;  anc[u] = u;
  for (int v : adj[u]) if (v != p){
    dfs(v, u);
    fa[find(v)] = find(u);              // finished region joins the parent
    anc[find(u)] = u;                   // ... whose stamp is u
  }
  for (auto &Q : qlist[u])
    if (col[Q.first] == 2)              // partner already dead -> answer now
      ans[Q.second] = anc[find(Q.first)];
  col[u] = 2;
}
int main(){
  // read n, q, edges; attach: qlist[u].push_back({v,i}); qlist[v].push_back({u,i});
  // dfs(1, 0);  print ans[0..q-1]
}
\end{lstlisting}

% ---------------------------------------------------------------
\subsection{Worked example: five queries on \(R\), every event traced}
% ---------------------------------------------------------------

Run \(\coderef{c21-tarjan}\) on the running tree \(R\) (\(\figref{c21-tree}\)),
children visited left-to-right, DFS preorder \(P=[1,2,4,8,5,3,6,7]\), with the
batch:
\[
  Q_1=(8,6),\quad Q_2=(4,5),\quad Q_3=(2,8),\quad Q_4=(6,7),\quad Q_5=(5,7).
\]
Expected answers (Ch.~16 by inspection of the picture):
\(\lca(8,6){=}1\), \(\lca(4,5){=}2\), \(\lca(2,8){=}2\), \(\lca(6,7){=}3\),
\(\lca(5,7){=}1\).

\begin{tracebox}[The full event log -- unions, stamps, and the five answer moments]
\textbf{Enter} 1,2,4,8 (GREY stack \([1,2,4,8]\)). \
8: no children; its lists \(Q_1,Q_3\): partners 6 WHITE, 2 GREY -- \emph{no
answer}; \textbf{BLACK}(8). \
Return 8\(\to\)4: \(\mathit{union}(4,8)\), set \(\{4,8\}\), stamp \(4\). \
4: lists \(Q_2,Q_3\): partner 5 WHITE, 2 GREY -- no answer;
\textbf{BLACK}(4). \
Return 4\(\to\)2: \(\mathit{union}(2,4)\), set \(\{2,4,8\}\), stamp \(2\). \
Enter 5; 5 dies: \(Q_2\): partner 4 \textbf{BLACK} \(\Rightarrow\)
\(\mathit{ans}_2=\mathit{anc}[\mathit{find}(4)]=\boxed{2}\) ✓;
\(Q_5\): partner 7 WHITE -- wait; \textbf{BLACK}(5). \
Return 5\(\to\)2: \(\mathit{union}(2,5)\), set \(\{2,4,5,8\}\), stamp \(2\). \
2 dies: \(Q_3\): partner 8 BLACK \(\Rightarrow\)
\(\mathit{ans}_3=\mathit{anc}[\mathit{find}(8)]=\boxed{2}\) ✓ (note 8's set is
\(\{2,4,5,8\}\) stamped \(2\) -- the ancestor case answered one union after
8's own death). \
Return 2\(\to\)1: \(\mathit{union}(1,2)\), set \(\{1,2,4,5,8\}\), stamp \(1\). \
Enter 3,6; 6 dies: \(Q_1\): partner 8 BLACK \(\Rightarrow\)
\(\mathit{ans}_1=\mathit{anc}[\mathit{find}(8)]=\boxed{1}\) ✓ -- the two
endpoints' regions met under 1; \(Q_4\): 7 WHITE -- wait. \
Return 6\(\to\)3: \(\mathit{union}(3,6)\), set \(\{3,6\}\), stamp \(3\). \
Enter 7; 7 dies: \(Q_4\): partner 6 BLACK \(\Rightarrow\)
\(\mathit{ans}_4=\mathit{anc}[\mathit{find}(6)]=\boxed{3}\) ✓ (siblings meet
under 3); \(Q_5\): partner 5 BLACK \(\Rightarrow\)
\(\mathit{ans}_5=\mathit{anc}[\mathit{find}(5)]=\boxed{1}\) ✓ (5's set merged
into 1 at 2's death). \
Return 7\(\to\)3: \(\mathit{union}(3,7)\); 3 dies (empty lists); return,
\(\mathit{union}(1,3)\); 1 dies.  Answers: \(1,2,2,3,1\) -- all five match.
\(\Box\)
\end{tracebox}

\begin{insightbox}[Read the invariant off the trace in one line]
At every answer moment, \(\mathit{find}(\text{partner})\)'s set was "everything
already finished in one branch", and the stamp was "the still-open node
directly above it" -- which, because GREY nodes are exactly the current
root-to-\(u\) stack, is precisely where the two endpoints' paths merge.  Tarjan
is not \emph{computing} ancestors; it is watching the DFS stack and naming the
merge point at the instant both sides stop moving.
\end{insightbox}

\begin{comparisonbox}[Tarjan vs.\ lifting (Ch.~16) vs.\ Euler+RMQ (Ch.~5)]
\begin{itemize}[nosep,leftmargin=1.4em]
  \item \textbf{Modality.} Tarjan: pure offline -- a query asked after an
        update, or answered out of batch order, is outside its contract.
        Lifting (\(\chref{c11}\)): online, \(O(\log n)\)/query, table
        \(\bigO{n\log n}\).  Euler tour + RMQ: offline-or-online after one
        \(\bigO{n}\) build plus \(\bigO{n\log n}\) (or linear Cartesian) RMQ.
  \item \textbf{Cost at \(q \approx n\) large:} Tarjan \(\bigO{(n+q)\alpha}\)
        \(<\) both \(\log\) methods in theory, and \emph{much} less memory --
        no table.  This is why contest mega-batches (SPOJ LCA, problem data with
        \(n,q\ge10^5\) and tight memory) hand themselves to it.
  \item \textbf{Fragility.} Tarjan needs all queries upfront \emph{and} a
        recursion-depth-safe DFS at \(n=2\cdot10^5\); lifting needs neither but
        pays the table.  Path Mo (\(\secref{c21-path}\)) needs \(\lca\) only to
        \emph{sort windows}, so batch its LCAs through Tarjan and the whole
        pipeline stays \(\bigO{(n+q)\sqrt{n+q}}\) with one \(\bigO{n}\) prep.
\end{itemize}
\end{comparisonbox}

\begin{complexitybox}[Tarjan]
One DFS: \(n-1\) unions; each query touched twice on attach, one find at the
answering endpoint (the WHITE/GREY partner check on the other endpoint costs
\(O(1)\)).  With path compression the DSU batch is \(\bigO{(n+q)\alpha(n)}\)
\(\approx\) linear.  \textbf{Space} \(\bigO{n+q}\) -- \emph{no lifting table},
which for \(n=10^6\) is the difference between fitting and not fitting.  The
recursion is the real risk: convert to an explicit stack when \(n\) is large.
\end{complexitybox}

\begin{warningbox}[Three ways Tarjan goes wrong]
\begin{enumerate}[nosep,leftmargin=1.6em]
  \item \textbf{Attaching the query to only one endpoint.} The answer fires at
        whichever endpoint turns BLACK \emph{second}; if the query sits only on
        the first one, it checks a WHITE/GREY partner, skips, and the query is
        silently never answered (stale \(0\) in \(\mathit{ans}\)).
  \item \textbf{Stamping after \(\mathit{find}(v)\) instead of \(\mathit{find}(u)\)}
        (or forgetting the stamp entirely): the set survives but its
        \(\mathit{anc}\) entry is wrong, and answers drift up or down by one
        level -- test on the \(Q_3=(2,8)\) ancestor shape, which dies first to
        this bug.
  \item \textbf{Checking \(\mathit{col}[\text{partner}] != \mathit{WHITE}\) instead
        of \(=\)BLACK:} a GREY partner means "query is being asked inside the
        partner's unfinished subtree"; answering from its stamp at that moment
        can look right (\(\mathit{anc}[\mathit{find}(\text{partner})]\) is the
        partner or its child) only because of an accident of that variant's
        invariants -- the clean contract is: \emph{answer exactly when the
        partner is BLACK}.  (The one-flag implementations you see on the web
        fold "GREY ancestor" into BLACK's formula -- derive it yourself before
        trusting it.)
\end{enumerate}
\end{warningbox}

\begin{hardeningbox}
\begin{itemize}[nosep,leftmargin=1.4em]
  \item \textbf{Same skeleton, other offline reads.} The union-on-return +
        death-time-scan pattern computes more than LCA: batch "nearest marked
        ancestor", offline subtree counts at death time -- the scan list is a
        general "what is visible from the stack right now" hook.
  \item \textbf{Iterative form.} Push \((u,\text{iterator state})\) frames; the
        union and the scan both happen on \emph{pop} -- one loop, no stack
        overflow at \(n=2\cdot10^5\).
  \item \textbf{Forest input?} Loop \(\mathit{dfs}\) over WHITE roots; queries
        across components must be answered separately (no LCA) -- check whether
        the statement guarantees connectivity before trusting \(\mathit{ans}\).
\end{itemize}
\end{hardeningbox}

\begin{recallbox}[27.4 in one breath]
\emph{One DFS, three colours; union a returned child into the parent and stamp
the parent; a query fires when its second endpoint turns BLACK, reading
\(\mathit{anc}[\mathit{find}(\text{partner})]\) -- the merge point of the two
endpoints' paths (\(\defref{c21-tarjan}\), \(\thmref{c21-tarjan}\),
\(\coderef{c21-tarjan}\)).  The five-query log on \(R\) returns
\(1,2,2,3,1\).  Batch modality \(\Rightarrow\) Tarjan; online \(\Rightarrow\)
Ch.~16 lifting.}
\selftest{
  \item What is \(\mathit{anc}[\mathit{find}(x)]\) mid-DFS? \emph{(The deepest GREY node whose subtree contains the finished region \(x\)'s set belongs to -- the invariant of \thmref{c21-tarjan}.)}
  \item Why attach queries to both endpoints? \emph{(The firing moment is the second endpoint's death; either endpoint can be second.)}
  \item State the cost difference against Ch.~16 at \(n=q=10^6\), memory-wise. \emph{(Tarjan O((n+q)alpha), space O(n+q), no table; lifting O((n+q)log n) with an O(n log n) table -- often the memory ceiling.)}
}
\end{recallbox}

% =====================================================================
%  27.5  Recognition and the method boundary -- the Mo signature
%  (offline + overlapping + non-invertible), and the six hand-offs:
%  invertible -> Euler+Fenwick (Ch.5); online path -> HLD (Ch.22) / LCT
%  (Ch.25); per-set fold -> virtual tree (Ch.26); small marked sets /
%  subtree folds -> DSU-on-tree (Ch.17); all-pairs distance shapes ->
%  centroid (Ch.24); LCA-only batches -> Tarjan (this chapter, 27.4).
%  WORKED pick examples only -- no solvedproblem (contract: two, spent).
%  Exports: sec:c21-boundary  def:c21-signature  def:c21-pick  lst:c21-kit
%  CF 86D (Powerful array) and CF 375D are PROSE references only -- their
%  URLs live in other chapters; nothing here links a foreign URL.
% =====================================================================
\section{The Method Boundary: When Mo Is Wrong}\label{sec:c21-boundary}

\begin{previouslybox}[Ch.~26 ended at a boundary -- it was this one]
The last chapter of this book's previous stop, virtual trees (\chref{c30}),
closed its pick table with an honest confession: "many overlapping offline
node-set queries whose statistic will not fold per-set" does not compress onto
any virtual tree, and belongs to Mo.  You now own Mo -- so the debt must be
collected in the other direction: \emph{which shapes do NOT belong to it}?  A
tool this comfortable (\(\bigO{(n+q)\sqrt n}\) answers from one sweep) is also
this easy to misfire: onto online queries (block order dies), onto invertible
statistics (a BIT is \(\sqrt n/\log n\) times cheaper), onto per-set folds a
virtual tree eats in \(\bigO{k\log k}\).  This section is the boundary contract
for the whole chapter.
\end{previouslybox}

\begin{topicintro}
Three questions, asked in order, decide everything.  \emph{Are all queries
known before any answer is needed?}  No \(\Rightarrow\) Mo is illegal -- online
path shapes go to HLD (\chref{c16}) or the link-cut tree (\chref{c23}).
\emph{Does the statistic telescope over a prefix?}  Yes \(\Rightarrow\) Mo is
overkill -- Euler flatten + Fenwick/segment answers in
\(\bigO{(n+q)\log n}\) (Ch.~18's CF 383C chose this door, \(\chref{c14}\); the
root-path batch CF 463E in \(\secref{c21-subtree}\) chose the stack instead).
\emph{Are the queried sets subtrees/paths of ONE fixed tree, many of them
overlapping heavily?}  Yes \(\Rightarrow\) this chapter's engine.  The three
negative escapes -- small marked sets (DSU-on-tree, \chref{c28}), per-set
compression (virtual tree, \chref{c30}), distance distributions over all pairs
(centroid, \chref{c18}) -- are worked through below, then \defref{c21-pick}
codifies the order of the questions.

\textbf{Prerequisites.} The whole chapter, plus the boundary sections of
Ch.~17/22/24/26 it hands work back to.

\textbf{After this section you can}
\begin{itemize}[nosep]
  \item recite the Mo signature (offline, overlapping, non-invertible, \(O(1)\) update); \hfill\lv{4}
  \item reject Mo correctly for online, invertible, and per-set-fold shapes; \hfill\lv{5}
  \item choose Mo vs.\ DSU-on-tree for subtree statistics with a reason; \hfill\lv{5}
  \item price each alternative before committing \(\bigO{(n+q)\sqrt n}\). \hfill\lv{4}
\end{itemize}
\end{topicintro}

% ---------------------------------------------------------------
\subsection{The signature}
% ---------------------------------------------------------------

\begin{definition}{The Mo signature (four tests)}{c21-signature}
A problem is Mo-shaped iff all four hold:
\begin{enumerate}[nosep,leftmargin=1.6em]
  \item \textbf{Offline:} every query is known upfront; no query's answer
        feeds an update (reordering the questions must be legal).
  \item \textbf{Overlapping ranges:} the query sets are ranges/paths of
        \emph{one fixed tree} that share elements heavily -- otherwise there
        is no work to amortise.
  \item \textbf{Non-invertible statistic:} distinct counts,
        \(\sum\mathrm{freq}^2\), mode, "how many values \(\ge l\) times" --
        no prefix/subtraction form exists.
  \item \textbf{\(O(1)\) (de)activation:} adding or removing \emph{one node}
        updates the statistic in \(O(1)\) (or \(O(\log)\) with a small
        structure -- then re-price).
\end{enumerate}
Fail any test and the table below names where the problem really lives.
\end{definition}

\begin{recognitionbox}[Question shapes: Mo versus each look-alike]
\begin{itemize}[nosep,leftmargin=1.4em]
  \item \textbf{"distinct colours on subtree\((v)\)" for many \(v\), all queries
        given at once.}
        \reckey{tests 1--4 all pass: flatten + path-free Mo.}
        \recuse{\item the 27.2 transfer: window over \([\tin(v),\tout(v)]\).}
        \recnotuse{\item DSU-on-tree also solves it (\chref{c28}) -- see the pick
        example below; price both, then move.}
  \item \textbf{"sum of values on subtree\((v)\)" with point updates
        interleaved.}
        \reckey{invertible \emph{and} online: Euler + Fenwick/segment,
        \(\bigO{\log n}\) per op (\chref{c24}).}
        \recuse{\item the difference/prefix form.}
        \recnotuse{\item Mo cannot see query \(i{+}1\) after an update at \(i\).}
  \item \textbf{"path \(u\)-\(v\) aggregate, answers needed online, edge updates
        allowed".}
        \reckey{the HLD contract (\chref{c16}); with topology-like operations,
        the LCT (\chref{c23}).}
        \recuse{\item decompose path \(\to\) \(O(\log n)\) segment ranges.}
        \recnotuse{\item Mo: test 1 fails immediately.}
  \item \textbf{Same path statistic but offline.}
        \reckey{path Mo (\(\secref{c21-path}\)): doubled tour + toggle.}
        \recuse{\item block order on tour windows.}
        \recnotuse{\item HLD when \(q\le 2\cdot10^5\) and updates exist -- Mo has
        lost test 1 again.}
  \item \textbf{"aggregate over each query's small marked set" (\(k\) marked,
        \(\sum k\) moderate).}
        \reckey{virtual tree (\chref{c30}): build \(\bigO{k\log k}\), fold.}
        \recuse{\item the sets are \emph{different} per query and do not overlap
        as ranges.}
        \recnotuse{\item Mo -- windows would shuffle arbitrary non-controlling
        node sets, no amortisation.}
  \item \textbf{"count pairs at distance \(d\)", "sum over all paths" (no query
        list at all).}
        \reckey{centroid decomposition (\chref{c18}) -- the object is the whole
        tree's path multiset.}
        \recuse{\item divide-and-conquer on paths.}
        \recnotuse{\item Mo -- there are no ranges to order.}
\end{itemize}
\end{recognitionbox}

% ---------------------------------------------------------------
\subsection{The pick, worked: Mo versus DSU-on-tree}
% ---------------------------------------------------------------

The sharpest boundary inside this chapter's own regime (offline, subtree,
non-invertible) is Mo vs.\ DSU-on-tree (\chref{c28}, small-to-large): both
answer "distinct colours in subtree\((v)\) for all \(v\)" -- CF 375D is the
textbook instance (its full treatment and URL live there; named in prose only
here, by contract).

\begin{center}\renewcommand{\arraystretch}{1.25}\setlength{\tabcolsep}{3pt}
\begin{tabular}{@{}p{0.20\linewidth} p{0.24\linewidth} p{0.24\linewidth} >{\raggedright\arraybackslash}p{0.18\linewidth}@{}}
\toprule
\textbf{Axis} & \textbf{Mo (this chapter)} & \textbf{DSU-on-tree (Ch.17)} & \textbf{Wins}\\
\midrule
Shape & any ranges/paths, \(q\) arbitrary & subtree queries, all \(v\) (or many) & Mo (generality) \\
Time & \(\bigO{(n+q)\sqrt n}\) & \(\bigO{n\log n}\) (subtrees only) & DSU (on subtrees) \\
Paths & native (27.3 toggle) & awkward (keeps subtrees, not paths) & Mo \\
Answers & stored offline, printed at end & available for every node at end & tie \\
Stat update & must be \(O(1)\) per move & must merge maps small-to-large & Mo (simpler contract) \\
Memory & \(\bigO{n+q}\) flat & \(\bigO{n}\) + per-node maps & Mo \\
\bottomrule
\end{tabular}
\end{center}

\begin{tracebox}[Worked pick: three statements, three verdicts]
\textbf{(a)} \(n,q\le10^5\); query: distinct \(\mathit{val}\) on path \(u\)-\(v\);
all queries given upfront.  Tests: 1 pass, 2 pass, 3 pass (distinct), 4 pass
(toggle) \(\Rightarrow\) \emph{path Mo}, \(\bigO{(n+q)\sqrt n}\) \(\approx\)
\(6\cdot10^7\) moves -- fits.  DSU-on-tree cannot hold a path.
\textbf{(b)} \(n\le10^5\); same statistic for \emph{every subtree} (query set
\(=\) all \(v\)).  Test 3 passes, but the subtree-all-queries specialisation is
\(\bigO{n\log n}\) with Ch.~17's keep-heavy-child trick: DSU-on-tree by
\(\sqrt n/\log\) factor.  Verdict \emph{DSU-on-tree} -- Mo still \emph{correct},
just slower; know the price you are paying.
\textbf{(c)} \(n,q\le2\cdot10^5\); subtree distinct, but queries arrive one per
line and an update \(a_x := y\) may follow each answer.  Test 1 fails
\(\Rightarrow\) neither Mo nor plain DSU: the honest door is sqrt-decomposition
\emph{over time} (rebuild DSU answers per block of operations) -- beyond this
chapter; recorded so the rejection names a destination, not a cliff. \(\Box\)
\end{tracebox}

\begin{definition}{The order of the four questions (pick protocol)}{c21-pick}
\emph{Online?} \(\Rightarrow\) HLD/LCT (\chref{c16}, \chref{c23}) or Euler+segment
for subtrees.  \emph{Invertible?} \(\Rightarrow\) flatten + Fenwick
(\chref{c24}); the parity-normalised CF 383C (Ch.~18, \(\chref{c14}\)) is the
template.  \emph{All query sets
small marked subgraphs?} \(\Rightarrow\) virtual tree (\chref{c30}).
\emph{Whole-tree path distribution?} \(\Rightarrow\) centroid (\chref{c18}).
\emph{Subtree-only, every node?} \(\Rightarrow\) price DSU-on-tree (\chref{c28})
first.  \emph{Still here?} The problem is Mo's by elimination -- and it is
\(\bigO{(n+q)\sqrt{n+q}}\) (array/subtree) or path Mo with the doubled tour.
\end{definition}

\begin{lstlisting}[caption={The pick protocol as a decision skeleton -- paste-ready comments for contest code; each branch names the chapter that owns the technique}, label={lst:c21-kit}]
// 1. do answers have to come BETWEEN updates?  -> NO Mo.
//      path ranges: HLD (Ch.22). dynamic topology: LCT (Ch.25).
// 2. can the statistic be written pref(r) - pref(l-1)? -> NO Mo.
//      subtree sums / parity-normalised adds: Euler + Fenwick (Ch.5).
// 3. is each query a SMALL marked set to fold over?     -> virtual tree (Ch.26)
// 4. no query list at all -- whole-tree pair counts?     -> centroid (Ch.24)
// 5. subtree-only, statistic mergeable small-to-large?   -> DSU on tree (Ch.17)
// 6. LCA batch only, offline, want linear-ish?           -> Tarjan (27.4)
// 7. otherwise: offline, overlapping ranges/paths, O(1) add/rm:
//        subtrees -> flatten P, Mo on [tin,tout]         (27.2)
//        paths    -> doubled tour E, toggle, lca split   (27.3)
//        B = (int)( (2*n) / sqrt(max(1,q)) ) ; odd-even R order.
\end{lstlisting}

\begin{insightbox}[Mo is the tool of last resort -- and that is a compliment]
Every rejection above is a statement that something \emph{cheaper} applies:
invertibility buys \(\log\), structure buys \(\log\), smallness buys \(k\log k\).
Mo is what remains when no structure exists -- but it always exists, needing
only test 4.  "I cannot compress it, so I will order it" is the mature
contest reflex: \(\bigO{(n+q)\sqrt n}\) with a 12-line engine beats a
\(\bigO{n\log^2 n}\) structure you cannot debug at minute 80.
\end{insightbox}

\begin{notebox}[Boundary homework this chapter leaves elsewhere]
Three shapes keep walking out of these doors and get settled by later
chapters: online path queries with updates arrive at HLD (\chref{c16}) and its
dynamic cousin (\chref{c23}); "tree isomorphism / canonical form" questions --
which subtree equals which -- belong to Chapter~28, next door, where hashing
supplies the equality test Mo never needed (Mo compares counts, not trees).
And if your "non-invertible statistic" is actually "frequency of frequencies"
\emph{on root paths}, the Summit at the end of this chapter answers it directly.
\end{notebox}

\begin{complexitybox}[The boundary, priced]
Same input class (offline subtree/path, non-invertible): Mo
\(\bigO{(n+q)\sqrt n}\) general \(\to\) DSU-on-tree \(\bigO{n\log n}\) (subtrees
only).  Invertible: \(\bigO{(n+q)\log n}\) Euler+Fenwick.  Online paths: HLD
\(\bigO{(n+q\log n)\log n}\)-ish with updates.  Per-set: virtual tree
\(\bigO{k\log k}\).  LCA batch: Tarjan \(\bigO{(n+q)\alpha}\).  The gaps are a
factor \(\sqrt n\) to \(\log^2 n\) -- enough to flip AC/TLE at \(n=2\cdot10^5\),
which is why the protocol order in \defref{c21-pick} is not negotiable.
\end{complexitybox}

\begin{recallbox}[27.5 in one breath]
\emph{Four tests make a Mo problem: offline, overlapping, non-invertible,
\(O(1)\) update (\(\defref{c21-signature}\)).  Each failed test is not a dead
end but a hand-off: online \(\to\) HLD/LCT; invertible \(\to\) Fenwick; small
marked \(\to\) virtual tree; all-pairs \(\to\) centroid; subtree-all \(\to\)
DSU-on-tree; LCA-only \(\to\) Tarjan.  Ask the questions in this order
(\(\defref{c21-pick}\), \(\coderef{c21-kit}\) -- and only then fire the
sweep.}
\selftest{
  \item Which single test does an interleaved-update problem fail, and what should you reach for on paths? \emph{(Test 1, offline; HLD (Ch.22), or LCT (Ch.25) for dynamic trees.)}
  \item Subtree distinct-colours for \emph{every} node: Mo or DSU-on-tree, and why? \emph{(DSU-on-tree -- the specialised O(n log n) beats O((n+q) sqrt n) when queries are exactly all subtrees.)}
  \item Name the CF 383C lesson in one sentence. \emph{(Scary updates can still be invertible after normalisation -- always look for the prefix before paying the sqrt tax; Ch.~18 solved it that way.)}
  \item CF 463E passed both the BIT and Mo doors yet was solved with neither. What did it walk through instead? \emph{(The third door: the recursion stack carries the root-path window; per-key push/undo table plus one rebuild per block of the \(\le 50\) updates.)}
}
\end{recallbox}

% =====================================================================
%  Chapter 27 review -- the 14-station revision contract.
%  Revision-only: tables, one-liners, pointers back via \secref/\thmref/
%  \defref/\figref/\coderef/\probref.  No re-derivations.
%  Summit: CF 1585E "Frequency Queries" (2400, L6) -- a rooted tree with an
%  integer a_i on each node and q offline questions (v,l,k): read the values on
%  the path v -> root, tally each value's frequency, keep values with freq >= l,
%  order the survivors by frequency INCREASING (ties arbitrary), and report the
%  k-th value (or -1).  This is the chapter's own "non-invertible offline
%  root-path statistic": a plain DFS from the root is the window (push a value on
%  entry, pop it on exit -- the exact O(1) add/remove primitive of 27.1, now
%  riding the recursion stack instead of Mo's two pointers), cnt[val] + a
%  frequency-of-frequency Fenwick select the k-th, and a per-frequency value
%  stack returns the answer value.  O((n+q) log n).  Reuses TWO chapter tools
%  (the Mo frequency table + the offline batch answering of Tarjan) and ONE
%  earlier-chapter tool (the Ch.5 depth/parent DFS).  Both official samples
%  hand-replayed BEFORE writing: [1 -1 1] and [1 1 2 1 -1].
%  Exports prob:c21-summit, lst:c21-summit, lst:c21-kit (kit lives in 27.5).
%  Mixed set EXACTLY 10, all live-verified this session (names/ratings/tags
%  re-pulled from the CF problemset API), un-owned, ASCENDING by rating; each
%  hint routes via the signature (def:c21-signature) or REJECTS Mo and names the
%  true tool (Fenwick/Euler, HLD, DSU-on-tree, centroid, LCA).  The canonical
%  tree-Mo problems 375D (Ch.17), 258E, 86D are owned elsewhere and NEVER
%  re-linked here.
% =====================================================================

\chaptersection{Chapter 27 Summary}

One sentence: \emph{when one fixed tree faces a huge offline batch of subtree or
path queries whose statistic has \textbf{no inverse} -- distinct colours,
\(\sum\mathrm{freq}^2\), "how many values appear \(\ge l\) times" -- you stop
being online and instead \emph{reorder} the queries so consecutive ones differ by
a single node, maintaining the statistic with an \(O(1)\) add/remove across a
\(\bigO{(n+q)\sqrt n}\) block-order sweep; and when the queries only want
\(\lca\), Tarjan's one-DFS \(+\) DSU answers the whole batch in
\(\bigO{(n+q)\alpha}\) (\(\secref{c21-moengine}\), \(\thmref{c21-complexity}\),
\(\thmref{c21-tarjan}\)).} Chapter~26 handed you the "many overlapping offline
node-set queries, non-invertible statistic" case at its method boundary; this
chapter is the answer to it. Each row below is a pointer, not a proof.

\begin{center}
\renewcommand{\arraystretch}{1.25}
\setlength{\tabcolsep}{3pt}
\begin{tabular}{@{}p{0.22\linewidth} p{0.53\linewidth} p{0.19\linewidth}@{}}
\toprule
\textbf{Result} & \textbf{Statement} & \textbf{Where} \\
\midrule
Block order & Sort \(q\) offline ranges by \((\lfloor L/B\rfloor,R)\), creep one
window \([curL,curR]\) with an \(O(1)\) \texttt{add}/\texttt{rm}; with
\(B=\Theta(n/\sqrt q)\) the total pointer movement is \(\bigO{(n+q)\sqrt n}\). &
\defref{c21-mo}, \thmref{c21-complexity} \\
XOR engine & Prefix-xor turns "subarray xor \(=k\)" into "count pairs with
\(px[p]\oplus px[q]=k\) in \([l-1,r]\)"; insert \(v\) forms \(cnt[v\oplus k]\) new
pairs -- Mo on \(px[0..n]\). & \probref{c21-xor}, \coderef{c21-xor} \\
Subtree flatten & subtree\((v)\) is the contiguous preorder range
\([\tin(v),\tout(v)]\); run the same sweep on it for a non-invertible subtree
statistic. & \defref{c21-flatten}, \thmref{c21-subtree-range} \\
Root-path stack & CF 463E's queried set is the current root path: a plain DFS
 carries the window -- per-prime \(\mathit{last}[p]\) with \(O(1)\) push/undo,
 one rebuild per block of the \(\le50\) updates. \emph{No Mo, no lifting.} & \probref{c21-caisa}, \coderef{c21-caisa} \\
Path toggle & A doubled tour writes each node at entry and exit; crossing an
occurrence \textbf{toggles} membership, and a path is
\([st(u),st(v)]\) if \(u\) ancestors \(v\), else \([en(u),st(v)]\) \emph{plus the
LCA at read time}. & \defref{c21-toggle}, \thmref{c21-pathsplit} \\
Tarjan LCA & One DFS, three colours; union a returned child, stamp the parent; a
query fires when its second endpoint turns BLACK, reading
\(\mathit{anc}[\mathit{find}(\text{partner})]\). \(\bigO{(n+q)\alpha}\), offline only. &
\defref{c21-tarjan}, \thmref{c21-tarjan} \\
The pick & Offline + overlapping + non-invertible + \(O(1)\) update \(\Rightarrow\)
Mo; else Fenwick (invertible), HLD/LCT (online), DSU-on-tree (subtree-all),
virtual tree (per-set), centroid (all-pairs). & \defref{c21-signature}, \defref{c21-pick} \\
\bottomrule
\end{tabular}
\end{center}

\medskip
{\small\textbf{Algorithms from this chapter.}}
\begin{center}
\renewcommand{\arraystretch}{1.2}
\setlength{\tabcolsep}{3pt}
\begin{tabular}{@{}p{0.36\linewidth} l l >{\raggedright\arraybackslash}p{0.20\linewidth}@{}}
\toprule
\textbf{Algorithm} & \textbf{Prep} & \textbf{Per run} & \textbf{Code} \\
\midrule
Generic Mo block-order sweep & \(\bigO{q\log q}\) sort & \(\bigO{(n+q)\sqrt n}\) & \coderef{c21-mo} \\
CF 617E: Mo over prefix xors & \(\bigO{n}\) prefix & \(\bigO{(n+m)\sqrt n}\) & \coderef{c21-xor} \\
Subtree flatten + Mo on ranges & \(\bigO{n}\) DFS & \(\bigO{(n+q)\sqrt n}\) & \coderef{c21-flatten} \\
CF 463E: DFS stack + last-seen primes & \(\bigO{A}\) sieve & \(\bigO{51\,n\cdot 7+q}\) & \coderef{c21-caisa} \\
Doubled-tour path toggle + LCA split & \(\bigO{n}\) tour & \(\bigO{(n+q)\sqrt n}\) & \coderef{c21-euler-toggle} \\
Tarjan offline LCA & \(\bigO{n}\) DFS+DSU & \(\bigO{q\alpha}\) & \coderef{c21-tarjan} \\
Recognition gate / boundary pick & -- & -- & \defref{c21-signature}, \defref{c21-pick} \\
Mo kit (one decision listing) & -- & -- & \coderef{c21-kit} \\
CF 1585E (Summit): root-path freq-of-freq & \(\bigO{n}\) build & \(\bigO{(n+q)\log n}\) & \coderef{c21-summit} \\
\bottomrule
\end{tabular}
\end{center}

\chaptersection{Master Decision Flowchart}

Post-study ladder: walk top to bottom and stop at the first branch that matches.
It decides \emph{whether the problem is a Mo-on-trees job at all} (offline batch,
overlapping subtree/path sets, non-invertible statistic, \(O(1)\) add/remove) and,
if so, \emph{which array} the sweep runs on (preorder flatten for subtrees, the
doubled tour for paths). Everything that fails the top gate is handed to a
cheaper or online engine: invertible ranges to Euler \(+\) Fenwick (\(\chref{c24}\)),
online paths to HLD (\(\chref{c16}\)) / link-cut (\(\chref{c23}\)), per-set folds
to virtual trees (\(\chref{c30}\)), whole-tree distance shapes to centroid
(\(\chref{c18}\)), and pure LCA batches to Tarjan (\(\secref{c21-tarjan}\)).

\begin{lstlisting}[language={}, basicstyle=\ttfamily\small, numbers=none, frame=single]
  Are ALL queries known up front, over ONE fixed tree, with NO updates between them?
  |
  |-- An update or answer must come BETWEEN queries (online / interleaved)?
  |     +-- YES: NOT Mo -- block order needs the whole batch.
  |         Online path range -> HLD (Ch.22, c16); dynamic tree -> LCT (Ch.25, c23);
  |         online subtree range -> Euler tour + segment tree (Ch.5, c24).
  |
  |-- Does the statistic telescope: answer = PREF(r) - PREF(l-1)?
  |     +-- YES: NOT Mo -- pay a LOG, not a SQRT.
  |         Sum / xor / depth-parity-normalised add -> Fenwick over the flatten
  |         (Ch.5, c24).  Template: Ch.18's CF 383C (parity BIT, c14).
  |
  |-- Queries ask ONLY for lca(u,v), offline, and you want it near-linear?
  |     +-- YES: Tarjan DFS+DSU, O((n+q) alpha), no lifting table (27.4, lst:c21-tarjan).
  |         (If you ALSO need online LCAs, keep Ch.16 lifting.)
  |
  |-- Each query a SMALL marked set to fold independently, sum k << n?
  |     +-- YES: virtual tree (Ch.26, c30) -- the sets do not overlap as ranges.
  |
  |-- No query list: a whole-tree pair / distance distribution (sum over all paths)?
  |     +-- YES: centroid decomposition (Ch.24, c18).
  |
  |-- Subtree-only, statistic merges small-to-large, query set = (most) nodes?
  |     +-- YES: price DSU on tree (Ch.17, c28) -- O(n log n) can beat O((n+q) sqrt n).
  |
  |-- STILL HERE: offline, overlapping subtree/path ranges, non-invertible O(1) statistic
  |     -> THIS chapter.  Choose the array the sweep runs on:
  |        subtrees -> preorder flatten P, Mo on [tin(v), tout(v)]       (27.2, lst:c21-flatten)
  |        paths    -> doubled tour E, TOGGLE, split at lca + add LCA    (27.3, lst:c21-euler-toggle)
  |        arrays   -> block order (L/B, R), O(1) add/rm                 (27.1, lst:c21-mo)
\end{lstlisting}

\chaptersection{Common Mistakes Checklist}

\begin{itemize}[nosep, leftmargin=1.4em]
  \item \textbf{Paying the \(\sqrt n\) tax on an invertible statistic.} If the
        answer is \(\mathrm{pref}(r)-\mathrm{pref}(l-1)\), a Fenwick beats Mo by a
        \(\sqrt n/\log n\) factor.  Ch.~18's CF 383C \emph{looks} like Mo but telescopes
        after the \((-1)^{\depth}\) normalisation (\(\chref{c14}\)).
  \item \textbf{Forgetting the empty prefix \(px[0]\) in CF 617E.} A subarray from
        the left edge pairs \(px[l-1]\) with a later prefix; if \(L=l-1\) never
        adds index \(0\), every left-edge subarray is under-counted
        (\(\coderef{c21-xor}\)).
  \item \textbf{\texttt{rm} updating in the wrong order.} \texttt{rm} must
        \emph{decrement} the frequency \emph{before} subtracting the formed pairs,
        or it removes a pair that is itself leaving (\(\coderef{c21-xor}\)).
  \item \textbf{Sorting queries by \((R,\text{block of }L)\) instead of the reverse.}
        The bound needs one pointer monotone per group and the other confined to
        width \(B\); swapping the keys gives \(O(nq)\) again
        (\(\thmref{c21-complexity}\), \(\defref{c21-mo}\)).
  \item \textbf{Missing the LCA in path Mo (case ii).} When \(\lca(u,v)\neq u\)
        the window represents the path \emph{minus} the LCA; forget the read-time
        \texttt{insert(lca)} and answers go wrong exactly when the LCA carries a
        new colour (\(\thmref{c21-pathsplit}\), walk 3 of \(\secref{c21-path}\)).
  \item \textbf{Double-adding the LCA in case (i).} When \(u\) is the ancestor the
        LCA is \emph{already} oddly represented; toggling it again removes it.
        The extra insert is for case (ii) only (\(\coderef{c21-euler-toggle}\)).
  \item \textbf{Attaching a Tarjan query to only one endpoint.} It fires at the
        \emph{second} endpoint's death; if the first endpoint is where it lives, it
        checks a WHITE/GREY partner and is silently never answered
        (\(\defref{c21-tarjan}\), \(\coderef{c21-tarjan}\)).
  \item \textbf{Running Mo with interleaved updates.} Block order is defined only
        for a fixed array; a point update between queries invalidates the whole
        sweep -- that is the HLD / sqrt-decomposition-over-time contract
        (\(\defref{c21-signature}\), \(\secref{c21-boundary}\)).
\end{itemize}

\chaptersection{WA Clinic}

Verdicts are symptoms; the causes are the short list this chapter drilled. Diagnose
top-down -- every leaf names a listing.

\begin{figure}[H]
\centering
\begin{tikzpicture}[diagram,
    v/.style={boxnode, fill=primaryblue!12, draw=primaryblue, text width=22mm, align=center},
    b/.style={boxnode, fill=difforange!12, draw=difforange, text width=12mm},
    l/.style={boxnode, fill=successgreen!10, draw=successgreen, text width=48mm,
              font=\scriptsize, align=left}]
  \node[v] (verdict) {\textbf{Verdict on the judge}};
  \node[b, below left=8mm and 5mm of verdict] (wa) {\textbf{WA}};
  \node[b, below right=8mm and 5mm of verdict] (tle) {\textbf{TLE / RE}};
  \node[l, below right=1mm and 2mm of wa.south] (wa1)
    {path-Mo case (ii) read without the extra LCA, or case (i) toggled the LCA twice --
     wrong exactly on the queries where the LCA adds a colour (\(\thmref{c21-pathsplit}\), \(\coderef{c21-euler-toggle}\))};
  \node[l, below left=1mm and 2mm of wa.south] (wa2)
    {617E dropped \(px[0]\) or ran \texttt{rm} before decrementing the count -- off-by-pairs
     on left-edge subarrays (\(\coderef{c21-xor}\)); or 463E answered after pushing its
     own primes -- every match degenerates to \(v\) itself (\(\coderef{c21-caisa}\))};
  \node[l, below right=1mm and 2mm of tle.south] (tle1)
    {Mo on an invertible statistic (a Fenwick would be \(\bigO{\log n}\)) or block size
     \(B\) set to a constant far from \(n/\sqrt q\) -- the sweep degrades (\(\thmref{c21-complexity}\), \(\defref{c21-pick}\))};
  \node[l, below left=1mm and 2mm of tle.south] (tle2)
    {Mo fired on an online / update-interleaved query (block order illegal), or a recursion
     DFS overflowed at \(n=2\cdot10^5\) -- go HLD / iterative (\(\defref{c21-signature}\), \(\secref{c21-tarjan}\))};
  \draw[flow] (verdict) -- (wa);
  \draw[flow] (verdict) -- (tle);
  \draw[flow] (wa) -- (wa1);
  \draw[flow] (wa) -- (wa2);
  \draw[flow] (tle) -- (tle1);
  \draw[flow] (tle) -- (tle2);
\end{tikzpicture}
\caption*{WA: the path Mo LCA handled wrong, or the prefix / parity normalisation
mis-set. TLE or RE: paying \(\sqrt n\) where \(\log n\) applies, a badly chosen block
size, or Mo on an online problem (and a deep recursion at \(n=2\cdot10^5\)).}
\end{figure}

\chaptersection{Quick Reference Card}

\begin{quickref}[Chapter 27 at a glance]
\textbf{Mo on trees \(=\) reorder offline subtree/path queries so neighbours differ
by one node, then sweep a window with an \(O(1)\) add/remove; all \(q\) answers in
\(\bigO{(n+q)\sqrt n}\). For \(q\) LCAs only, Tarjan's DFS \(+\) DSU does the batch
in \(\bigO{(n+q)\alpha}\).}

\medskip
\textbf{Engine (27.1).} Block size \(B=n/\sqrt q\); sort by
\((\lfloor L/B\rfloor,R)\); creep \([curL,curR]\) with \texttt{add}/\texttt{rm};
read the statistic. CF 617E: window over prefix-xors, insert \(v\) adds
\(cnt[v\oplus k]\) pairs (\(\thmref{c21-complexity}\), \(\coderef{c21-mo}\),
\(\coderef{c21-xor}\)).

\medskip
\textbf{Subtrees (27.2).} Flatten: \(P[\tin(v)]=v\), subtree\((v)=[\tin(v),\tout(v)]\);
Mo on those ranges for a distinct/freq statistic. Invertible signed updates skip Mo
(Ch.~18's CF 383C parity BIT, \(\chref{c14}\)); root-path batches with \(\le50\)
updates ride the DFS stack: per-prime \(\mathit{last}[p]\) push/undo, one rebuild per
block (\(\defref{c21-flatten}\), \(\thmref{c21-subtree-range}\),
\(\coderef{c21-caisa}\)).

\medskip
\textbf{Paths (27.3).} Doubled tour \(E[1..2n]\); toggling gives represented\(=\)odd
occurrences; window \([st(u),st(v)]\) if \(u\) ancestors \(v\), else \([en(u),st(v)]\)
\emph{plus the LCA at read time} (\(\defref{c21-toggle}\), \(\thmref{c21-pathsplit}\),
\(\coderef{c21-euler-toggle}\)).

\medskip
\textbf{Tarjan (27.4).} WHITE/GREY/BLACK DFS; union a returned child,
\(\mathit{anc}[\mathit{find}(u)]=u\); when \(u\) turns BLACK answer its queries whose
partner is BLACK via \(\mathit{anc}[\mathit{find}(\text{partner})]\)
(\(\thmref{c21-tarjan}\), \(\coderef{c21-tarjan}\)).

\medskip
\textbf{The pick (27.5).} Offline \emph{and} overlapping \emph{and} non-invertible
\emph{and} \(O(1)\) update \(\Rightarrow\) Mo; invertible \(\to\) Fenwick
(\(\chref{c24}\)); online \(\to\) HLD/LCT (\(\chref{c16},\chref{c23}\)); per-set
\(\to\) virtual tree (\(\chref{c30}\)); subtree-all \(\to\) DSU-on-tree
(\(\chref{c28}\)); all-pairs \(\to\) centroid (\(\chref{c18}\)); LCA-batch \(\to\)
Tarjan (\(\defref{c21-signature}\), \(\defref{c21-pick}\)).
\end{quickref}

\chaptersection{Write-From-Memory Drills}

Five signature skeletons, each lifted from a listing already in the chapter. Cover the
marked line, reproduce from memory, then check the answer one line below.

\begin{lstlisting}[caption={Drill 1 -- the Mo block-order sweep (from lst:c21-mo)}]
sort(queries, cmp);                    // cmp: block of L, then R
curL = 0; curR = -1;                   // start EMPTY
for (qu : queries){
    while (curL > qu.l) add(--curL);   while (curR < qu.r) add(++curR);
    while (curL < qu.l) rm(curL++);    while (curR > qu.r) rm(curR--);
    ans[qu.idx] = cur;
}
\end{lstlisting}
\noindent{\scriptsize Key: grow first, then shrink, in that order keeps the window
valid; \texttt{add}/\texttt{rm} must be \(O(1)\) and mutually inverse; total movement
\(\bigO{(n+q)\sqrt n}\) with \(B=n/\sqrt q\) (\(\thmref{c21-complexity}\)).}

\begin{lstlisting}[caption={Drill 2 -- CF 617E add/rm over prefix xors (from lst:c21-xor)}]
void add(int t){ int v=px[t]; cur += cnt[v ^ k]; cnt[v]++; }   // pairs formed, then store
void rm (int t){ int v=px[t]; cnt[v]--;  cur -= cnt[v ^ k]; }   // un-store, then pairs lost
// query (l,r) -> window [l-1, r] over px[0..n]; remember px[0]=0
\end{lstlisting}
\noindent{\scriptsize Key: a subarray xor is \(px[q]\oplus px[p]\); inserting \(v\)
creates \(cnt[v\oplus k]\) good pairs; \texttt{rm} decrements \emph{first} so the
leaving element is not counted against itself; \(cur\) is \code{long long}
(\(\probref{c21-xor}\)).}

\begin{lstlisting}[caption={Drill 3 -- subtree flatten + Mo range (from lst:c21-flatten)}]
void dfs(int u,int p){ tin[u]=++t; P[t]=u;
    for(int v:adj[u]) if(v!=p){ dep[v]=dep[u]+1; dfs(v,u);} tout[u]=t; }
// subtree(v) aggregate == Mo over array range [ tin[v] , tout[v] ] on P
//   add(pos) inserts node P[pos] into the statistic;  rm(pos) erases it.
\end{lstlisting}
\noindent{\scriptsize Key: preorder makes every subtree one contiguous block, so the
array engine runs unchanged on \([\tin(v),\tout(v)]\); only \texttt{add} now touches
\(P[\text{pos}]\) (\(\defref{c21-flatten}\), \(\thmref{c21-subtree-range}\)).}

\begin{lstlisting}[caption={Drill 4 -- path window + toggle (from lst:c21-euler-toggle)}]
if (st[u] > st[v]) swap(u,v);                  // keep st[u] < st[v]
if (is_ancestor(u,v)){ L=st[u]; R=st[v]; extraLca=false; }
else                 { L=en[u]; R=st[v]; extraLca=true; }   // window misses the lca
// add(pos): x=eu[pos]; held[x]^=1; (held[x]? insert(x): erase(x));   rm = add
// if extraLca:  insert(lca(u,v)); read cur; erase(lca(u,v));
\end{lstlisting}
\noindent{\scriptsize Key: a path is a doubled-tour window whose represented set is the
odd-occurrence nodes; when \(u\) is the ancestor the LCA is already inside, otherwise it
is \emph{zero} occurrences and must be added at read time (\(\thmref{c21-pathsplit}\)).}

\begin{lstlisting}[caption={Drill 5 -- Tarjan offline LCA core (from lst:c21-tarjan)}]
void dfs(int u,int p){
  col[u]=1; fa[u]=u; anc[u]=u;
  for(v : adj[u]) if(v!=p){ dfs(v,u); fa[find(v)]=find(u); anc[find(u)]=u; }
  for(Q : qlist[u]) if(col[Q.other]==BLACK) ans[Q.id]=anc[find(Q.other)];
  col[u]=BLACK;
}   // attach every query (u,v) to BOTH u and v
\end{lstlisting}
\noindent{\scriptsize Key: after each child returns its finished region merges up and
re-stamps \(u\); the query fires at the second endpoint's death, reading the stamp of
the partner's set -- the merge point of the two paths (\(\thmref{c21-tarjan}\)).}

\chaptersection{Complete Problem Checklist}

Every problem this chapter offers -- two fully solved in-section, worked traces on the
running example \(R\), and one Summit on its own CF samples. Tick \emph{only} after an
unaided solve; a hinted attempt counts as half.

\begin{center}
\footnotesize
\setlength{\tabcolsep}{4pt}
\renewcommand{\arraystretch}{1.25}
\begin{tabular}{@{}c p{0.16\linewidth} l p{0.20\linewidth} >{\raggedright\arraybackslash}p{0.26\linewidth} c@{}}
\toprule
\textbf{\#} & \textbf{Problem} & \textbf{Source} & \textbf{Key technique} & \textbf{What to look for} & \textbf{Done?} \\
\midrule
\multicolumn{6}{l}{\textbf{Section 27.1 -- the Mo engine (array)}} \\
W27.1 & block order on \(n{=}9,B{=}3\) & worked & \lv{3} sort \((L/B,R)\), creep & 13 moves; \(R\) monotone per block, \(L\) jumps \(\le B\) & \(\square\) \\
W27.1b & XOR and Favorite Number & CF 617E & \lv{4} Mo over prefix xors & \(px[0]\) counted; \texttt{rm} decrements first; samples \(\to\) 7,0 and 9,4,4 & \(\square\) \\
\midrule
\multicolumn{6}{l}{\textbf{Section 27.2 -- subtrees}} \\
W27.2 & subtree\((2)\) of \(R\) as a range & worked & \lv{3} preorder flatten & \(\subtree(2)=P[2..5]=\{2,4,8,5\}\); three distinct colours & \(\square\) \\
W27.2b & Caisa and Tree & CF 463E & \lv{4} DFS stack + last-seen primes & read before push (strict ancestor); \(\le50\) updates \(\to\) 51 rebuilds; sample \(\to\) -1,1,2,-1,1 & \(\square\) \\
\midrule
\multicolumn{6}{l}{\textbf{Section 27.3 -- paths (worked only)}} \\
W27.3 & paths \((4,8),(8,6),(2,3)\) on \(R\) & worked & \lv{5} toggle + LCA split & \((4,8)\to2\); \((8,6)\to3\) with LCA; \((2,3)\to3\) vs lying 2 & \(\square\) \\
\midrule
\multicolumn{6}{l}{\textbf{Section 27.4 -- Tarjan (worked only)}} \\
W27.4 & five LCAs on \(R\) & worked & \lv{4} DFS + DSU + colours & union-on-return; fire at 2nd BLACK; answers \(1,2,2,3,1\) & \(\square\) \\
\midrule
\multicolumn{6}{l}{\textbf{Section 27.5 -- recognition / boundary (worked only)}} \\
W27.5 & three specs, three verdicts & worked & \lv{5} the four tests & path-distinct \(\to\) Mo; subtree-all \(\to\) DSU; updates \(\to\) reject & \(\square\) \\
\midrule
\multicolumn{6}{l}{\textbf{Summit (chapter end)}} \\
W27 & Frequency Queries & CF 1585E & \lv{6} root-path freq-of-freq & DFS stack window + BIT k-th; samples \(\to\) 1 -1 1 / 1 1 2 1 -1 & \(\square\) \\
\bottomrule
\end{tabular}
\end{center}

\chaptersection{Recommended Study Order}

Four focused days. Day~1 the engine and CF 617E; day~2 the flatten, CF 463E and its
contrast to Mo; day~3 the doubled-tour toggle and Tarjan; day~4 the boundary pick, the
Summit and the mixed set.

\begin{center}
\small
\renewcommand{\arraystretch}{1.25}
\setlength{\tabcolsep}{3pt}
\begin{tabular}{@{}llp{0.55\linewidth}@{}}
\toprule
\textbf{Day} & \textbf{Focus} & \textbf{Tasks} \\
\midrule
1 & Read + do 27.1 & run the block order on \(n{=}9,B{=}3\) to 13 moves (\(\coderef{c21-mo}\)); prove \(\bigO{(n+q)\sqrt n}\) (\(\thmref{c21-complexity}\)); then CF 617E, replaying the prefix-xor window to \(7,0\) and \(9,4,4\) (\(\probref{c21-xor}\)) \\
2 & Read + do 27.2 & flatten subtree\((2)\) of \(R\) and read its three colours (\(\coderef{c21-flatten}\)); then CF 463E, carrying \(\mathit{last}[p]\) on the DFS stack and replaying to \(-1,1,2,-1,1\) (\(\probref{c21-caisa}\)) -- then say out loud why it needed \emph{neither} Mo nor lifting \\
3 & Read + do 27.3, 27.4 & walk the three path windows and the LCA-parity trap (\(\thmref{c21-pathsplit}\)); run the five-query Tarjan log to \(1,2,2,3,1\) (\(\thmref{c21-tarjan}\)) \\
4 & Read + do 27.5, Summit & place the three specs on the pick (\(\defref{c21-pick}\)); then CF 1585E, replaying both samples (\(\probref{c21-summit}\)); finish with the first 3 mixed entries \\
\bottomrule
\end{tabular}
\end{center}

\chaptersection{Difficulty Progression}

Tiers are competency gates -- each row names the \emph{skill} the tier unlocks.

\begin{center}
\small
\renewcommand{\arraystretch}{1.3}
\setlength{\tabcolsep}{3pt}
\begin{tabular}{@{}lp{0.30\linewidth}p{0.42\linewidth}@{}}
\toprule
\textbf{Tier} & \textbf{Problems} & \textbf{What passing this tier proves} \\
\midrule
\lv{3} \ \(\bigstar\bigstar\bigstar\)
  & block order by hand; subtree flatten
  & you can creep a window, count pointer moves, and read \([\tin(v),\tout(v)]\) \\
\lv{4} \ \(\bigstar\bigstar\bigstar\bigstar\)
  & CF 617E, CF 463E, Tarjan batch
  & you code an \(O(1)\) \texttt{add}/\texttt{rm}, carry a push/undo table on the DFS stack, and fire a Tarjan query at the right death \\
\lv{5} \ \(\bigstar\bigstar\bigstar\bigstar\bigstar\)
  & path toggle + LCA split, the boundary pick
  & you represent a path as a doubled-tour window, handle the parity trap, and \emph{reject} Mo to the right tool \\
\lv{6} \ \(\bigstar\bigstar\bigstar\bigstar\bigstar\bigstar\)
  & CF 1585E (Summit), mixed at CF~2300--2400
  & you recast the DFS stack as the Mo window and select a k-th-by-frequency answer with a frequency-of-frequency structure \\
\bottomrule
\end{tabular}
\end{center}

\chaptersection{Reduction Gallery}

Every entry is ``one question shape \(\to\) one move (or an honest hand-off).'' Each
line collapses a problem family onto a listing already in the chapter.

\begin{itemize}[nosep, leftmargin=1.4em]
  \item ``count subarrays with \(\mathit{xor}=k\), offline'' \(\to\) prefix-xor pairs,
        Mo window over \(px[0..n]\), \(O(1)\) insert/delete
        (\(\probref{c21-xor}\), \(\coderef{c21-xor}\)).
  \item ``distinct colours / \(\sum\mathrm{freq}^2\) over subtree\((v)\), offline''
        \(\to\) preorder flatten + Mo on \([\tin(v),\tout(v)]\)
        (\(\thmref{c21-subtree-range}\), \(\coderef{c21-flatten}\); CF 375D is the
        sibling -- its URL lives in Ch.17, \chref{c28}).
  \item ``add with sign flipping by depth, read a node'' (CF 383C) \(\to\) multiply
        every stored quantity by \((-1)^{\depth}\), one difference BIT
        (solved in Ch.~18, \(\chref{c14}\)).
  \item ``deepest ancestor sharing a factor, rare updates'' (CF 463E) \(\to\) the
        DFS stack \emph{is} the root-path window; per-prime last-seen push/undo,
        one rebuild per static block (\(\probref{c21-caisa}\), \(\coderef{c21-caisa}\)).
  \item ``distinct colours on path \(u\)-\(v\), offline'' \(\to\) doubled-tour toggle
        window + read-time LCA insert (\(\thmref{c21-pathsplit}\),
        \(\coderef{c21-euler-toggle}\)).
  \item ``answer \(q\) LCAs offline, no lifting table'' \(\to\) Tarjan DFS \(+\) DSU,
        fire at the second endpoint's death (\(\thmref{c21-tarjan}\),
        \(\coderef{c21-tarjan}\)).
  \item ``frequency histogram of the root-to-\(v\) path, k-th by count'' (CF 1585E)
        \(\to\) the DFS stack \emph{is} the window; cnt[val] + freq-of-frequency BIT
        (\(\probref{c21-summit}\), \(\coderef{c21-summit}\)).
  \item ``invertible range sum/xor on a subtree or path'' \(\to\) \emph{not} Mo:
        Euler \(+\) Fenwick / segment (\(\chref{c24}\), \(\defref{c21-pick}\)).
  \item ``online path range with updates'' \(\to\) \emph{not} Mo: HLD (\chref{c16})
        or link-cut (\chref{c23}); ``per-set fold'' \(\to\) virtual tree (\chref{c30}).
  \item ``subtree statistic for every node, mergeable'' \(\to\) price DSU-on-tree
        (\chref{c28}); ``all-pairs distance'' \(\to\) centroid (\chref{c18})
        (\(\secref{c21-boundary}\)).
\end{itemize}

\chaptersection{Mixed Problem Set}

Ten entries, each live-verified against the current Codeforces problem set (names,
ratings and tags re-pulled from the \texttt{problemset.problems} API this session) and
absent from every other roster in the book (the canonical tree-Mo problems CF 375D in
Ch.~17, CF 258E and CF 86D are owned elsewhere and not re-linked here). None says
``use Mo''; several deliberately \emph{do not} want it, and the drill is to run them
through the four tests (\(\defref{c21-signature}\)) and say so. Each hint names the
reduction (or the rejection) and the listing it leans on, ordered ascending by rating.

\begin{mixedset}
  \practice{Alyona and a tree}{CF 739B}
    {https://codeforces.com/problemset/problem/739/B}{CF 1900 L4}
    {Tags \emph{binary search, data structures, dfs, graphs, trees}. For each node \(v\),
     count ancestors \(u\) with \(\dist(u,v)\le a_v\). This is a \emph{root-path} sweep --
     the DFS stack \emph{is} the window (the Summit's shape, \(\probref{c21-summit}\)) --
     but the count telescopes along the depth axis: binary-lift to the highest valid
     ancestor, then a difference/BIT over depth. So it lives in world~1 (invertible,
     \(\secref{c21-subtree}\)) rather than Mo: \(\bigO{n\log n}\), not
     \(\bigO{(n+q)\sqrt n}\). Price both; the honest answer is the depth-difference.}
  \practice{Tree Queries}{CF 1328E}
    {https://codeforces.com/problemset/problem/1328/E}{CF 1900 L4}
    {Each query names \(a,b,x,y\): does \emph{some} path from \(a\) to \(b\) use the tree
     edge \((x,y)\)? Root once, orient the edge to its deeper end \(c\); the edge is on
     the \(a\)-\(b\) path iff exactly one of \(a,b\) lies inside subtree\((c)\) -- two
     \(\tin/\tout\) containment tests. Online per query and prefix-shaped: test~1 fails
     and test~3 fails (\(\defref{c21-signature}\)); reject Mo, keep
     \(\chref{c24}+\chref{c11}\). The same containment arithmetic 27.2's flatten hands
     over for free.}
  \practice{Passable Paths (easy version)}{CF 1702G1}
    {https://codeforces.com/problemset/problem/1702/G1}{CF 1900 L4}
    {Each query names a \emph{set} of vertices: is one simple path passable through all
     of them? Pick the set's two farthest vertices (diameter by \(\dist\) via
     \(\lca\)) and test every \(w\) for \(\dist(x,w)+\dist(w,y)=\dist(x,y)\) -- a
     membership test in pure distance arithmetic (\(\thmref{c11-distf}\)). Queries are
     independent checks, not overlapping ranges to amortise (test~2 empty) and nothing
     histogram-like is carried (test~3 fails). Reject; go to LCA (\(\chref{c11}\)).}
  \practice{Number of Simple Paths}{CF 1454E}
    {https://codeforces.com/problemset/problem/1454/E}{CF 2000 L4}
    {A connected graph on \(n\) vertices and \(n\) edges -- exactly one cycle -- count
     ordered pairs joined by \emph{exactly one} simple path. There is \emph{no query
     list} at all (test~2 empty): the answer is one global count from the cycle's
     attached subtree sizes, a parity-free inclusion-exclusion over the cycle.
     Structurally the cleanest rejection in the set -- count by structure, never sweep
     (\(\defref{c21-pick}\); subtree sizes from \(\chref{c24}\)). Use it to feel that
     ``unicyclic graph'' does \emph{not} mean offline queries.}
  \practice{Beautiful Subarrays}{CF 665E}
    {https://codeforces.com/problemset/problem/665/E}{CF 2100 L5}
    {Count subarrays with at most \(k\) even elements. The set-up is the closest call in
     the list -- offline, overlapping ranges over one array -- but the statistic is a
     count of evens, fully \emph{invertible}: a two-pointer window slides it in
     \(\bigO n\), a prefix count or BIT in \(\bigO{n\log n}\). Test~3 fails, so even
     \(\bigO{(n+q)\sqrt n}\) Mo is a tax on the array version of the lesson 27.2 taught
     with CF~383C (\(\defref{c21-signature}\)). Drill: say \emph{why} the window here
     never needs a block sort.}
  \practice{Timofey and a Black-White Tree}{CF 1790F}
    {https://codeforces.com/problemset/problem/1790/F}{CF 2100 L5}
    {Nodes start white and turn black one by one; at time \(i\), before recolouring
     node~\(i\), report how many black nodes are reachable from \(i\) through black
     nodes only. This is component \emph{size} under a growing activation prefix --
     connectivity, not a reorderable range: the answers must respect the time prefix
     (test~1 fails). The standard move is reverse-time DSU merging black nodes for the
     late times plus a thresholded direct DFS for the early ones (sqrt-decomposition
     over \emph{time}, not over ranges). Keep the DSU instinct from
     \(\secref{c21-tarjan}\); reject Mo's block order (\(\defref{c21-pick}\)).}
  \practice{We Need More Bosses}{CF 1000E}
    {https://codeforces.com/problemset/problem/1000/E}{CF 2100 L5}
    {Query \((s,t)\): the maximum number of bridges any walk from \(s\) to \(t\) can
     force. Compress the 2-edge-connected components; what remains is a \emph{bridge
     tree}, and the answer is just the number of tree-edges on the \(s\)-\(t\) path --
     a depth-prefix sum with an \(\lca\) correction. The distance decomposes over
     prefixes (test~3 fails) and each query answers online: reject Mo, keep bridges
     \(+\) lifting (\(\chref{c11}\)). The compression step is the chapter's one
     inversion: run Tarjan-family DFS ideas \emph{for} the structure, not against it.}
  \practice{Tree Master}{CF 1806E}
    {https://codeforces.com/problemset/problem/1806/E}{CF 2200 L5}
    {Tags \emph{brute force, data structures, dfs, trees}. Feed it to the four tests
     before coding: is the whole query set known up front, do the queried sets overlap,
     does the statistic refuse to telescope, and is one element cheap to add and drop?
     A \emph{data structures + dfs} tree problem with a \emph{brute force} tag is the
     tell for ``a fold that a smarter structure beats'' -- so first look for the
     invertible/structural route (Fenwick over the flatten, or DSU-on-tree for
     subtree-all, \(\chref{c28}\)); only if all four tests fire does Mo earn its
     \(\sqrt n\) (\(\defref{c21-signature}\), \(\secref{c21-boundary}\)).}
  \practice{Weights Division (hard version)}{CF 1399E2}
    {https://codeforces.com/problemset/problem/1399/E2}{CF 2200 L5}
    {Tags \emph{binary search, dfs and similar, greedy, sortings, trees}. Weighted
     tree rooted at \(1\), threshold \(S\): repeatedly halve an edge (floor), each
     halving costs \(1\), and every root-to-leaf sum must end up \(\le S\); minimise
     the operations. There is \emph{no query list} to sweep (test~2 empty) and no
     window statistic at all: the leverage of halving edge \(e\) is
     \(\mathit{leaves}(e)\cdot(w-\lfloor w/2\rfloor)\) -- a bottom-up leaf-count
     fold plus a greedy/priority queue (the tree-DP family, \(\chref{c10}\)).
     Reject Mo the way 27.5 rejects every fold-shaped objective
     (\(\defref{c21-pick}\)); the \(\sqrt n\) has nothing to amortise here.}
  \practice{Yasya and the Mysterious Tree}{CF 1980G}
    {https://codeforces.com/problemset/problem/1980/G}{CF 2300 L5}
    {Tags \emph{bitmasks, data structures, dfs, graphs, greedy, strings, trees}. Path
     questions where the state per edge is a small \emph{parity bitmask}. A bitmask fold
     along \(u\)-to-\(v\) is invertible (xor of root-to-node masks, \(\lca\) correction),
     so a single path query is answered in \(O(\log)\) from prefixes -- the opposite of
     the ``no inverse'' Mo regime. Gate test~3 fails; use LCA + bitmask prefix
     (\(\chref{c11}\)). Reach for Mo only if the statistic were a distinct-count, which
     it is not (\(\defref{c21-signature}\)).}
\end{mixedset}

\chaptersection{The Summit -- CF 1585E Frequency Queries}

\begin{notebox}[Why this is the chapter top rope]
Every problem before it put a \emph{histogram} on a subtree or a path and let Mo sweep.
The Summit asks for something finer: on the path \(v\to\) root, look at the
\emph{frequency of each value}, keep the values seen at least \(l\) times, and report the
\(k\)-th one ordered \emph{by frequency}. That is a statistic on a statistic -- it has no
inverse and no single running scalar, so naive Mo (which carries one number per window)
does not fit. The rescue is the deepest idea of the chapter wearing a new hat: a plain
DFS \emph{already} walks the root-to-\(v\) path as a stack, pushing a value on entry and
popping it on exit -- the exact \(O(1)\) \texttt{add}/\texttt{rm} of \(\secref{c21-moengine}\),
now riding the recursion instead of two pointers. Layer a frequency-of-frequency Fenwick
on top and the \(k\)-th-by-count is a select in \(O(\log n)\). It reuses \textbf{two}
tools of this chapter (the Mo frequency table and Tarjan's ``answer offline queries
attached to a node during one DFS'') plus \textbf{one} from Ch.~5 (the depth/parent DFS),
and is hand-replayed to both judge outputs \texttt{1 -1 1} and \texttt{1 1 2 1 -1}.
\end{notebox}

\solvedproblem{c21-summit}{Frequency Queries}{CF 1585E}
  {https://codeforces.com/problemset/problem/1585/E}{CF 2400 L6}
  {one DFS from the root turns the recursion stack into the Mo window (push a value on
   entry, pop on exit -- \(O(1)\) add/remove); maintain \(\mathit{cnt}[x]\) = occurrences
   of value \(x\) on the current root path and a frequency-of-frequency structure
   \(\mathit{bucket}[f]\) = number of distinct values seen exactly \(f\) times, kept as a
   Fenwick so ``\(k\)-th distinct value by increasing frequency among those with freq
   \(\ge l\)' ' is a prefix-sum select \(O(\log n)\); a per-frequency value stack returns
   the answer value; queries are attached to their node and answered at that DFS instant
   (Tarjan's offline discipline)}

\begin{problemstatement}[Frequency Queries (statement, constraints and both samples verified against the live CF page)]
Vertex \(1\) roots a tree on \(n\) nodes; node \(i\) carries an integer \(a_i\)
(\(1\le a_i\le n\)). There are \(q\) questions; question \((v,l,k)\) is answered by:
write the sequence of values on the shortest path from \(v\) up to the root (both
included); count how many times each value occurs; \textbf{drop} every value occurring
fewer than \(l\) times; \textbf{deduplicate} and \textbf{order the remaining distinct
values by their occurrence count in the original path sequence, increasing} (ties in any
order); the answer is the \(k\)-th value in that ordered list, or \(-1\) if fewer than
\(k\) values survive.

\textbf{Input.} Multiple test cases, total \(\sum n,\sum q\le 10^6\), \(1\le t\le10^6\).
Per case: \(n,q\) (\(1\le n,q\le10^6\)); the \(n\) values \(a_1..a_n\); the \(n-1\)
parents \(p_2..p_n\); then \(q\) lines \(v,l,k\) (\(1\le v,l,k\le n\)).

\textbf{Output.} For each question print its answer; ties may print any valid value.

\textbf{Samples (both official).} Case~1: \(n{=}3,q{=}3\), \(a{=}[1,1,1]\), parents
\([1,2]\) (chain \(1\!-\!2\!-\!3\)), questions \((3,1,1),(3,1,2),(3,2,1)\) \(\to\)
\texttt{1 -1 1}. Case~2: \(n{=}5,q{=}5\), \(a{=}[1,2,1,1,2]\), parents \([1,1,2,2]\),
questions \((3,1,1),(2,1,2),(4,1,1),(4,2,1),(4,2,2)\) \(\to\) \texttt{1 1 2 1 -1}.
\end{problemstatement}

\begin{attemptgate}[Attempt it yourself first -- \lv{6}, budget a real sitting]
The queried set is always a root path, and you need the whole frequency \emph{histogram}
of that path, not one scalar.  Why does that break plain Mo, and what does a DFS already
give you for free?
\peek{a root-to-\(v\) path is exactly the node-stack of a DFS when the search is at \(v\).
Enter a node: its value's count goes up by one. Leave: down by one. That is an \(O(1)\)
\texttt{add}/\texttt{rm} -- the Mo primitive -- but the ``window'' is the recursion depth,
so you need no block order at all.}
\nudge{Attach every question to its node \(v\) and answer it the instant the DFS is at
\(v\) (Tarjan's offline discipline, \(\secref{c21-tarjan}\)). Maintain
\(\mathit{cnt}[\mathit{val}]\) and \(\mathit{bucket}[f]=\#\{\text{values with }\mathit{cnt}=f\}\).
The ordered list you must select from is ``values with \(\mathit{cnt}\ge l\), by
\(\mathit{cnt}\) up'' -- a selection over the \(\mathit{bucket}\) array.}
\fullhint{Put \(\mathit{bucket}[\cdot]\) in a Fenwick so its prefix sum
\(S(f)=\sum_{g\le f}\mathit{bucket}[g]\) is ``\#distinct values with freq \(\le f\)'' and
advances by \(\pm1\) whenever a value moves between frequencies. For \((v,l,k)\):
\(\mathit{base}=S(l-1)\), \(\mathit{total}=S(n)\); if \(\mathit{total}-\mathit{base}<k\)
answer \(-1\); else find the smallest \(f\) with \(S(f)\ge \mathit{base}+k\) (Fenwick
lower-bound, \(O(\log n)\)) and output \emph{any} value whose current count is \(f\) --
maintain a per-frequency swap-remove stack \(\mathit{valAt}[f]\) to read one in \(O(1)\).
Moving value \(x\) from count \(c\) to \(c{+}1\): \(\mathit{bit}.\mathrm{add}(c,-1),
\mathit{bit}.\mathrm{add}(c{+}1,+1)\) and relocate \(x\) between \(\mathit{valAt}[c]\)
and \(\mathit{valAt}[c{+}1]\) (skip the \(c=0\) add). On leaving, mirror it with \(c\to
c{-}1\). One DFS answers all \(q\) in \(\bigO{(n+q)\log n}\); the structures fully unwind
to empty when the DFS returns, so the next test case starts clean.}
\end{attemptgate}

\solstep{Understanding}
Two reframings collapse the problem. First, ``the sequence on the path \(v\to\) root'' is
not an arbitrary range needing a doubled tour -- it is the \emph{ancestor stack} the DFS
holds while it sits at \(v\). So instead of Mo's two pointers creeping over a static
array, the window is the current recursion path, and it is entered/exited one node at a
time for free: entering \(x\) does \(\mathit{cnt}[a_x]{+}{+}\), leaving does
\(\mathit{cnt}[a_x]{-}{-}\), each an \(O(1)\) add/remove of exactly the \(\secref{c21-moengine}\)
flavour. Second, the answer is a \emph{selection on the frequency histogram}: among
distinct values with \(\mathit{cnt}\ge l\), ordered by \(\mathit{cnt}\) ascending, take
the \(k\)-th. That histogram is captured by \(\mathit{bucket}[f]\) -- ``how many distinct
values occur exactly \(f\) times'' -- and the selection is a prefix-sum lower-bound over
\(\mathit{bucket}\). A Fenwick maintains \(\mathit{bucket}\) under \(\pm1\) moves and
answers the select in \(O(\log n)\); the offline ``answer each question at its node''
discipline is Tarjan's (\(\secref{c21-tarjan}\)), not Mo's block order. The Summit is
therefore the chapter's thesis -- cheap single-element updates plus a non-invertible
histogram -- solved by the \emph{other} window mechanism this chapter introduced.

\solstep{Key Observations}
\begin{observation}
\textbf{A root path is a DFS stack.} While the DFS is at \(v\), the multiset of values
with \(\mathit{cnt}>0\) is exactly the values on the path \(v\to\) root. Entering/leaving
a node changes \(\mathit{cnt}\) of one value by one, so the window updates are \(O(1)\)
and come in matched enter/leave pairs -- everything unwinds to empty at the end of the
DFS, which is why no explicit reset is needed between test cases.
\end{observation}
\begin{observation}
\textbf{Ordering by frequency is selection on \(\mathit{bucket}\).} The deduplicated,
freq-increasing list is the values grouped by their count \(f=l,l+1,\dots\), group \(f\)
holding \(\mathit{bucket}[f]\) values. The \(k\)-th survivor is any value in the group
\(f^\ast\) where \(f^\ast\) is the smallest \(f\ge l\) with
\(\sum_{g=l}^{f}\mathit{bucket}[g]\ge k\). With \(S(f)=\sum_{g\le f}\mathit{bucket}[g]\)
that is ``smallest \(f\) with \(S(f)\ge S(l-1)+k\)'': a Fenwick lower-bound in \(O(\log n)\),
and ties let us emit any member of group \(f^\ast\) (read from \(\mathit{valAt}[f^\ast]\)).
\end{observation}

\solstep{Think It With Me}
\begin{thinkbox}
\deadend{Run Mo on paths: treat the root path as a doubled-tour window}
{Path Mo (\(\secref{c21-path}\)) answers a statistic that is a single running scalar
(distinct count, \(\sum\mathrm{freq}^2\)) -- you read one number when the window is right.
Here the answer needs the \emph{entire} frequency histogram re-sorted and a \(k\)-th
element selected, which Mo's scalar window cannot express; you would re-scan all
frequencies per query (\(\bigO{n}\) each, \(\tle\)). And a root path is already
contiguous in the ancestor stack -- the doubled-tour toggle is needless overhead.
\textbf{TLE / wrong model}.}
\deadend{Keep \(\mathit{bucket}[f]\) in a plain array and scan \(f=l,l+1,\dots\) per query}
{Each query then costs \(O(\text{number of distinct frequencies})=O(n)\) in the worst
case, and \(q\le 10^6\), so \(O(nq)\). The cumulative-count selection over
\(\mathit{bucket}\) is a \emph{prefix-sum lower-bound} -- exactly what a Fenwick does in
\(O(\log n)\). You must promote the histogram array to a searchable structure.
\textbf{TLE}.}
\deadend{Recompute the path value multiset for every query from \(v\) up through parents}
{Walking \(v\to\) root per query is \(\bigO{\depth}\) and rebuilding a frequency map from
scratch kills \(q\) -- no work is shared between sibling queries. The whole point of the
DFS view is that moving from a node to its child changes the histogram by \emph{one}
value; re-walking per query throws that amortisation away. \textbf{TLE}.}
\aha{Do one DFS from the root and let the recursion stack be the window: \(\mathit{enter}(x)\)
bumps \(\mathit{cnt}[a_x]\) (moving value \(a_x\) from frequency bucket to bucket in a
Fenwick and between two \(\mathit{valAt}\) stacks) and \(\mathit{leave}(x)\) undoes it.
Answer every query attached to \(v\) at the DFS instant: \(\mathit{base}=S(l-1)\),
\(\mathit{total}=S(n)\), if \(\mathit{total}-\mathit{base}<k\) print \(-1\), else
Fenwick-lower-bound for the smallest \(f\) with \(S(f)\ge\mathit{base}+k\) and print any
member of \(\mathit{valAt}[f]\). All in \(\bigO{(n+q)\log n}\).}
\end{thinkbox}

\solstep{The Approach Ladder}
\begin{center}\renewcommand{\arraystretch}{1.25}\setlength{\tabcolsep}{3pt}
\begin{tabular}{@{}c p{0.23\linewidth} p{0.29\linewidth} l c@{}}
\toprule
\textbf{\#} & \textbf{Approach} & \textbf{Idea} & \textbf{Time} & \textbf{Verdict}\\
\midrule
1 & \lv{3}\ walk \(v\to\) root per query & rebuild freq map from scratch & \(\bigO{q\cdot n}\) & \tle \\
2 & \lv{4}\ path Mo, doubled tour & scalar window, wrong shape & \(\bigO{(n+q)\sqrt n}\) & \textcolor{warningred}{\ding{55}\ \textbf{WA}} \\
3 & \lv{5}\ DFS stack window + array histogram & right window, linear select & \(\bigO{(n+q)\cdot n}\) & \tle \\
4 & \lv{6}\ \textbf{DFS stack + freq-of-freq Fenwick} & push/pop + lower-bound select & \(\bigO{(n+q)\log n}\) & \ac \\
\bottomrule
\end{tabular}\end{center}
Rung~4 is the Summit: rungs~1--3 each keep the right \emph{goal} but the wrong
\emph{object} (a re-walked path, a scalar Mo window, an unsearchable histogram); realising
that a root path \emph{is} a DFS stack and that ``\(k\)-th by frequency'' \emph{is} a
Fenwick lower-bound over the frequency-of-frequency array is the whole win.

\solstep{Solution}
Build child lists from the parent array. DFS from node \(1\): on \(\mathit{enter}(u)\)
move value \(a_u\) from its current frequency to the next (\(\mathit{cnt}\) up), applying
the matched \(\mathit{bucket}\) Fenwick \(-1/+1\) and relocating \(a_u\) between
\(\mathit{valAt}\) stacks; then answer every query attached to \(u\) by
\(\mathit{base}=S(l-1),\ \mathit{total}=S(n)\), emitting \(-1\) if
\(\mathit{total}-\mathit{base}<k\) else the value \(=\mathit{valAt}[f^\ast].\!\)back where
\(f^\ast=\)Fenwick-lower-bound\((\mathit{base}+k)\). Recurse into children, then
\(\mathit{leave}(u)\) mirrors the move. Because every enter is matched by a leave, all
structures return to empty when the DFS ends, so per test case you only clear the
adjacency and query lists. Read answers in input order.

\solstep{Proof of Correctness}
\begin{proofbox}
\textbf{(Window $=$ path.)} At the DFS instant at \(v\), the entered-but-not-yet-left nodes
are exactly the root-to-\(v\) stack, so \(\{x:\mathit{cnt}[x]>0\}\) equals the multiset of
values on the path \(v\to\) root; enter/leave maintain this by induction on DFS events.
\textbf{(Histogram correctness.)} \(\mathit{bucket}[f]=|\{x:\mathit{cnt}[x]=f\}|\) is
updated \(-1\) on the old count and \(+1\) on the new whenever a single value moves, so it
always reflects the current path's frequency-of-frequency, and the Fenwick stores it, giving
\(S(f)=\#\{x:\mathit{cnt}[x]\le f\}\). \textbf{(Selection.)} The required ordered list is the
values with \(\mathit{cnt}\ge l\), ascending by count: group \(l\) first, then \(l{+}1\), etc.,
each group contributing \(\mathit{bucket}\) members in arbitrary (tie) order. The \(k\)-th
element lies in the smallest group \(f^\ast\ge l\) with \(\sum_{g=l}^{f^\ast}\mathit{bucket}[g]\ge k\),
i.e.\ \(S(f^\ast)\ge S(l-1)+k\) and \(S(f^\ast{-}1)<S(l-1)+k\) -- the Fenwick lower-bound of
\(\mathit{base}+k\). Existence fails exactly when \(S(n)-\mathit{base}<k\), giving \(-1\);
otherwise any member of group \(f^\ast\) (top of \(\mathit{valAt}[f^\ast]\), nonempty because
\(\mathit{bucket}[f^\ast]\ge1\)) is a valid answer. \textbf{(Unwind.)} Each node's enter is
mirrored by its leave, so after \(\mathit{dfs}(1)\) all counts, buckets and stacks are empty:
no inter-test contamination. Hence every printed answer matches the specification. \(\Box\)
\end{proofbox}

\solstep{Implementation}
\begin{lstlisting}[caption={CF 1585E Summit: a DFS from the root uses the recursion stack as the Mo window (enter bumps a value's frequency, leave undoes it); a frequency-of-frequency Fenwick selects the k-th distinct value by increasing count, and per-frequency value stacks hand back an answer value. O((n+q) log n). ASCII only, identifiers without underscores.}, label={lst:c21-summit}]
#include <bits/stdc++.h>
using namespace std;
const int MAXN=1000005;
int T,n,q,highest,a[MAXN];
vector<int> kid[MAXN];
struct QQ{ int l,k,id; };
vector<QQ> atNode[MAXN];
int ansArr[MAXN];
int cntVal[MAXN];            // occurrences of value on current root path
int bitFen[MAXN];            // Fenwick over bucket[f] = #values with freq f
vector<int> valAt[MAXN];     // the values currently having frequency f
int posIn[MAXN];             // index of a value inside valAt[its frequency]
void bitAdd(int i,int d){ for(;i<=n;i+=i&-i) bitFen[i]+=d; }
int bitSum(int i){ int s=0; for(;i>0;i-=i&-i) s+=bitFen[i]; return s; }
int bitLower(int target){ int pos=0;
  for(int step=highest; step; step>>=1){ int np=pos+step;
    if(np<=n && bitFen[np]<target){ pos=np; target-=bitFen[np]; } }
  return pos+1; }
void attachVal(int v,int f){ valAt[f].push_back(v); posIn[v]=valAt[f].size()-1; }
void detachVal(int v,int f){ int p=posIn[v]; int last=valAt[f].back();
  valAt[f][p]=last; posIn[last]=p; valAt[f].pop_back(); }
void moveVal(int v,int from,int to){         // value v changes count from->to
  detachVal(v,from); bitAdd(from,-1);
  attachVal(v,to);   bitAdd(to,1); }
void enter(int u){ int v=a[u], c=cntVal[v];
  if(c>0) moveVal(v,c,c+1); else { attachVal(v,1); bitAdd(1,1); }
  cntVal[v]=c+1; }
void leave(int u){ int v=a[u], c=cntVal[v];
  if(c>1) moveVal(v,c,c-1); else { detachVal(v,1); bitAdd(1,-1); }
  cntVal[v]=c-1; }
void dfs(int u){
  enter(u);
  for(const QQ &Q : atNode[u]){
    int base=bitSum(Q.l-1), total=bitSum(n);
    if(total-base < Q.k) ansArr[Q.id]=-1;
    else { int f=bitLower(base+Q.k); ansArr[Q.id]=valAt[f].back(); }
  }
  for(int w : kid[u]) dfs(w);
  leave(u);
}
int main(){
  scanf("%d",&T);
  while(T--){
    scanf("%d%d",&n,&q);
    for(int i=1;i<=n;i++) scanf("%d",&a[i]);
    for(int i=2;i<=n;i++){ int p; scanf("%d",&p); kid[p].push_back(i); }
    for(int i=1;i<=q;i++){ int v,l,k; scanf("%d%d%d",&v,&l,&k); atNode[v].push_back({l,k,i}); }
    highest=1; while((highest<<1)<=n) highest<<=1;
    dfs(1);
    for(int i=1;i<=q;i++) printf("%d ",ansArr[i]);
    printf("\n");
    for(int i=1;i<=n;i++){ kid[i].clear(); atNode[i].clear(); } // window auto-unwinds
  }
}
\end{lstlisting}

\solstep{Trace on the Official Samples}
\begin{tracebox}[Sample 1: chain \(1\!-\!2\!-\!3\), all values \(1\) -- the histogram is one spike]
At node~3 the root path is \(1,2,3\), values \([1,1,1]\), so \(\mathit{cnt}[1]=3\) and
\(\mathit{bucket}[3]=1\) (S(1)=S(2)=0, S(3)=1). \((3,1,1)\): base=S(0)=0, total=S(3)=1,
\(1{-}0\ge1\), target=1, bitLower(1)=3, valAt[3]=\{1\} \(\Rightarrow\) \(\boxed{1}\).
\((3,1,2)\): total\(-\)base\(=1<2\) \(\Rightarrow\) \(\boxed{-1}\). \((3,2,1)\): base=S(1)=0,
total=1, target=1 \(\Rightarrow f^\ast=3\) \(\Rightarrow\) \(\boxed{1}\). Output
\texttt{1 -1 1} -- match. \(\Box\)
\end{tracebox}

\begin{tracebox}[Sample 2: the trap that ties decide, and the LCA-instant histogram]
Tree \(1\!\to\!\{2,3\}\), \(2\!\to\!\{4,5\}\); \(a=[1,2,1,1,2]\).
\((3,1,1)\) at node 3, path values \([1,1]\): \(\mathit{cnt}[1]=2\), bucket\([2]=1\),
S(0)=0, total=S(5)=1, target=1, \(f^\ast=2\), val=1 \(\Rightarrow\boxed{1}\).
\((2,1,2)\) at node 2, path \([1,2]\): \(\mathit{cnt}[1]=\mathit{cnt}[2]=1\),
bucket\([1]=2\), base=S(0)=0, total=2, target=2, \(f^\ast=1\) (S(1)=2\(\ge\)2),
valAt[1]=\(\{1,2\}\) \(\Rightarrow\) \(\boxed{1\text{ (or }2\text{)}}\), a legal tie pick.
\((4,1,1)\) at node 4, path \([1,2,1]\): \(\mathit{cnt}[1]=2,\mathit{cnt}[2]=1\),
bucket\([1]=1,\)bucket\([2]=1\), target=1 \(\Rightarrow f^\ast=1\), val=2 \(\Rightarrow\boxed{2}\).
\((4,2,1)\): base=S(1)=1, total=2, target=2, \(f^\ast=2\), val=1 \(\Rightarrow\boxed{1}\).
\((4,2,2)\): total\(-\)base\(=2-1=1<2\) \(\Rightarrow\boxed{-1}\). Output
\texttt{1 1 2 1 -1} (the tie on query~2 printed as 1) -- match. \(\Box\)
\end{tracebox}

\solstep{Complexity}
\begin{complexitybox}
Each node is entered and left once; every move does two Fenwick updates and one
\(\mathit{valAt}\) relocate, \(O(\log n)\) amortised. Each query is answered with two
Fenwick sums and one lower-bound, \(O(\log n)\). Total \(\bigO{(n+q)\log n}\) per test,
\(\bigO{(\sum n+\sum q)\log n}\) over all cases -- inside the 4\,s limit at
\(\sum n,\sum q\le10^6\). \textbf{Space} \(\bigO{n+q}\): the child/query lists, the
Fenwick, \(\mathit{cnt}\), and the per-frequency \(\mathit{valAt}\) stacks whose total
size never exceeds the number of distinct values on the current path (\(\le n\)). No node
is ever re-walked and no histogram rebuilt -- the DFS-stack window and the
frequency-of-frequency Fenwick are the entire machinery.
\end{complexitybox}

\begin{hardeningbox}
\begin{itemize}[nosep,leftmargin=1.4em]
  \item \textbf{Do not run Mo here.} The answer is not a single running scalar; it needs
        the whole histogram re-sorted. Path Mo's window (\(\secref{c21-path}\)) gives the
        right \emph{set} but the wrong \emph{read}, and re-scanning frequencies per query
        is \(\tle\) (dead end 1 and 2).
  \item \textbf{The lower-bound lands in \(f^\ast\ge l\) for free.} Because
        \(\mathit{target}=\mathit{base}+k > S(l-1)\), the smallest \(f\) with
        \(S(f)\ge\mathit{target}\) is automatically \(\ge l\) -- you never special-case
        the ``drop counts below \(l\)'' step; it is absorbed by \(\mathit{base}=S(l-1)\).
  \item \textbf{The window auto-unwinds -- rely on it.} Matched enter/leave return every
        \(\mathit{cnt}\), bucket and \(\mathit{valAt}\) to empty at the end of the DFS, so
        per test you clear only adjacency and query lists; do \emph{not} memset the whole
        \(\bigO{n}\) arrays (the sum of \(n\) is \(10^6\), but a full-array clear per tiny
        test across \(t\le10^6\) cases would re-introduce an \(O(t\cdot n)\) blow-up).
  \item \textbf{Recursion depth at \(n=10^6\).} A chain overflows the call stack; convert
        \(\mathit{dfs}\) to an explicit stack (push enter, then children, then leave) --
        the same caution as Tarjan (\(\secref{c21-tarjan}\)).
\end{itemize}
\end{hardeningbox}

\begin{recallbox}[Chapter 27 in ten seconds]
\emph{When a huge offline batch of subtree/path queries has a statistic with no inverse,
block-sort the ranges and sweep one window with an \(O(1)\) add/remove --
\(\bigO{(n+q)\sqrt n}\) (\(\thmref{c21-complexity}\), \(\coderef{c21-mo}\)); a subtree is
the preorder range \([\tin(v),\tout(v)]\) (\(\thmref{c21-subtree-range}\)) and a path is a
doubled-tour toggle window split at the \(\lca\) plus the read-time LCA
(\(\thmref{c21-pathsplit}\), \(\coderef{c21-euler-toggle}\)); CF~617E is the array engine on
prefix xors (\(\probref{c21-xor}\)); CF~463E carries the root-path window on the DFS
stack with a per-prime push/undo table (\(\probref{c21-caisa}\)), while the
parity-normalised invertible case (Ch.~18's CF~383C, \(\chref{c14}\)) skips Mo
entirely; Tarjan's one DFS \(+\) DSU answers \(q\) LCAs in
\(\bigO{(n+q)\alpha}\) (\(\thmref{c21-tarjan}\)); and the four tests decide Mo versus
Fenwick / HLD / DSU-on-tree / virtual tree / centroid (\(\defref{c21-signature}\),
\(\defref{c21-pick}\)). The Summit CF~1585E (\(\probref{c21-summit}\)) turns the DFS stack
itself into the window and a frequency-of-frequency Fenwick into the selector.}
\selftest{
  \item State the four Mo tests and why CF~463E needs none of them. \hfill\lv{4}
  \item Give the doubled-tour window for path \((u,v)\) when \(u\) is not the ancestor, and the extra step. \hfill\lv{5}
  \item Why does Tarjan beat lifting only for \emph{offline} batches, and at what cost? \hfill\lv{4}
  \item Which structure turns ``\(k\)-th value by increasing frequency'' into one \(O(\log n)\) query? \hfill\lv{6}
}
\end{recallbox}

\chaptersection{Tool Kit Inventory}

The machine assembled through Chapter~27. Earlier chapters' rows are carried forward;
this chapter's are appended. Pointers only.

\begin{center}
\small
\renewcommand{\arraystretch}{1.25}
\setlength{\tabcolsep}{3pt}
\begin{tabular}{@{}p{0.24\linewidth} l c p{0.22\linewidth} >{\raggedright\arraybackslash\scriptsize}p{0.17\linewidth}@{}}
\toprule
\textbf{Tool} & \textbf{Since} & \textbf{Lv} & \textbf{Trigger phrase} & \textbf{Where} \\
\midrule
\multicolumn{5}{l}{\itshape Carried forward} \\
LCA + \(\dist\), binary lifting & 16 & \lv{4} & ``the split point of a path'' & \chref{c11}, \thmref{c11-distf} \\
Euler \(\tin/\tout\), flatten & 5 & \lv{3} & ``a subtree is one array range'' & \chref{c24} \\
Fenwick / difference array & 5/10 & \lv{3} & ``invertible range add, point query'' & \chref{c24}, \chref{c10} \\
DSU (find/union, path compress) & 17 & \lv{3} & ``merge finished regions'' & \chref{c28} \\
DSU on tree / small-to-large & 17 & \lv{4} & ``subtree statistic for (almost) every node'' & \chref{c28} \\
\midrule
\multicolumn{5}{l}{\itshape Chapter 27 -- new this chapter} \\
Mo block order + two pointers & 27.1 & \lv{4} & ``offline ranges, non-invertible \(O(1)\) update'' & \defref{c21-mo}, \coderef{c21-mo} \\
Prefix-xor pair counting & 27.1 & \lv{4} & ``count subarrays with \(\mathit{xor}=k\)'' & \probref{c21-xor}, \coderef{c21-xor} \\
Subtree flatten \(+\) Mo & 27.2 & \lv{4} & ``aggregate over subtree\((v)\), offline'' & \thmref{c21-subtree-range}, \coderef{c21-flatten} \\
Root-path stack + last-seen table & 27.2 & \lv{4} & ``deepest ancestor sharing a key'' & \probref{c21-caisa}, \coderef{c21-caisa} \\
Doubled-tour toggle + LCA split & 27.3 & \lv{5} & ``statistic on path \(u\)-\(v\), offline'' & \thmref{c21-pathsplit}, \coderef{c21-euler-toggle} \\
Tarjan offline LCA & 27.4 & \lv{4} & ``answer \(q\) LCAs, no lifting table'' & \thmref{c21-tarjan}, \coderef{c21-tarjan} \\
Recognition gate / boundary pick & 27.4/27.5 & \lv{5} & ``is this even a Mo job?'' & \defref{c21-signature}, \defref{c21-pick} \\
Mo kit (decision skeleton) & 27 & \lv{5} & ``the seven questions in one place'' & \coderef{c21-kit} \\
Summit: DFS-stack window + freq-of-freq BIT & Summit & \lv{6} & ``root path, order values by count'' & \coderef{c21-summit} \\
\bottomrule
\end{tabular}
\end{center}

\chaptersection{Related Topics}

\begin{itemize}[nosep, leftmargin=1.4em]
  \item \chref{c30} (Ch.~26, Virtual Trees) --- backward, \textbf{the hand-off this chapter
        answers}: Ch.~26's boundary (\(\secref{c30-boundary}\)) forwarded ``many overlapping
        offline node-set queries with a non-invertible statistic'' here; \(\secref{c21-moengine}\)
        opens by discharging exactly that debt. The virtual tree re-folds \emph{per marked
        set}; Mo \emph{shares work across} the batch.
  \item \chref{c11} (Ch.~16, LCA) --- tool: \(\lca\) splits the path window (\(\secref{c21-path}\))
        and is the quantity Tarjan (\(\secref{c21-tarjan}\)) computes offline; \(\thmref{c11-distf}\)
        is used, never re-derived.
  \item \chref{c24} (Ch.~5, Euler Tour / \(\tin/\tout\)) --- tool: the subtree flatten
        (\(\defref{c21-flatten}\)) is Ch.~5's array; the doubled tour of \(\secref{c21-path}\)
        extends it, and Ch.~18's difference-Fenwick solve of CF~383C
        (\(\chref{c14}\)) is Ch.~5's invertible primitive.
  \item \chref{c28} (Ch.~17, DSU on tree) --- contrast and boundary: DSU-on-tree and Mo both
        answer subtree distinct-colour folds (CF~375D, whose URL lives there); \(\secref{c21-boundary}\)
        prices \(\bigO{n\log n}\) DSU against \(\bigO{(n+q)\sqrt n}\) Mo and picks per shape.
  \item \chref{c16} (Ch.~22, HLD) / \chref{c23} (Ch.~25, Link-Cut) --- boundary: every
        \emph{online} path or update-interleaved query leaves Mo (which is offline-only)
        for these (\(\defref{c21-signature}\)).
  \item \chref{c18} (Ch.~24, Centroid Decomposition) --- boundary: whole-tree pair/distance
        distributions with no query list go to centroid, not a window sweep
        (\(\secref{c21-boundary}\)); the Summit's frequency-of-frequency also echoes centroid's
        ``fold over a global distribution''.
  \item \chref{c22} (Ch.~28, Tree Isomorphism and Hashing) --- \textbf{next chapter}: where Mo
        compares \emph{counts}, hashing compares \emph{structures}; the note in
        \(\secref{c21-boundary}\) hands the ``which subtree equals which'' question forward.
\end{itemize}

\begin{lastword}[Last Word --- Chapter 27]
\emph{``When a statistic has no inverse, do not try to be clever about prefixes -- be
clever about \textbf{order}. Block-sort the offline queries so consecutive windows differ
by one node, keep an \(O(1)\) add/remove, and a \(\bigO{nq}\) nightmare becomes
\(\bigO{(n+q)\sqrt n}\); the tree enters through two doors, the preorder flatten (a subtree
is one range) and the doubled Euler tour (a path is a toggle window split at the
\(\lca\)). And when the batch only wants \(\lca\), Tarjan's DFS \(+\) DSU answers it in
near-linear time with no lifting table at all.''}

\medskip
Chapter~26 handed this chapter a problem it could not compress, and Chapter~27 collected
it. Section 27.1 built the engine and proved \(\bigO{(n+q)\sqrt n}\)
(\(\thmref{c21-complexity}\)), cashing it in on CF~617E's prefix-xor pairs
(\(\probref{c21-xor}\)); 27.2 bolted on the tree via the subtree flatten
(\(\thmref{c21-subtree-range}\)), handed the scary depth-parity update to Ch.~18's
Fenwick as the invertible case it is (\(\chref{c14}\)), and opened a third door with
CF~463E -- the recursion stack carrying a per-prime push/undo table through
fifty-one rebuilds (\(\probref{c21-caisa}\)); 27.3 invented the \term{toggle} so a path -- which no
flatten can make contiguous -- becomes a window, and disarmed the LCA-parity trap by hand
on three walks of \(R\) (\(\thmref{c21-pathsplit}\)); 27.4 gave the offline \(\lca\) with
Tarjan's WHITE/GREY/BLACK DSU run traced to \(1,2,2,3,1\) (\(\thmref{c21-tarjan}\)); and
27.5 drew the boundary so you reach for the sweep only after Fenwick, HLD, DSU-on-tree,
virtual tree and centroid have each been offered the job and declined it
(\(\defref{c21-pick}\)). The Summit CF~1585E (\(\probref{c21-summit}\)) fused it all: the
DFS stack \emph{is} the Mo window, and a frequency-of-frequency Fenwick reads the answer.

\medskip
Carry three reflexes forward:
\begin{enumerate}[nosep, leftmargin=1.6em]
  \item \emph{Look for the inverse first.} If \(\mathrm{pref}(r)-\mathrm{pref}(l-1)\)
        exists -- even after a \((-1)^{\depth}\) normalisation -- Mo is a \(\sqrt n\) tax
        you chose to pay (\(\chref{c14}\)).
  \item \emph{Choose the array before the algorithm.} Subtree \(\to\) preorder flatten;
        path \(\to\) doubled tour with a toggle; neither is ``the tree'' -- the sweep runs
        on a line (\(\thmref{c21-subtree-range}\), \(\thmref{c21-pathsplit}\)).
  \item \emph{Respect modality.} Mo is offline and update-free; one interleaved update and
        the block order is illegal -- that is HLD's or the LCT's contract, not yours
        (\(\defref{c21-signature}\)).
\end{enumerate}

\medskip
And three habits: get \(B\approx n/\sqrt q\) right and verify the \texttt{add}/\texttt{rm}
are true inverses before trusting the sweep; add the path-Mo LCA at \emph{read} time only
in case (ii); and when a whole batch is LCA, reach for Tarjan (\(\bigO{(n+q)\alpha}\), no
table) before building an \(O(n\log n)\) lifting structure. You now own the one engine that
answers tree queries by \emph{reordering the questions} rather than by folding the tree --
the offline, amortised companion to every prefix, path, and per-set structure before it.

\medskip
\textbf{Next:} \chref{c22} (Tree Isomorphism and Hashing) changes the question entirely.
Mo and Tarjan compare \emph{counts} on sets of nodes; the next chapter asks when two trees
are \emph{the same tree} -- canonical forms (AHU), subtree hashing for structural equality,
centre-based unrooted comparison, and counting distinct subtrees -- the equality test this
chapter never needed, and the tool behind cutting-point and symmetry problems like CF~1800G
and CF~1252F.
\end{lastword}


\closechapter
